% lain-spec.tex — The Lain Programming Language: Language Specification
%
% Compile with:
%   latexmk -pdf lain-spec.tex
% or:
%   pdflatex lain-spec.tex  (run twice for cross-references)
%
% Structure:
%   Frontmatter  — Foreword, Introduction, TOC
%   Clauses 1–4  — Administrative (Scope, References, Terms, Conformance)
%   Clauses 5–20 — Normative language specification
%   Annex A–B    — Normative (Grammar, Diagnostics)
%   Annex C–D    — Informative (Rationale, Comparison)

\documentclass[11pt,a4paper,twoside,openright]{book}
\usepackage{lain-spec}

% =========================================================================
% Document metadata
% =========================================================================
\title{%
  \vspace{-2cm}
  {\LARGE\bfseries\color{lain-blue} The Lain Programming Language}\\[0.8em]
  {\large Language Specification}\\[0.4em]
  {\normalsize Version 0.1 --- Draft}
}
\author{}
\date{\today}

% =========================================================================
% Begin document
% =========================================================================
\begin{document}

% -------------------------------------------------------------------------
% Front matter
% -------------------------------------------------------------------------
\frontmatter

\maketitle
\thispagestyle{empty}

\vfill
\noindent
\begin{tabular}{ll}
  \textbf{Document status:} & Draft --- not for normative use \\
  \textbf{Language version:} & Lain 1.0 \\
  \textbf{Date:}             & \today \\
\end{tabular}
\vfill
\noindent\textit{%
  This document is the authoritative specification of the Lain programming
  language.  It defines the syntax, semantics, and constraints that every
  conforming implementation shall enforce.
}
\clearpage

\tableofcontents
\clearpage

% 00-frontmatter.tex — Foreword and Introduction (non-normative)

% =========================================================================
% FOREWORD (non-normative)
% =========================================================================
\chapter*{Foreword}
\addcontentsline{toc}{chapter}{Foreword}

This document specifies the Lain programming language, Version~1.0.

Lain is a statically-typed, compiled programming language for systems
programming.  It guarantees memory safety, type safety, and resource safety
entirely at compile time with zero runtime overhead.

This specification is structured analogously to ISO/IEC~9899 (the C
standard).  Clauses~1 through~4 are administrative; Clauses~5 through~20
are normative.  Annexes~A and~B are normative; Annexes~C and~D are
informative.

Text marked \textbf{Constraint} describes a condition whose violation
requires an implementation to issue a diagnostic message.  Text marked
\idbehav{} denotes behavior that is valid but whose exact result is
determined by the implementation and shall be documented by it.  Text
marked \undefbehav{} denotes a situation for which this specification
imposes no requirements.

\begin{note}
  Undefined behavior in Lain can arise only inside \lain{unsafe} blocks or
  through direct C interoperability.  All code in the safe fragment of Lain
  has defined behavior.
\end{note}

% =========================================================================
% INTRODUCTION (non-normative)
% =========================================================================
\chapter*{Introduction}
\addcontentsline{toc}{chapter}{Introduction}

\section*{Design philosophy}

Lain is built on five foundational pillars:

\begin{description}[leftmargin=2em]
  \item[P1 --- Assembly-Speed Performance]
    Every language feature compiles to C code that is equivalent in
    performance to hand-written C.  There are no hidden allocations, no
    garbage collector, and no reference-counting overhead.

  \item[P2 --- Zero-Cost Memory Safety]
    Memory safety is enforced entirely at compile time through a linear type
    system, a borrow checker, and value range analysis.  A conforming
    implementation that rejects all constraint violations guarantees, for
    safe programs, the absence of use-after-free, double-free, null
    dereference, and buffer overflow at runtime.

  \item[P3 --- Static Verification]
    The verification properties required by this specification are decidable
    and polynomial in program size.  No SMT solver or unbounded fixpoint
    iteration is required.

  \item[P4 --- Determinism]
    Pure functions (\lain{func}) are referentially transparent and
    guaranteed to terminate.  The same inputs always produce the same
    outputs.

  \item[P5 --- Syntactic Simplicity]
    Lain has no implicit copies, no hidden conversions (beyond the defined
    integer-widening rules), and no magic methods.  Ownership is always
    explicit.
\end{description}

\section*{Relationship to C}

Lain compiles to C99.  The C code generated by a conforming Lain
implementation is the \emph{translation output}; this specification does not
prescribe the form of that output beyond the observable behavior of the
resulting program.

\section*{How to read this specification}

Each normative clause is structured as follows: a prose description of the
syntax (with a reference to the formal grammar in Annex~A), followed by the
semantics, followed by \textbf{Constraint} blocks listing every condition
that requires a diagnostic message.  Non-normative \textbf{Note} and
\textbf{Example} environments appear throughout to aid understanding.


% -------------------------------------------------------------------------
% Main matter — Normative clauses
% -------------------------------------------------------------------------
\mainmatter

% --- Administrative clauses ---
% 01-scope.tex — Clause 1: Scope
\chapter{Scope}
\label{sec:scope}
\index{scope (of document)}

\pnum
This document specifies the Lain programming language.  It defines the
syntax, the semantic rules, and the constraints that every conforming Lain
implementation shall enforce.  It serves as the single authoritative
reference for language users and implementors alike.

\pnum
Lain is a statically-typed, compiled programming language designed for
systems programming, embedded systems, and safety-critical software.  It
achieves memory safety, type safety, and resource safety entirely at compile
time with zero runtime overhead.

\pnum
This specification covers:

\begin{itemize}
  \item lexical structure and grammar (\S\,\ref{sec:05-lexical});
  \item type system and type inference (\S\,\ref{sec:07-types});
  \item variable declarations, scoping, and initialization
        (\S\,\ref{sec:06-basics}, \S\,\ref{sec:10-declarations});
  \item expressions, operators, and evaluation order
        (\S\,\ref{sec:08-expressions});
  \item statements and control flow (\S\,\ref{sec:09-statements});
  \item functions, procedures, and purity (\S\,\ref{sec:12-functions});
  \item ownership, borrowing, and the linear type system
        (\S\,\ref{sec:11-ownership});
  \item value range constraints and static verification
        (\S\,\ref{sec:13-vra});
  \item compile-time evaluation and generics (\S\,\ref{sec:14-generics});
  \item pattern matching and exhaustiveness checking
        (\S\,\ref{sec:15-match});
  \item module system (\S\,\ref{sec:16-modules});
  \item C interoperability (\S\,\ref{sec:17-interop});
  \item unsafe code boundaries (\S\,\ref{sec:18-unsafe});
  \item memory model (\S\,\ref{sec:19-memory});
  \item diagnostic messages (Annex~B);
  \item standard library (\S\,\ref{sec:20-stdlib}).
\end{itemize}

\pnum
This specification does \emph{not} cover:

\begin{itemize}
  \item the internal architecture of any particular compiler implementation;
  \item build system integration or tooling;
  \item IDE or editor support;
  \item operating system interfaces beyond what the standard library provides.
\end{itemize}

\section{Compilation model}
\label{sec:scope.compilation}
\index{compilation model}

\pnum
Lain uses a source-to-C transpilation model.  A Lain source file (extension
\texttt{.ln}) is processed by a conforming implementation to produce C99
source code, which is subsequently compiled by a C99 compiler to produce an
executable.  The observable behavior of the program is defined by this
specification; the form of the intermediate C99 code is
\idbehav{the generated C99 code need not be human-readable}.

\pnum
A conforming implementation shall perform the following translation phases
in the following order:

\begin{enumerate}
  \item \textbf{Lexical analysis} --- source text is converted to a token
        stream as described in \S\,\ref{sec:05-lexical}.
  \item \textbf{Parsing} --- the token stream is reduced to an abstract
        syntax tree according to the grammar in Annex~A.
  \item \textbf{Name resolution} --- identifiers are resolved to
        declarations as described in \S\,\ref{sec:06-basics}.
  \item \textbf{Type checking} --- types are inferred and verified as
        described in \S\,\ref{sec:07-types} through
        \S\,\ref{sec:10-declarations}.
  \item \textbf{Exhaustiveness checking} --- pattern matching completeness
        is verified as described in \S\,\ref{sec:15-match}.
  \item \textbf{Use analysis (NLL)} --- last-use points are computed for
        all variables.
  \item \textbf{Linearity and borrow checking} --- ownership and borrowing
        rules are verified as described in \S\,\ref{sec:11-ownership}.
  \item \textbf{Value range analysis} --- bounds safety and numeric
        constraints are verified as described in \S\,\ref{sec:13-vra}.
  \item \textbf{Code emission} --- well-formed programs are translated to
        C99.
\end{enumerate}

\pnum
A program that completes all phases without a diagnostic message is
\textbf{well-formed} and its translation output is defined.

\begin{note}
  Phases~3--8 collectively form \emph{semantic analysis}.  An
  implementation may interleave these phases as long as the observable
  behavior is equivalent to executing them in the order listed.
\end{note}

% 02-references.tex — Clause 2: Normative references
\chapter{Normative references}
\label{sec:references}
\index{normative references}

\pnum
The following documents are referred to in the text in such a way that some
or all of their content constitutes requirements of this document.  For
dated references, only the edition cited applies.

\begin{description}[leftmargin=2em, style=nextline]
  \item[ISO/IEC 9899:2018 (C17)]
    \textit{Information technology --- Programming languages --- C}.\\
    International Organization for Standardization, Geneva, 2018.\\
    Referenced for: C99/C17 calling conventions (\S\,\ref{sec:17-interop}),
    standard integer types and their representations
    (\S\,\ref{sec:07-types}), and undefined behavior categories
    (\S\,\ref{sec:18-unsafe}).

  \item[IEC 60559:2020 (IEEE 754-2019)]
    \textit{Floating-point arithmetic}.\\
    International Electrotechnical Commission, Geneva, 2020.\\
    Referenced for: semantics of \lain{f32} and \lain{f64} arithmetic
    (\S\,\ref{sec:07-types}).
\end{description}

\pnum
The following documents are informative references.  They are not required
for conformance but provide useful background.

\begin{description}[leftmargin=2em, style=nextline]
  \item[ISO/IEC 14882:2020 (C++20)]
    \textit{Programming languages --- C++}.\\
    Referenced for: vocabulary and structure conventions used in this
    document.

  \item[The Rust Reference]
    \textit{The Rust Reference}, \url{https://doc.rust-lang.org/reference/}.\\
    Referenced for: linear type and ownership terminology.
\end{description}

% 03-terms.tex — Clause 3: Terms, definitions and symbols
\chapter{Terms, definitions, and symbols}
\label{sec:terms}
\index{terms and definitions}

\pnum
For the purposes of this document, the following terms and definitions
apply.  Terms defined here are used with these precise meanings throughout
all normative clauses.  Where a term is used in its ordinary English sense
rather than as a defined term, it appears without emphasis.

\begin{note}
  Cross-references of the form \textit{(see \S\,N)} point to the clause
  where the concept is treated normatively.
\end{note}

% -------------------------------------------------------------------------
\section{General language terms}
\label{sec:terms.general}
% -------------------------------------------------------------------------

\subsection*{abstract syntax tree (AST)}
The in-memory tree representation of a Lain program produced by the parser
and consumed by semantic analysis.

\subsection*{behavior}
External appearance or action of a program as observed by a user or
another program.  See also: \textit{defined behavior},
\textit{implementation-defined behavior}, \textit{unspecified behavior},
\textit{undefined behavior}.

\subsection*{conforming implementation}
An implementation that correctly processes all programs that satisfy the
rules of this specification and issues a diagnostic message for every
program that violates a \textbf{Constraint} (see \S\,\ref{sec:conformance}).

\subsection*{conforming program}
A program that uses only the syntax and semantics defined in this
specification, does not rely on undefined behavior, and does not depend on
implementation-defined extensions.

\subsection*{constraint}
A requirement imposed on the source text.  If a program violates a
constraint, the implementation \textit{shall} emit a diagnostic message.
Constraints appear in \textbf{Constraint} environments throughout this
specification.

\subsection*{defined behavior}
Behavior that is fully and precisely specified by this document for every
conforming program.  Contrast with \textit{undefined behavior} and
\textit{implementation-defined behavior}.

\subsection*{diagnostic message}
\index{diagnostic message}
A human-readable message issued by the implementation to report a
constraint violation.  Every diagnostic shall include a source location
(file name, line number) and an error code of the form \texttt{[EXXX]} or
\texttt{[VRA]} (see Annex~B).

\subsection*{ill-formed program}
A program that violates one or more of the syntactic or semantic rules of
this specification.  An ill-formed program is not a conforming program and
a conforming implementation shall reject it with a diagnostic.

\subsection*{implementation}
The system (compiler, linker, runtime) that translates a Lain source program
into an executable.  The reference implementation is the \textit{lain}
compiler.

\subsection*{implementation-defined behavior}
\index{implementation-defined behavior}
Behavior that is valid but not fully specified by this document; the
implementation shall determine the exact behavior and document it.

\subsection*{program}
A collection of one or more Lain source files (\texttt{.ln}) that together
constitute a translation unit or a set of modules.

\subsection*{source file}
\index{source file}
A UTF-8 encoded text file with the extension \texttt{.ln} containing Lain
source text.

\subsection*{translation output}
The C99 source code emitted by a conforming implementation for a
well-formed Lain program.

\subsection*{undefined behavior}
\index{undefined behavior}
Behavior for which this specification imposes no requirements.  A program
that exhibits undefined behavior is erroneous.  In the safe fragment of
Lain, no operation produces undefined behavior.  Undefined behavior may
arise only inside \lain{unsafe} blocks or through direct C
interoperability.

\subsection*{unspecified behavior}
\index{unspecified behavior}
Behavior for which this specification provides two or more possibilities
without prescribing which shall occur.  A conforming program shall not
depend on unspecified behavior.

\subsection*{well-formed program}
A program that satisfies all syntactic and semantic rules of this
specification and is therefore accepted by a conforming implementation
without diagnostic.

% -------------------------------------------------------------------------
\section{Type system terms}
\label{sec:terms.types}
% -------------------------------------------------------------------------

\subsection*{algebraic data type (ADT)}
\index{algebraic data type}
A composite type that is either a product type (\textit{struct}) or a sum
type (\textit{enum}).

\subsection*{constrained type}
\index{constrained type}
A type annotated with a value range constraint of the form
\textit{T op bound}, e.g.\ \lain{int >= 0}
(see \S\,\ref{sec:07-types}).

\subsection*{integer rank}
\index{integer rank}
A total ordering on integer types that determines the result type of
mixed-type arithmetic.  The rank system is defined in
\S\,\ref{sec:07-types}.

\subsection*{integer widening}
\index{integer widening}
An implicit conversion from an integer type of lower rank to one of
higher rank, permitted in arithmetic and assignment contexts
(see \S\,\ref{sec:07-types}).

\subsection*{linear type}
\index{linear type}
A type whose values must be used exactly once.  Struct types and resource
types are linear.  Primitive types (\lain{int}, \lain{bool}, etc.) are
unrestricted (see \S\,\ref{sec:11-ownership}).

\subsection*{nominal type}
A type identified by its declared name rather than its structure.

\subsection*{ownership mode}
\index{ownership mode}
An annotation on a variable or parameter that specifies the kind of
access the holder has: \lain{shared} (read-only), \lain{var}
(read-write), or owned (no annotation on a linear type, transferred with
\lain{mov}) (see \S\,\ref{sec:11-ownership}).

\subsection*{primitive type}
\index{primitive type}
One of the built-in types: \lain{int}, \lain{u8}, \lain{u16}, \lain{u32},
\lain{u64}, \lain{usize}, \lain{i8}, \lain{i16}, \lain{i32}, \lain{i64},
\lain{isize}, \lain{f32}, \lain{f64}, \lain{bool}.  Primitive types are
unrestricted (non-linear).

\subsection*{structural type equivalence}
Two types are structurally equivalent if they have the same structure.
Lain uses nominal equivalence for user-defined types; the name is part of
the type identity.

\subsection*{unrestricted type}
\index{unrestricted type}
A type whose values may be used any number of times without explicit
consumption.  All primitive types are unrestricted.  Contrast with
\textit{linear type}.

% -------------------------------------------------------------------------
\section{Ownership and memory terms}
\label{sec:terms.ownership}
% -------------------------------------------------------------------------

\subsection*{active borrow}
\index{borrow!active}
A borrow that is currently in the ACTIVE phase; its subject cannot be moved
or mutated while the borrow is active (see \S\,\ref{sec:11-ownership}).

\subsection*{arena}
\index{arena}
A memory allocation strategy in which all allocations are made from a
contiguous block; the entire block is freed at once.  Lain's compiler
allocates all AST nodes from arenas.

\subsection*{borrow}
\index{borrow}
A temporary access to a value without taking ownership.  A shared borrow
grants read-only access; a mutable borrow grants read-write access
(see \S\,\ref{sec:11-ownership}).

\subsection*{borrow phase}
\index{borrow!phase}
Either RESERVED (the borrow has been registered but is not yet active,
during argument evaluation) or ACTIVE (the borrow is in force)
(see \S\,\ref{sec:11-ownership}).

\subsection*{consumed variable}
\index{variable!consumed}
A linear variable whose ownership has been transferred (via \lain{mov}).
Reading or transferring a consumed variable is a constraint violation.

\subsection*{linear state}
\index{linear state}
The tracking state of a linear variable: one of \textit{Uninit},
\textit{Unconsumed}, or \textit{Consumed} (see \S\,\ref{sec:11-ownership}).

\subsection*{move}
\index{move}
The transfer of ownership of a value from one variable or storage location
to another, performed by the \lain{mov} operator.  After a move, the
source is in the Consumed state.

\subsection*{non-lexical lifetime (NLL)}
\index{non-lexical lifetime}
A lifetime that ends at the variable's \textit{last use} rather than at the
end of its lexical scope.  Lain uses NLL to minimize the duration of borrows
(see \S\,\ref{sec:11-ownership}).

\subsection*{read-write lock invariant}
\index{read-write lock invariant}
The aliasing invariant: at any program point, for any owned value, there
exist either \textit{N}~shared borrows (N~$\geq$~0) and no mutable borrow,
or exactly one mutable borrow and no shared borrows
(see \S\,\ref{sec:11-ownership}).

\subsection*{resource}
\index{resource}
A value that represents an external entity (file handle, socket, etc.)
whose lifetime must be explicitly managed.  Resources are modelled as
linear types.

\subsection*{two-phase borrow}
\index{two-phase borrow}
A borrow that begins in the RESERVED phase during argument evaluation and
is promoted to ACTIVE after the call expression is complete
(see \S\,\ref{sec:11-ownership}).

% -------------------------------------------------------------------------
\section{Control flow and purity terms}
\label{sec:terms.purity}
% -------------------------------------------------------------------------

\subsection*{decreasing measure}
\index{decreasing measure}
An expression in a \lain{while decreasing} statement that the implementation
shall verify is non-negative when the loop condition holds and strictly
decreases on every iteration (see \S\,\ref{sec:09-statements}).

\subsection*{\texttt{func}}
\index{func}
A pure, terminating function.  A \lain{func} shall not call a \lain{proc}
or contain an unbounded \lain{while} loop (see \S\,\ref{sec:12-functions}).

\subsection*{purity}
\index{purity}
The property of a \lain{func}: its result depends only on its inputs and it
has no observable side effects on any state outside its own stack frame.

\subsection*{\texttt{proc}}
\index{proc}
A procedure that may have side effects.  A \lain{proc} may call other
\lain{proc}s, perform I/O, and contain unbounded \lain{while} loops
(see \S\,\ref{sec:12-functions}).

\subsection*{termination}
\index{termination}
The property that a function always reaches a \lain{return} statement in
a finite number of steps.  All \lain{func}s shall be verified to terminate
by the implementation (see \S\,\ref{sec:12-functions}).

% -------------------------------------------------------------------------
\section{Verification terms}
\label{sec:terms.verification}
% -------------------------------------------------------------------------

\subsection*{bounds check}
\index{bounds check}
A runtime test that verifies an array index is within the valid range
\texttt{[0, N-1]}.  A conforming implementation shall eliminate bounds
checks when VRA can statically prove the index is in range.

\subsection*{comptime}
\index{comptime}
Compile-time evaluation.  A \lain{comptime} expression or function is
evaluated by the compiler at compile time rather than at runtime
(see \S\,\ref{sec:14-generics}).

\subsection*{exhaustiveness}
\index{exhaustiveness}
The property of a \lain{case} expression in which every possible value of
the scrutinee type is covered by at least one arm
(see \S\,\ref{sec:15-match}).

\subsection*{in guard}
\index{in guard}
A use of the \lain{in} operator as a loop or conditional guard.  An in
guard of the form \lain{idx in arr} proves to the VRA that
\texttt{idx} $\in$ \texttt{[0, arr.len-1]} in the guarded body
(see \S\,\ref{sec:08-expressions}, \S\,\ref{sec:13-vra}).

\subsection*{interval}
\index{interval}
A pair $[\ell, u]$ representing the set of integers $\{x \mid \ell \leq x
\leq u\}$.  VRA represents value ranges as intervals with $-\infty$ and
$+\infty$ sentinels.

\subsection*{monomorphization}
\index{monomorphization}
The process by which a generic function is instantiated for each concrete
set of type arguments, producing a separate, non-generic translation for
each instantiation (see \S\,\ref{sec:14-generics}).

\subsection*{return constraint}
\index{return constraint}
A value range constraint on the return value of a function, declared in the
function signature as \textit{T op bound}
(see \S\,\ref{sec:13-vra}).

\subsection*{value range analysis (VRA)}
\index{value range analysis (VRA)}
A static analysis that tracks integer value ranges through the program using
interval arithmetic, used to eliminate bounds checks and prove
constraint satisfaction (see \S\,\ref{sec:13-vra}).

% -------------------------------------------------------------------------
\section{Module and name resolution terms}
\label{sec:terms.modules}
% -------------------------------------------------------------------------

\subsection*{import}
\index{import}
A directive (\lain{import module.path}) that makes the names defined in
another module available in the current module
(see \S\,\ref{sec:16-modules}).

\subsection*{module}
\index{module}
The unit of compilation in Lain.  Each source file defines one module.

\subsection*{qualified name}
\index{qualified name}
A name of the form \textit{module.identifier} that refers to a declaration
in a specific module (see \S\,\ref{sec:06-basics}).

\subsection*{Universal Function Call Syntax (UFCS)}
\index{UFCS}
\index{Universal Function Call Syntax}
A syntactic sugar by which \lain{x.f(args)} is desugared to
\lain{f(x, args)} before name resolution.  UFCS allows methods to be
defined as free functions (see \S\,\ref{sec:12-functions}).

% 04-conformance.tex — Clause 4: Conformance
\chapter{Conformance}
\label{sec:conformance}
\index{conformance}

\section{Conforming implementation}
\label{sec:conformance.impl}
\index{conforming implementation}

\pnum
A \textbf{conforming implementation} shall:

\begin{enumerate}
  \item accept, without diagnostic, every \textit{conforming program};
  \item reject, with at least one diagnostic message, every program that
        violates any constraint stated in this specification;
  \item translate accepted programs such that their observable behavior
        conforms to the semantic rules stated in this specification;
  \item implement all features described as normative in Clauses~5
        through~20;
  \item document every behavior described in this specification as
        \idbehav{the exact value or behavior is implementation-defined}.
\end{enumerate}

\pnum
A conforming implementation may:

\begin{enumerate}
  \item issue additional diagnostic messages (warnings) beyond those
        required by this specification, provided they do not cause a
        conforming program to be rejected;
  \item accept programs through implementation-defined extensions, provided
        those extensions are documented and do not alter the semantics of
        any conforming program;
  \item omit runtime bounds checks when value range analysis
        (\S\,\ref{sec:13-vra}) proves the access is within bounds.
\end{enumerate}

\begin{note}
  The reference implementation of Lain (\textit{lain} compiler, version
  corresponding to this specification) is a conforming implementation.
  Where this specification marks a behavior as implementation-defined, the
  reference implementation's documented behavior takes precedence for
  programs compiled by that implementation.
\end{note}

\section{Conforming program}
\label{sec:conformance.program}
\index{conforming program}

\pnum
A \textbf{conforming program} is a program that:

\begin{enumerate}
  \item consists only of syntax and semantics defined in this
        specification;
  \item does not rely on undefined behavior;
  \item does not rely on unspecified behavior to produce a specific outcome;
  \item does not depend on any implementation-defined extension.
\end{enumerate}

\pnum
A \textbf{strictly conforming program} additionally:

\begin{enumerate}
  \item does not use any behavior documented as implementation-defined,
        even when that behavior is deterministic on a given implementation.
\end{enumerate}

\section{Diagnostic messages}
\label{sec:conformance.diagnostics}
\index{diagnostic message!requirements}

\pnum
A diagnostic message issued for a constraint violation shall include:

\begin{enumerate}
  \item the source file name;
  \item the line number of the offending construct;
  \item the diagnostic code in the form \texttt{[EXXX]} (see Annex~B);
  \item a human-readable description of the violation.
\end{enumerate}

\pnum
A conforming implementation may issue additional diagnostic information
(column number, source context, suggested fix) beyond the minimum required
by \S\,\ref{sec:conformance.diagnostics}.

\begin{note}
  The reference implementation additionally prints the source line and a
  caret (\texttt{\^{}}) pointing to the column of the offending token.
\end{note}

\section{Implementation limits}
\label{sec:conformance.limits}
\index{implementation limits}

\pnum
A conforming implementation shall support at least the following minimum
translation limits.  Programs that require more than these limits invoke
\idbehav{the behavior when a translation limit is exceeded}.

\begin{longtable}{@{}lll@{}}
\toprule
\textbf{Resource} & \textbf{Minimum} & \textbf{Clause} \\
\midrule
\endfirsthead
\multicolumn{3}{l}{\textit{(continued)}}\\
\toprule
\textbf{Resource} & \textbf{Minimum} & \textbf{Clause} \\
\midrule
\endhead
\bottomrule
\endlastfoot
Identifier length (bytes)                      & 255         & \S\,\ref{sec:05-lexical} \\
Block nesting depth                            & 64          & \S\,\ref{sec:09-statements} \\
Function parameters                            & 127         & \S\,\ref{sec:10-declarations} \\
Array elements                                 & $2^{31}-1$  & \S\,\ref{sec:07-types} \\
Module import depth                            & 32          & \S\,\ref{sec:16-modules} \\
\lain{case} arms per match expression          & 256         & \S\,\ref{sec:15-match} \\
Struct fields                                  & 256         & \S\,\ref{sec:07-types} \\
Enum variants                                  & 256         & \S\,\ref{sec:07-types} \\
Comptime evaluation depth                      & 256         & \S\,\ref{sec:14-generics} \\
Scope nesting depth                            & 64          & \S\,\ref{sec:06-basics} \\
\end{longtable}

\section{Safe and unsafe fragments}
\label{sec:conformance.safe-unsafe}
\index{safe fragment}
\index{unsafe fragment}

\pnum
The \textbf{safe fragment} of a Lain program consists of all code that does
not appear inside an \lain{unsafe} block and does not use raw pointer
operations.

\pnum
A conforming implementation shall guarantee, for any well-formed program
in the safe fragment, that execution produces no:

\begin{itemize}
  \item use of a variable after its ownership has been transferred (use-after-move);
  \item double transfer of ownership of the same value;
  \item access to an uninitialized variable;
  \item violation of the read-write lock invariant
        (\S\,\ref{sec:terms.ownership});
  \item buffer overrun that has been statically proven safe by value range
        analysis.
\end{itemize}

\pnum
The \textbf{unsafe fragment} comprises all code inside \lain{unsafe}
blocks.  Within the unsafe fragment, the programmer is responsible for
upholding the memory safety invariants listed in
\S\,\ref{sec:18-unsafe}; the implementation shall not be required to
verify them.

\begin{constraint}
  The constraints of the ownership and linearity system
  (\S\,\ref{sec:11-ownership}) are \emph{not} suspended inside
  \lain{unsafe} blocks.  Only raw pointer operations and the bounds-check
  requirement are relaxed.  An implementation shall continue to enforce
  \diagcode{E001}--\diagcode{E005} inside \lain{unsafe} blocks.
\end{constraint}


% --- Language specification ---
% 05-lexical.tex — Clause 5: Lexical conventions
\chapter{Lexical conventions}
\label{sec:05-lexical}
\index{lexical conventions}

\pnum
A Lain source file is translated in two conceptual steps: (1) the source
text is decomposed into tokens as described in this clause; (2) the token
stream is parsed into an abstract syntax tree according to the grammar.
Implementation may combine these steps.

\section{Source file encoding}
\label{sec:lex.encoding}
\index{source file!encoding}

\pnum
A Lain source file shall be a sequence of bytes encoded in UTF-8 (Unicode
Standard, Annex D).  The byte order mark (U+FEFF) shall not appear in a
source file.  If it does, behavior is \idbehav{whether a leading BOM is
silently skipped or causes a diagnostic}.

\pnum
Multi-byte UTF-8 sequences are permitted inside string literals and comments
but shall not appear in identifiers.  Identifiers shall consist solely of
ASCII characters (see \S\,\ref{sec:lex.identifiers}).

\pnum
Line endings are normalized: a carriage return (U+000D) followed
immediately by a line feed (U+000A) is treated as a single line feed.
A lone carriage return is treated as a line feed.

\section{Comments}
\label{sec:lex.comments}
\index{comment}

\pnum
A \textit{line comment} begins with \tok{//} and extends to the first
following line feed character.  The line feed is not part of the comment.

\pnum
A \textit{block comment} begins with \tok{/*} and ends with the nearest
following \tok{*/}.  Block comments may be nested: each opening \tok{/*}
inside a block comment opens a new nesting level requiring its own closing
\tok{*/}.

\pnum
Comments are stripped during lexical analysis and have no effect on the
semantics of the enclosing program.

\begin{constraint}
  An unterminated block comment (one for which no matching \tok{*/} is
  found before the end of the source file) is \illformed{}~(\diagcode{E100}).
\end{constraint}

\section{Tokens}
\label{sec:lex.tokens}
\index{token}

\pnum
The result of lexical analysis is a sequence of tokens.  Whitespace and
comments serve solely to separate tokens.  The grammar production for
\gramref{token} is given in Annex~A (\S\,\ref{gram:token}).

\pnum
Where the token stream is ambiguous, the lexer shall apply the
\textit{maximal munch} rule: at each position, the longest sequence of
characters that forms a valid token is consumed.

\begin{example}
\begin{laincode}
..=    // one token: range-inclusive, NOT '.' followed by '.='
>=     // one token: greater-than-or-equal, NOT '>' followed by '='
\end{laincode}
\end{example}

\pnum
Keywords (\S\,\ref{sec:lex.keywords}) take priority over identifiers: the
character sequence \texttt{func} at a point where an identifier is expected
is always the keyword \tok{func}.

\section{Whitespace}
\label{sec:lex.whitespace}
\index{whitespace}

\pnum
The whitespace characters are: space (U+0020) and horizontal tab (U+0009).
They serve only to separate adjacent tokens.

\pnum
Line feeds (U+000A), after normalization (\S\,\ref{sec:lex.encoding}), are
significant: they act as statement terminators in contexts where a statement
may follow (see \S\,\ref{sec:09-statements}).  Multiple consecutive line
feeds are treated as a single statement boundary.

\begin{constraint}
  A semicolon (\tok{;}) shall not appear in a Lain source file.  If one is
  present, the implementation shall emit~\diagcode{E100}.
\end{constraint}

\begin{rationale}
  Newlines as statement terminators eliminate the class of bugs caused by
  missing or accidental semicolons.  They also prevent multi-statement lines,
  keeping statement boundaries unambiguous.
\end{rationale}

\section{Keywords}
\label{sec:lex.keywords}
\index{keyword}

\pnum
The following identifiers are reserved as keywords.  They shall not be
used as user-defined identifiers.  The grammar production is
\gramref{keyword} (Annex~A).

\begin{longtable}{@{}lll@{}}
\toprule
\textbf{Keyword} & \textbf{Purpose} & \textbf{Clause} \\
\midrule
\endfirsthead
\multicolumn{3}{l}{\textit{(continued)}}\\
\toprule
\textbf{Keyword} & \textbf{Purpose} & \textbf{Clause} \\
\midrule
\endhead
\bottomrule
\endlastfoot
\tok{and}         & Logical AND operator           & \S\,\ref{sec:08-expressions} \\
\tok{as}          & Type cast                      & \S\,\ref{sec:08-expressions} \\
\tok{case}        & Pattern matching               & \S\,\ref{sec:15-match} \\
\tok{comptime}    & Compile-time evaluation        & \S\,\ref{sec:14-generics} \\
\tok{decreasing}  & Bounded-loop termination measure & \S\,\ref{sec:09-statements} \\
\tok{defer}       & Deferred execution             & \S\,\ref{sec:09-statements} \\
\tok{else}        & Alternative branch             & \S\,\ref{sec:09-statements} \\
\tok{false}       & Boolean literal                & \S\,\ref{sec:lex.literals} \\
\tok{for}         & Range-based iteration          & \S\,\ref{sec:09-statements} \\
\tok{func}        & Pure function declaration      & \S\,\ref{sec:12-functions} \\
\tok{if}          & Conditional                    & \S\,\ref{sec:09-statements} \\
\tok{import}      & Module import                  & \S\,\ref{sec:16-modules} \\
\tok{in}          & Range iteration / bounds guard & \S\,\ref{sec:08-expressions} \\
\tok{module}      & Module declaration             & \S\,\ref{sec:16-modules} \\
\tok{mov}         & Ownership transfer             & \S\,\ref{sec:11-ownership} \\
\tok{mut}         & Mutable borrow expression      & \S\,\ref{sec:11-ownership} \\
\tok{not}         & Logical NOT                    & \S\,\ref{sec:08-expressions} \\
\tok{or}          & Logical OR operator            & \S\,\ref{sec:08-expressions} \\
\tok{proc}        & Procedure declaration          & \S\,\ref{sec:12-functions} \\
\tok{ref}         & Shared borrow expression       & \S\,\ref{sec:11-ownership} \\
\tok{return}      & Return statement               & \S\,\ref{sec:09-statements} \\
\tok{struct}      & Struct type declaration        & \S\,\ref{sec:10-declarations} \\
\tok{enum}        & Enum type declaration          & \S\,\ref{sec:10-declarations} \\
\tok{true}        & Boolean literal                & \S\,\ref{sec:lex.literals} \\
\tok{unsafe}      & Unsafe block                   & \S\,\ref{sec:18-unsafe} \\
\tok{var}         & Mutable binding / mutable borrow & \S\,\ref{sec:10-declarations} \\
\tok{while}       & Loop statement                 & \S\,\ref{sec:09-statements} \\
\end{longtable}

\pnum
The following identifiers are \textit{reserved keywords}: they are
recognized by the lexer but have no current semantics.  They are reserved
for future use.

\begin{longtable}{@{}ll@{}}
\toprule
\textbf{Reserved keyword} & \textbf{Intended future purpose} \\
\midrule
\endfirsthead
\multicolumn{2}{l}{\textit{(continued)}}\\
\toprule
\textbf{Reserved keyword} & \textbf{Intended future purpose} \\
\midrule
\endhead
\bottomrule
\endlastfoot
\tok{end}     & Block terminator \\
\tok{export}  & Module visibility control \\
\tok{extern}  & C interoperability declarations \\
\tok{type}    & Type alias / type parameter \\
\end{longtable}

\begin{constraint}
  An identifier that is identical to a keyword or reserved keyword shall
  not be used in any context that requires an identifier.  Such use is
  \illformed{}~(\diagcode{E100}).
\end{constraint}

\section{Identifiers}
\label{sec:lex.identifiers}
\index{identifier}

\pnum
An identifier consists of ASCII letters, decimal digits, and underscores,
and shall begin with a letter or underscore.  The grammar production is
\gramref{identifier} (Annex~A).

\pnum
Identifiers are case-sensitive: \texttt{foo}, \texttt{Foo}, and
\texttt{FOO} are three distinct identifiers.

\pnum
A conforming implementation shall support identifiers of at least 255
characters.  The behavior when an identifier exceeds this length is
\idbehav{whether oversized identifiers are truncated or cause a diagnostic}.

\begin{constraint}
  An identifier shall not be identical to a keyword
  (\S\,\ref{sec:lex.keywords}).
\end{constraint}

\begin{constraint}
  An identifier shall begin with a letter or underscore, not a digit.
  A sequence of characters that begins with a digit is an integer
  literal~(\S\,\ref{sec:lex.literals}), not an identifier.
\end{constraint}

\section{Literals}
\label{sec:lex.literals}
\index{literal}

\subsection{Integer literals}
\label{sec:lex.literals.integer}
\index{integer literal}

\pnum
An integer literal denotes a constant integer value.  Four notations are
supported (grammar: \gramref{integer-literal}):

\begin{longtable}{@{}lll@{}}
\toprule
\textbf{Notation} & \textbf{Prefix} & \textbf{Example} \\
\midrule
\endfirsthead
\toprule
\textbf{Notation} & \textbf{Prefix} & \textbf{Example} \\
\midrule
\endhead
\bottomrule
\endlastfoot
Decimal   & (none) & \lain{42}, \lain{1_000_000} \\
Hexadecimal & \lain{0x} & \lain{0xFF}, \lain{0x1A_3B} \\
Binary    & \lain{0b} & \lain{0b1010_0101} \\
Octal     & \lain{0o} & \lain{0o755} \\
\end{longtable}

\pnum
An underscore (\texttt{\_}) may appear between any two digits of an integer
literal as a visual separator.  It has no effect on the value.

\pnum
The type of an integer literal is inferred from context (see
\S\,\ref{sec:07-types}).  In the absence of type context, the default type
is \lain{int}.

\begin{constraint}
  An integer literal whose value does not fit in the type determined by
  context is \illformed{}~(\diagcode{E100}).
\end{constraint}

\begin{example}
\begin{laincode}
var a = 255         // int, value 255
var b u8 = 255      // u8, value 255 -- OK
var c u8 = 256      // [E100]: 256 does not fit in u8
var d = 0xFF        // int, value 255
var e = 0b1111_0000 // int, value 240
var f = 1_000_000   // int, value 1000000
\end{laincode}
\end{example}

\subsection{Floating-point literals}
\label{sec:lex.literals.float}
\index{floating-point literal}

\pnum
A floating-point literal consists of a decimal integer part, a decimal
point, and a fractional part.  An optional exponent part
\texttt{e}\textit{N} or \texttt{E}\textit{N} may follow (grammar:
\gramref{float-literal}).

\pnum
The default type of a floating-point literal is \lain{f64}.

\begin{example}
\begin{laincode}
var x = 3.14      // f64
var y f32 = 1.0   // f32
var z = 2.5e3     // f64, value 2500.0
\end{laincode}
\end{example}

\subsection{Boolean literals}
\label{sec:lex.literals.bool}
\index{boolean literal}

\pnum
The boolean literals are \tok{true} and \tok{false} (grammar:
\gramref{bool-literal}).  Both have type \lain{bool}.

\subsection{Character literals}
\label{sec:lex.literals.char}
\index{character literal}

\pnum
A character literal is a single byte value enclosed in single quotes.  It
has type \lain{u8}.  Escape sequences (\S\,\ref{sec:lex.escape}) are
permitted.

\begin{example}
\begin{laincode}
var a = 'A'    // u8, value 65
var n = '\n'   // u8, value 10
var z = '\0'   // u8, value 0
\end{laincode}
\end{example}

\begin{constraint}
  A character literal shall contain exactly one character or one escape
  sequence.  An empty or multi-character character literal is
  \illformed{}~(\diagcode{E100}).
\end{constraint}

\subsection{String literals}
\label{sec:lex.literals.string}
\index{string literal}

\pnum
A string literal is a sequence of zero or more characters enclosed in
double quotes.  Escape sequences (\S\,\ref{sec:lex.escape}) are permitted
(grammar: \gramref{string-literal}).

\pnum
The type of a string literal of $N$ visible characters is
\lain{Fixed_u8_}\textit{N}\texttt{+1} (an array of $N{+}1$ bytes,
null-terminated).  This type is implicitly coercible to \lain{[u8]} (a
byte slice).

\pnum
The \texttt{.len} property of a string literal value equals $N$ (excluding
the null terminator).

\begin{example}
\begin{laincode}
var s = "hello"   // type: Fixed_u8_6 (5 chars + null)
                  // s.len == 5
\end{laincode}
\end{example}

\pnum
String literals are stored in the read-only data segment of the translation
output.  Modifying a string literal through a pointer is \undefbehav{the
result of writing through a pointer derived from a string literal}.

\subsection{Escape sequences}
\label{sec:lex.escape}
\index{escape sequence}

\pnum
The following escape sequences are recognized in string and character
literals:

\begin{longtable}{@{}lll@{}}
\toprule
\textbf{Escape} & \textbf{Byte value} & \textbf{Description} \\
\midrule
\endfirsthead
\toprule
\textbf{Escape} & \textbf{Byte value} & \textbf{Description} \\
\midrule
\endhead
\bottomrule
\endlastfoot
\texttt{\textbackslash n}  & 0x0A & Line feed (newline) \\
\texttt{\textbackslash t}  & 0x09 & Horizontal tab \\
\texttt{\textbackslash r}  & 0x0D & Carriage return \\
\texttt{\textbackslash 0}  & 0x00 & Null byte \\
\texttt{\textbackslash\textbackslash} & 0x5C & Backslash \\
\texttt{\textbackslash "} & 0x22 & Double quote \\
\texttt{\textbackslash '} & 0x27 & Single quote \\
\end{longtable}

\begin{constraint}
  A backslash not followed by one of the characters in the table above is
  an unrecognized escape sequence and is \illformed{}~(\diagcode{E100}).
\end{constraint}

\section{Operators and punctuators}
\label{sec:lex.operators}
\index{operator}

\pnum
The complete set of operator and punctuator tokens is listed below.  The
grammar productions \gramref{operator} and \gramref{punctuator} are in
Annex~A.

\subsection{Arithmetic operators}

\begin{longtable}{@{}ll@{}}
\toprule
\textbf{Token} & \textbf{Meaning} \\
\midrule
\endfirsthead
\toprule
\textbf{Token} & \textbf{Meaning} \\
\midrule
\endhead
\bottomrule
\endlastfoot
\tok{+}  & Addition; unary identity \\
\tok{-}  & Subtraction; unary negation \\
\tok{*}  & Multiplication \\
\tok{/}  & Division \\
\tok{\%} & Modulo (remainder) \\
\end{longtable}

\subsection{Comparison operators}

\begin{longtable}{@{}ll@{}}
\toprule
\textbf{Token} & \textbf{Meaning} \\
\midrule
\endfirsthead
\toprule
\textbf{Token} & \textbf{Meaning} \\
\midrule
\endhead
\bottomrule
\endlastfoot
\tok{==} & Equal \\
\tok{!=} & Not equal \\
\tok{<}  & Less than \\
\tok{>}  & Greater than \\
\tok{<=} & Less than or equal \\
\tok{>=} & Greater than or equal \\
\end{longtable}

\subsection{Logical operators}

\begin{longtable}{@{}ll@{}}
\toprule
\textbf{Token} & \textbf{Meaning} \\
\midrule
\endfirsthead
\toprule
\textbf{Token} & \textbf{Meaning} \\
\midrule
\endhead
\bottomrule
\endlastfoot
\tok{and} & Logical AND (short-circuit) \\
\tok{or}  & Logical OR (short-circuit) \\
\tok{not} & Logical NOT (prefix) \\
\tok{!}   & Logical NOT (alternate prefix form) \\
\end{longtable}

\subsection{Bitwise operators}

\begin{longtable}{@{}ll@{}}
\toprule
\textbf{Token} & \textbf{Meaning} \\
\midrule
\endfirsthead
\toprule
\textbf{Token} & \textbf{Meaning} \\
\midrule
\endhead
\bottomrule
\endlastfoot
\tok{\&}   & Bitwise AND \\
\tok{|}    & Bitwise OR \\
\tok{\^{}} & Bitwise XOR \\
\tok{<<}   & Left shift \\
\tok{>>}   & Right shift \\
\end{longtable}

\subsection{Assignment operators}

\begin{longtable}{@{}ll@{}}
\toprule
\textbf{Token} & \textbf{Meaning} \\
\midrule
\endfirsthead
\toprule
\textbf{Token} & \textbf{Meaning} \\
\midrule
\endhead
\bottomrule
\endlastfoot
\tok{=}    & Assignment \\
\tok{+=}   & Add and assign \\
\tok{-=}   & Subtract and assign \\
\tok{*=}   & Multiply and assign \\
\tok{/=}   & Divide and assign \\
\tok{\%=}  & Modulo and assign \\
\tok{\&=}  & Bitwise AND and assign \\
\tok{|=}   & Bitwise OR and assign \\
\tok{\^{}=}& Bitwise XOR and assign \\
\tok{<<=}  & Left shift and assign \\
\tok{>>=}  & Right shift and assign \\
\end{longtable}

\subsection{Other operators and punctuators}

\begin{longtable}{@{}ll@{}}
\toprule
\textbf{Token} & \textbf{Meaning} \\
\midrule
\endfirsthead
\toprule
\textbf{Token} & \textbf{Meaning} \\
\midrule
\endhead
\bottomrule
\endlastfoot
\tok{.}    & Member access; module path separator \\
\tok{..}   & Exclusive range (half-open interval) \\
\tok{->}   & Function return type separator \\
\tok{(}    & Open parenthesis \\
\tok{)}    & Close parenthesis \\
\tok{\{}   & Open brace \\
\tok{\}}   & Close brace \\
\tok{[}    & Open bracket \\
\tok{]}    & Close bracket \\
\tok{,}    & List separator \\
\tok{:}    & Type annotation; match arm separator \\
\tok{@}    & Attribute / annotation prefix \\
\end{longtable}

\section{Token disambiguation}
\label{sec:lex.disambiguation}
\index{token disambiguation}

\pnum
Several token sequences have context-dependent meaning that is resolved by
the parser, not the lexer.  The lexer produces the same token regardless of
context:

\begin{longtable}{@{}lll@{}}
\toprule
\textbf{Symbol} & \textbf{Context 1} & \textbf{Context 2} \\
\midrule
\endfirsthead
\toprule
\textbf{Symbol} & \textbf{Context 1} & \textbf{Context 2} \\
\midrule
\endhead
\bottomrule
\endlastfoot
\tok{*} & Multiplication (binary)            & (future: pointer dereference) \\
\tok{-} & Subtraction (binary)               & Negation (unary prefix) \\
\tok{.} & Member access                      & Module path separator \\
\tok{\&}& Borrowed match prefix              & Bitwise AND \\
\end{longtable}

\section{Integer arithmetic and overflow}
\label{sec:lex.overflow}
\index{integer overflow}

\pnum
Integer arithmetic in Lain uses two's complement representation for all
signed integer types.  Overflow wraps modulo $2^N$ where $N$ is the bit
width of the type.

\pnum
Unsigned integer arithmetic is modular (wrapping) by definition.

\pnum
A conforming implementation that emits C99 shall compile the translation
output with \texttt{-fwrapv} or equivalent, ensuring that signed integer
overflow is two's complement wrapping and not undefined behavior in the C
sense.

\begin{note}
  This differs from ISO C, where signed integer overflow is undefined
  behavior.  In Lain, signed overflow is defined as wrapping.  Programs
  that rely on wrapping behavior are well-formed.
\end{note}

% 06-basics.tex — Clause 6: Basic concepts
\chapter{Basic concepts}
\label{sec:06-basics}
\index{basic concepts}

\section{Declarations and definitions}
\label{sec:basics.declarations}
\index{declaration}
\index{definition}

\pnum
A \textit{declaration} introduces a name into a scope and associates it with
an entity: a variable, function, procedure, type, or module.  Every name
shall be declared before it is used, with the exception of names that are
resolved by Universal Function Call Syntax (\S\,\ref{sec:12-functions}).

\pnum
A \textit{definition} is a declaration that also allocates storage (for
variables) or provides a body (for functions and procedures).  In Lain,
every declaration at the local level is also a definition.

\begin{constraint}
  Redeclaring a name already visible in the current or an enclosing scope
  --- a duplicate parameter name, or a local that repeats or shadows an
  outer binding --- is \illformed{}~(\diagcode{E013}).  Lain forbids
  shadowing (\S\,\ref{sec:basics.scopes}).
\end{constraint}

\begin{constraint}
  Use of an undeclared identifier is \illformed{}~(\diagcode{E106}).  The
  reference implementation issues this diagnostic when the translation output is
  generated --- after name resolution, monomorphization, and UFCS desugaring are
  complete --- so that a name still unbound at that point is unambiguously
  undeclared.
\end{constraint}

\begin{example}
\begin{laincode}
proc f() {
    var x = 1
    var x = 2   // [E013]: redeclaration of 'x' is forbidden
}
\end{laincode}
\end{example}

\section{Scopes}
\label{sec:basics.scopes}
\index{scope}

\pnum
A \textit{scope} is a region of program text within which a declared name
can be used.  Lain has the following scope levels:

\begin{description}[leftmargin=2em]
  \item[Module scope]
    Names declared at the top level of a source file are in module scope.
    Module-scope names are accessible throughout the module and, by
    default, accessible from other modules that import this module.

  \item[Function/procedure scope]
    Parameter names are in the scope of the function or procedure body.

  \item[Block scope]
    Names declared inside a block (\lain{\{ \}}) are in the scope of that
    block.  They are not accessible outside the enclosing block.
\end{description}

\pnum
Scopes are \textit{lexically nested}.  Lain \emph{forbids shadowing}
entirely: a declaration whose name is already visible in the current scope
or in any enclosing scope --- a function parameter shadowed by a local, or a
name redeclared in a nested block --- is \illformed{}~(\diagcode{E013}).

\begin{rationale}
  Disallowing shadowing removes a class of bugs in which an inner binding
  silently captures a name the reader expects to denote an outer one.  Every
  identifier in a function body therefore resolves to exactly one
  declaration.
\end{rationale}

\pnum
A conforming implementation shall support at least 64 levels of nested
block scope (\S\,\ref{sec:conformance.limits}).

\section{Name resolution}
\label{sec:basics.resolution}
\index{name resolution}

\pnum
Name resolution maps an identifier in source text to its declaration.
The lookup algorithm is:

\begin{enumerate}
  \item Search the innermost block scope.
  \item If not found, search progressively wider enclosing block scopes.
  \item If not found, search the function/procedure scope (parameters).
  \item If not found, search the module scope.
  \item If not found, search imported module scopes (in import order).
  \item If not found, attempt UFCS resolution (\S\,\ref{sec:12-functions}).
  \item If still not found, the identifier is undeclared and the program is
        \illformed{}~(\diagcode{E106}).
\end{enumerate}

\pnum
Because shadowing is forbidden (\S\,\ref{sec:basics.scopes}), a name resolves
to at most one declaration along the scope chain: a lookup that would find a
second binding of the same name is rejected at the point of the redeclaration
(\diagcode{E013}), not silently resolved by an innermost-wins rule.

\section{Storage duration}
\label{sec:basics.storage}
\index{storage duration}

\pnum
Every object has a \textit{storage duration} that determines the lifetime
of its storage.

\begin{description}[leftmargin=2em]
  \item[Automatic storage duration]
    Variables declared inside a block or function body.  Storage is
    allocated on entry to the block and deallocated on exit.

  \item[Static storage duration]
    Variables declared at module scope.  Storage persists for the lifetime
    of the program.  \idbehav{whether module-scope variables are zero-
    initialized before \lain{proc main} is called}.

  \item[Arena storage duration]
    Objects allocated from an explicit \lain{Arena}.  Storage persists
    until the arena is freed as a whole.
\end{description}

\section{Objects and values}
\label{sec:basics.objects}
\index{object}

\pnum
An \textit{object} is a region of storage interpreted according to a type.
Every declared variable corresponds to an object.

\pnum
Every variable has:
\begin{itemize}
  \item a \textit{declared type} (\S\,\ref{sec:07-types});
  \item an \textit{ownership mode}: owned (default for linear types),
        \lain{shared} (read-only reference), or \lain{var} (mutable
        reference) (\S\,\ref{sec:11-ownership});
  \item a \textit{linear state}: Uninit, Unconsumed, or Consumed (for
        linear types only) (\S\,\ref{sec:11-ownership}).
\end{itemize}

\section{Lvalues and rvalues}
\label{sec:basics.lvalues}
\index{lvalue}
\index{rvalue}

\pnum
An \textit{lvalue} is an expression that designates an object with a
persistent storage location: a variable identifier, a field access on an
lvalue, or an array index expression on an lvalue.

\pnum
An \textit{rvalue} is an expression whose value has no persistent storage
location: literals, the result of arithmetic operations, or the result of
a function call.

\pnum
Assignment requires the left-hand side to be an lvalue.

\begin{constraint}
  An assignment whose left-hand side is an rvalue is
  \illformed{}~(\diagcode{E100}).
\end{constraint}

\section{Alignment}
\label{sec:basics.alignment}
\index{alignment}

\pnum
Every type has a required alignment.  The alignment of a type is
\idbehav{the same as the alignment of the corresponding C99 type}.
Struct layout follows C99 struct layout rules, including padding
(\S\,\ref{sec:07-types}).

\pnum
Stack-allocated variables are aligned to their type's required alignment.
The implementation shall not insert additional padding between local
variables unless required by the hardware.

% 07-types.tex — Clause 7: Types
\chapter{Types}
\label{sec:07-types}
\index{type}

\pnum
Lain is statically typed.  Every variable, expression, and parameter has a
type determined at compile time.  There is no runtime type information
(RTTI).  The grammar productions for types are in Annex~A (\gramref{type}).

\section{Type taxonomy}
\label{sec:types.taxonomy}
\index{type!taxonomy}

\pnum
The types of Lain are organized as follows:

\begin{itemize}
  \item \textit{Primitive types}: integer, floating-point, boolean, void
        (\S\,\ref{sec:types.primitive});
  \item \textit{Array types}: fixed-size, compile-time-sized
        (\S\,\ref{sec:types.array});
  \item \textit{Slice types}: dynamically-sized views into contiguous memory,
        including sized slice types with static length constraints
        (\S\,\ref{sec:types.slice});
  \item \textit{Struct types}: named product types (\S\,\ref{sec:types.struct});
  \item \textit{Enum types}: simple enumerations (\S\,\ref{sec:types.enum});
  \item \textit{ADT types}: sum types with associated data (\S\,\ref{sec:types.adt});
  \item \textit{Pointer types}: raw memory addresses, unsafe only
        (\S\,\ref{sec:types.pointer});
  \item \textit{Constrained types}: types with value range annotations
        (\S\,\ref{sec:types.constrained}).
\end{itemize}

\pnum
Types are classified as \textit{linear} or \textit{unrestricted}.  Linear
types must be consumed exactly once (see \S\,\ref{sec:11-ownership}).
Unrestricted types may be used any number of times.

\pnum
All primitive types are unrestricted.  A struct type is linear if and only
if at least one of its fields is of a linear type.  This property is
transitive.

\section{Primitive types}
\label{sec:types.primitive}
\index{primitive type}

\subsection{Integer types}
\label{sec:types.integer}
\index{integer type}

\pnum
The integer types are listed below.  Sizes marked \textit{(platform)} are
\idbehav{determined by the target platform as in the C99 standard}.

\begin{longtable}{@{}lllll@{}}
\toprule
\textbf{Type} & \textbf{Signedness} & \textbf{Bit width} & \textbf{Range} & \textbf{C99 type} \\
\midrule
\endfirsthead
\multicolumn{5}{l}{\textit{(continued)}}\\
\toprule
\textbf{Type} & \textbf{Signedness} & \textbf{Bit width} & \textbf{Range} & \textbf{C99 type} \\
\midrule
\endhead
\bottomrule
\endlastfoot
\lain{u8}    & unsigned & 8    & $[0,\, 2^8-1]$       & \texttt{uint8\_t}  \\
\lain{u16}   & unsigned & 16   & $[0,\, 2^{16}-1]$    & \texttt{uint16\_t} \\
\lain{u32}   & unsigned & 32   & $[0,\, 2^{32}-1]$    & \texttt{uint32\_t} \\
\lain{u64}   & unsigned & 64   & $[0,\, 2^{64}-1]$    & \texttt{uint64\_t} \\
\lain{usize} & unsigned & (platform) & $[0,\, 2^N-1]$ & \texttt{size\_t}   \\
\lain{i8}    & signed   & 8    & $[-2^7,\, 2^7-1]$    & \texttt{int8\_t}   \\
\lain{i16}   & signed   & 16   & $[-2^{15},\, 2^{15}-1]$ & \texttt{int16\_t} \\
\lain{i32}   & signed   & 32   & $[-2^{31},\, 2^{31}-1]$ & \texttt{int32\_t} \\
\lain{i64}   & signed   & 64   & $[-2^{63},\, 2^{63}-1]$ & \texttt{int64\_t} \\
\lain{isize} & signed   & (platform) & $[-2^{N-1},\, 2^{N-1}-1]$ & \texttt{ptrdiff\_t} \\
\lain{int}   & signed   & 32 & $[-2^{31},\, 2^{31}-1]$ & \texttt{int32\_t} \\
\end{longtable}

\pnum
\lain{int} is a documented alias of \lain{i32}, identical in every
semantic respect (boundary checks, niche optimization, codegen). Use
either spelling — \lain{int} when bit-width is incidental, \lain{i32}
when explicit width matters.

\pnum
Generalized integer types \lain{iN}/\lain{uN} are admitted for any
$N \in [1, 64]$. Hardware-aligned widths
($N \in \{1, 8, 16, 32, 64\}$) map to standard C99 stdint types.
Logical widths (other $N$) occupy the smallest hardware-aligned
container that fits $N$ bits when standalone; in a \tok{[packed]}
struct they occupy exactly $N$ bits.

\pnum
\lain{usize} and \lain{isize} are the platform-native unsigned and signed
integer types whose width equals the pointer width.  Their exact width is
\idbehav{the pointer width on the target platform}.

\subsection{Hardware-aligned vs logical integer types}
\label{sec:types.iN-categories}
\index{integer type!hardware-aligned}
\index{integer type!logical}

\pnum
The \lain{iN}/\lain{uN} types fall into two semantic categories:

\textbf{Hardware-aligned} (\(N \in \{1, 2, 4, 8, 16, 32, 64\}\), with
the common subset being \(\{8, 16, 32, 64\}\)): these correspond to
fundamental C99 stdint widths (\texttt{uint8\_t}, \texttt{int32\_t},
etc.). They are the canonical choice for C interop and for any context
where exact storage size matters at the byte boundary.

\textbf{Logical} (other \(N\)): these are \textit{bit-width quantities}.
Their semantics are context-dependent:

\begin{itemize}
  \item \textbf{Standalone}: storage uses the smallest hardware-aligned
        container fitting \(N\) bits (e.g., \lain{u12} maps to
        \texttt{uint16\_t}). The natural range is
        \([-2^{N-1},\,2^{N-1}-1]\) signed or \([0,\,2^N-1]\) unsigned.
  \item \textbf{Inside a packed struct} (future feature,
        \S\,\ref{sec:types.packed-future}): the type occupies exactly
        \(N\) bits of storage with bit-exact layout.
\end{itemize}

\pnum
\textbf{Distinction from refinement.}  Logical \lain{iN}/\lain{uN} types
differ from a refinement constraint on \lain{int}:

\begin{itemize}
  \item \lain{int >= 0 and <= 7} declares an \lain{int} with a range
        constraint. Storage is the \lain{int} default (\texttt{int32\_t});
        bit-width is not specified.
  \item \lain{u3} declares a 3-bit quantity. Standalone its range is
        \([0, 7]\) and its storage is \texttt{uint8\_t}; in a packed
        struct (future) it occupies exactly 3 bits.
\end{itemize}

The two forms are interchangeable when only the range matters; they
differ when bit-exact storage layout matters.

\begin{example}
\begin{laincode}
// Logical bit-width: 12-bit quantity, container = uint16_t.
var x u12 = 4095

// Refinement constraint on int: range [0, 4095], container = int32_t.
var y int >= 0 and <= 4095 = 4095

// Refinement type alias: a named refinement type (\S\,\ref{sec:types.alias-refinement}).
type Pressure = int >= 0 and <= 1000
var p Pressure = 500
\end{laincode}
\end{example}

\subsection{Refinement type aliases}
\label{sec:types.alias-refinement}
\index{refinement type}
\index{type alias!refinement}

\pnum
A type alias may carry a refinement constraint that narrows the underlying
type's value range:

\begin{laincode}
type Percent = int >= 0 and <= 100
type Pressure = int >= 0 and <= 1000
type NonZero = int != 0
\end{laincode}

\pnum
Variables declared with a refinement type alias inherit its constraint.
Assignments whose source value range violates the alias constraint are
\illformed{}~(\diagcode{E086}).

\begin{example}
\begin{laincode}
type Percent = int >= 0 and <= 100

proc main() int {
    var ok Percent = 75              // [E086 not raised]: 75 fits [0,100]
    // var bad Percent = 150         // [E086]: 150 > 100
    return 0
}
\end{laincode}
\end{example}

\subsection{Packed structs}
\label{sec:types.packed-future}
\index{packed struct}
\index{bitfield}

\pnum
The \tok{[packed]} attribute (\S\,\ref{sec:decl.attributes}) on a struct
declaration enables bit-exact memory layout. Each \lain{iN}/\lain{uN}
field occupies exactly \(N\) bits within a single scalar container.

\begin{laincode}
[packed]
type GpioModer {
    pin0 u2     // bit 0-1
    pin1 u2     // bit 2-3
    pin2 u2     // bit 4-5
    pin3 u2     // bit 6-7
}
// total: 8 bits -> uint8_t container, bit-exact layout
\end{laincode}

\pnum
\textbf{Layout rules}:
\begin{itemize}
  \item Fields are laid out in declaration order, starting at bit 0.
  \item Each field occupies exactly its declared bit width.
  \item Total bit width must be \(\leq 64\); container is the smallest
        of \texttt{uint8\_t}, \texttt{uint16\_t}, \texttt{uint32\_t}, or
        \texttt{uint64\_t} that fits the total.
\end{itemize}

\pnum
\textbf{Access semantics}:
\begin{itemize}
  \item Construction: \lain{GpioModer(1, 2, 3, 0)} packs the four
        2-bit fields into the container.
  \item Read: \lain{r.pin1} extracts bits 2-3 via shift-and-mask.
  \item Write: \lain{r.pin1 = v} updates bits 2-3, leaving other
        fields unchanged.
\end{itemize}

\begin{constraint}
  Inside a \tok{[packed]} struct, fields shall be \lain{iN}/\lain{uN}
  types only. Nested structs, pointers, and slices are not permitted.
  Violation is \illformed{}~(\diagcode{E121}).
\end{constraint}

\begin{constraint}
  The total bit width of fields in a \tok{[packed]} struct shall not
  exceed 64. Otherwise \illformed{}~(\diagcode{E121}).
\end{constraint}

\begin{note}
  This feature targets embedded use cases: hardware registers,
  packed network protocols, codec implementations. The reference
  implementation emits a scalar typedef plus inline getter/setter
  functions for each field.
\end{note}

\subsection{Integer rank and widening}
\label{sec:types.rank}
\index{integer rank}
\index{integer widening}

\pnum
The integer types are partially ordered by \textit{rank}:

\begin{longtable}{@{}lll@{}}
\toprule
\textbf{Rank} & \textbf{Types} & \textbf{Bit width} \\
\midrule
\endfirsthead
\toprule
\textbf{Rank} & \textbf{Types} & \textbf{Bit width} \\
\midrule
\endhead
\bottomrule
\endlastfoot
1 & \lain{u8}, \lain{i8}             & 8 \\
2 & \lain{u16}, \lain{i16}           & 16 \\
3 & \lain{u32}, \lain{i32}, \lain{int} & 32 \\
4 & \lain{u64}, \lain{i64}           & 64 \\
5 & \lain{usize}, \lain{isize}       & platform \\
\end{longtable}

\pnum
\textbf{Implicit widening}: a value of a lower-rank integer type may be
used where a higher-rank integer type is expected, without an explicit cast.
The result type of a binary arithmetic or bitwise operation is the
higher-rank operand type.

\begin{example}
\begin{laincode}
var a u8  = 42
var b int = a          // u8 -> int: implicit widening, OK
var c = a + 1000       // u8 + int -> int (wider wins)
\end{laincode}
\end{example}

\begin{constraint}
  An implicit conversion from a higher-rank integer type to a lower-rank
  integer type (narrowing) is \illformed{}.  Narrowing requires an explicit
  \lain{as} cast (\S\,\ref{sec:08-expressions}).
\end{constraint}

\begin{constraint}
  An implicit conversion between integer types and floating-point types is
  \illformed{}.  Such conversions require an explicit \lain{as} cast.
\end{constraint}

\begin{constraint}
  An implicit conversion between integer types and \lain{bool} is
  \illformed{}.  Such conversions require an explicit \lain{as} cast.
\end{constraint}

\subsection{Floating-point types}
\label{sec:types.float}
\index{floating-point type}

\pnum
\lain{f32} and \lain{f64} are 32-bit and 64-bit floating-point types
conforming to IEC~60559 (IEEE~754).  They correspond to C's \texttt{float}
and \texttt{double} respectively.  \lain{float} is a documented alias
of \lain{f32}.

\pnum
The default type of a floating-point literal is \lain{f64}
(\S\,\ref{sec:lex.literals.float}).

\pnum
Arithmetic on \lain{f32} and \lain{f64} follows IEEE~754 semantics:
overflow produces $\pm\infty$; the result of $0.0/0.0$ and similar
operations is NaN.  These are defined behaviors in Lain.

\subsection{Boolean type}
\label{sec:types.bool}
\index{boolean type}

\pnum
\lain{bool} has exactly two values: \lain{true} and \lain{false}.
\lain{bool} is stored as a single byte; its exact representation is
\idbehav{whether \lain{true} is 1 or any non-zero value}.

\begin{constraint}
  There are no implicit \textbf{value} conversions between \lain{bool}
  and any integer type for general arithmetic or assignment.
\end{constraint}

\pnum
The condition expression of \lain{if} and \lain{while} accepts either
a value of type \lain{bool} or any integer type.  When an integer
value is used in condition position, zero is treated as \lain{false}
and any non-zero value is treated as \lain{true}.  This relaxation is
ergonomic only and does not extend to other contexts.

\begin{example}
\begin{laincode}
var n int = read_count()
if n { ... }        // OK: integer in condition, non-zero is true
var b bool = n      // [E100]: not a value conversion
\end{laincode}
\end{example}

\subsection{The void type}
\label{sec:types.void}
\index{void type}

\pnum
\lain{void} represents the absence of a value.  It is used exclusively
in pointer types (\lain{*void}) for C interoperability (see
\S\,\ref{sec:17-interop}).

\begin{constraint}
  A variable declaration with type \lain{void} is \illformed{}.
\end{constraint}

\section{Array types}
\label{sec:types.array}
\index{array type}

\pnum
An array type \lain{[T; N]} denotes a contiguous sequence of $N$ values of
type~\lain{T}.  $N$ shall be a compile-time constant integer greater than
zero (grammar: \gramref{array-type}).

\pnum
The \texttt{.len} property of an array type yields $N$ as a compile-time
constant of type \lain{int}.

\pnum
Array elements are laid out contiguously in memory in increasing index
order, with element~0 at the base address.  The stride between elements is
\idbehav{equal to \texttt{sizeof(T)} in C99, including any intra-element
padding}.

\begin{constraint}
  An array size $N \leq 0$ is \illformed{}~(\diagcode{E100}).
\end{constraint}

\begin{constraint}
  Every array index expression shall satisfy
  $0 \leq \textit{idx} < \textit{arr.len}$.
  An access for which value range analysis cannot prove this is a
  constraint violation~(\diagcode{E014}).
\end{constraint}

\begin{example}
\begin{laincode}
var arr int[5] = [1, 2, 3, 4, 5]
var x = arr[2]         // OK: 2 is statically in [0, 4]
// var y = arr[5]      // [E014]: 5 >= arr.len
\end{laincode}
\end{example}

\section{Slice types}
\label{sec:types.slice}
\index{slice type}

\pnum
A slice type \lain{*T[]} is a reference to a contiguous sequence of values
of type~\lain{T} whose length is determined at runtime.  Its runtime
representation is a pair \texttt{(ptr: *T, len: usize)}.  The \lain{*}
prefix is mandatory: all view/reference types use \lain{*} consistently.

\pnum
The full family of reference-to-array types:

\begin{center}
\begin{tabular}{@{}lll@{}}
\toprule
\textbf{Syntax} & \textbf{Semantics} & \textbf{Runtime cost} \\
\midrule
\lain{T[N]}      & inline array (N elements owned) & N $\times$ sizeof(T) \\
\lain{*T[]}      & view, runtime length             & ptr + usize \\
\lain{*T[N]}     & view, compile-time length        & ptr only \\
\lain{*T[:S]}    & view, sentinel-terminated         & ptr only \\
\lain{*T}        & raw pointer (no length)           & ptr only \\
\bottomrule
\end{tabular}
\end{center}

\pnum
A slice has two properties accessible via dot notation:
\begin{itemize}
  \item \texttt{.len}: the number of elements, type \lain{usize};
  \item \texttt{.data}: a pointer to the first element, type \lain{*T}.
\end{itemize}

\pnum
A string literal of $N$ characters has static type \lain{Fixed\_u8\_{N+1}}
(a null-terminated fixed-size array of $N{+}1$ bytes) and is implicitly
coercible to \lain{*u8[]}.

\subsection{Sized slice types}
\label{sec:types.sized-slice}
\index{sized slice type}
\index{slice type!sized}

\pnum
A \textit{sized slice type} is a slice type annotated with a size constraint
on its length.  The constraint is part of the parameter or variable
declaration, not of the nominal type itself.  The syntax is:

\begin{center}
\begin{tabular}{@{}ll@{}}
\toprule
\textbf{Syntax} & \textbf{Meaning} \\
\midrule
\lain{*T[]}         & unsized dynamic slice (no length constraint) \\
\lain{*T[n]}        & length exactly \lain{n}; \textit{n} is a size expression \\
\lain{*T[>= n]}     & length $\geq$ \lain{n}; \textit{n} is a size expression \\
\lain{*T[> n]}      & length $>$ \lain{n}; \textit{n} is a size expression \\
\bottomrule
\end{tabular}
\end{center}

\pnum
A \textit{size expression} (\textit{size\_expr}) may be any of:

\begin{itemize}
  \item an integer literal (e.g., \lain{3}, \lain{0});
  \item an identifier that names a numeric variable in scope (e.g., \lain{n});
  \item a member expression of the form \lain{p.len}, where \lain{p} is a
        slice or array parameter in scope (e.g., \lain{a.len}, \lain{src.len});
  \item a binary arithmetic expression over the above
        (e.g., \lain{a.len + b.len}, \lain{src.len - 1}).
\end{itemize}

\begin{example}
\begin{laincode}
// Exact-length constraint: out must have length == a.len + b.len
func concat(a i32[], b i32[], out i32[a.len + b.len]) { ... }

// Lower-bound constraint: arr must have at least n elements
func first(arr i32[> 0]) i32 { return arr[0] }

// Derived constraint using member .len
func slide(src i32[], out i32[src.len - 1]) { ... }
\end{laincode}
\end{example}

\pnum
\textbf{Constraint semantics.}
\begin{itemize}
  \item \lain{T[expr]}: the value range analysis (VRA, \S\,\ref{sec:13-vra})
        records the constraint \texttt{len == expr} and uses it to prove
        subsequent bounds checks inside the function body.
  \item \lain{T[>= expr]}: VRA records \texttt{len >= expr}.
  \item \lain{T[> expr]}: VRA records \texttt{len > expr}, i.e.\ \texttt{len >= expr + 1}.
\end{itemize}

\pnum
\textbf{Call-site verification.}
At every call site, the compiler verifies that the argument satisfies the
formal parameter's size constraint (\diagcode{E087}); see
\S\,\ref{sec:vra.sized-slice-callsite}.

\pnum
\textbf{C ABI decomposition (Fase~7).}
When a function parameter has a sized or unsized dynamic slice type, the
compiler decomposes it into a pair of C parameters:

\begin{itemize}
  \item a length parameter \texttt{size\_t \_\_len\_NAME};
  \item a data pointer \texttt{const T*} (shared) or \texttt{T* restrict} (mutable).
\end{itemize}

If the parameter's size expression is an arithmetic expression over other
parameters' lengths (e.g., \lain{out i32[a.len + b.len]}), no
\texttt{\_\_len\_out} parameter is emitted; the length is re-derived from
\texttt{\_\_len\_a + \_\_len\_b} at every use site.

\begin{example}
\begin{laincode}
// Source:
proc vadd(var out i32[], a i32[out.len], b i32[out.len]) { ... }

// C emission:
// void vadd(size_t __len_out, int32_t * restrict out,
//           const int32_t* a, const int32_t* b);
\end{laincode}
\end{example}

\section{Struct types}
\label{sec:types.struct}
\index{struct type}

\pnum
A struct type groups named fields into a single named product type.  The
grammar production is \gramref{struct-declaration} (Annex~A).

\pnum
Field order is the declaration order.  The memory layout of a struct is
\idbehav{equal to the C99 struct layout, including compiler-inserted padding
for alignment}.

\pnum
A struct type~$S$ is linear if at least one field of $S$ has a linear type
(transitively).

\begin{constraint}
  \textbf{Q-008.}
  Every field of a struct or enum variant whose declared type is linear
  shall be annotated \lain{mov}.  Omitting \lain{mov} on a linear field
  is \illformed{}~(\diagcode{E083}).  This requirement makes linearity
  locally visible in the declaration: a reader of a struct definition
  does not need to inspect the field's type definition to know whether
  the struct is linear.
\end{constraint}

\begin{example}
\begin{laincode}
type Handle {
    mov ptr *int            // explicit `mov` on linear field
}

type Session {
    mov h Handle            // OK: Handle is linear, mov required
    name [u8]               // unrestricted field, no mov needed
}

// Ill-formed: missing mov on a linear field
type Bad {
    h Handle                // [E083]
}
\end{laincode}
\end{example}

\pnum
\textbf{Struct construction}: all fields shall be provided in declaration
order.  The type of each argument shall be compatible with the declared
field type (\S\,\ref{sec:types.compat}).

\begin{constraint}
  A struct constructor expression with a number of arguments different from
  the number of fields is \illformed{}~(\diagcode{E100}).
\end{constraint}

\begin{constraint}
  Writing to a field of a struct variable requires the variable to be
  declared as \lain{var} (mutable).  Writing to an immutable binding's
  field is \illformed{}.
\end{constraint}

\begin{constraint}
  Using \lain{==} or \lain{!=} on a struct type is \illformed{}.
  Field-by-field comparison must be performed explicitly.
\end{constraint}

\section{Enum types}
\label{sec:types.enum}
\index{enum type}

\pnum
An enum type is a type with a finite, named set of values.  Each variant
carries no associated data.  The grammar production is
\gramref{enum-declaration} (Annex~A).

\pnum
The underlying representation of an enum type is \idbehav{an integer type
large enough to hold all variant tags; the reference implementation uses
\texttt{int}}.

\pnum
Enum values are accessed as \lain{TypeName.VariantName}.

\begin{constraint}
  Using \lain{==} or \lain{!=} to compare enum values is \illformed{}.
  Enum values shall be compared using \lain{case} pattern matching
  (\S\,\ref{sec:15-match}).
\end{constraint}

\begin{note}
  The constraint on \lain{==} for enums reflects that the underlying tag
  representation is implementation-defined.  Pattern matching is always
  sound.
\end{note}

\section{Algebraic data types (ADTs)}
\label{sec:types.adt}
\index{algebraic data type}

\pnum
An ADT is a sum type in which each variant may carry associated fields.
A type declaration is classified as an ADT if at least one variant has
fields.  The grammar production is \gramref{enum-declaration} (Annex~A,
using the record-variant form).

\pnum
ADTs use a niche-based representation whenever a zero-cost encoding is
available; otherwise they fall back to a tagged union (a tag field of
implementation-defined integer type plus a union of variant payload
structs).

\begin{note}
  \textbf{D-Niche aggressive multi-slot.}
  When the variant payload admits a \textit{niche}---a bit pattern that
  cannot occur in any valid value of the payload type---that niche is
  used to encode tag-less variants.  Niche sources, in priority order:
  zero-page pointer slots for pointers/slices; VRA-proven unused values
  for constrained integers; bit patterns outside the declared variants
  for \lain{bool} or small enums; niche of any recursive field. See
  \S\,\ref{sec:mem.adt-layout} for the full algorithm.

  If no niche or only a partial niche can be derived, the compiler
  emits \diagcode{W120} (best-effort warning, never an error) and
  falls back to a tag+union representation. The programmer may
  resolve the warning by constraining the payload, by changing to a
  payload type with larger sentinel space, or by accepting the tag
  byte. See \S\,\ref{sec:mem.adt-layout} for the full layout
  algorithm and target dependencies.
\end{note}

\pnum
ADT values are constructed as \lain{TypeName.VariantName(args)} for
variants with fields, or \lain{TypeName.VariantName} for unit variants.

\pnum
ADT values are destructured via \lain{case} pattern matching
(\S\,\ref{sec:15-match}).

\section{Pointer types}
\label{sec:types.pointer}
\index{pointer type}

\pnum
A pointer type \lain{*T} or \lain{const *T} is a raw memory address.
Pointer operations require an \lain{unsafe} block (\S\,\ref{sec:18-unsafe}).

\pnum
The ownership modes for pointer types are:

\begin{longtable}{@{}llll@{}}
\toprule
\textbf{Lain syntax} & \textbf{Mode} & \textbf{C99 type} & \textbf{Linear?} \\
\midrule
\endfirsthead
\toprule
\textbf{Lain syntax} & \textbf{Mode} & \textbf{C99 type} & \textbf{Linear?} \\
\midrule
\endhead
\bottomrule
\endlastfoot
\lain{*T}     & shared (read-only) & \texttt{const T*}  & no \\
\lain{var *T} & mutable            & \texttt{T*}        & no \\
\lain{mov *T} & owned              & \texttt{T*}        & yes \\
\end{longtable}

\begin{constraint}
  Dereferencing a pointer or taking the address of a variable shall occur
  only inside an \lain{unsafe} block.  Pointer operations in safe code are
  \illformed{}~(\diagcode{E012}).
\end{constraint}

\section{Constrained types}
\label{sec:types.constrained}
\index{constrained type}

\pnum
A constrained type is a type annotated with a value range constraint:

\begin{center}
\textit{T} \textit{op} \textit{bound}
\end{center}

where \textit{op} $\in \{$\lain{>=}, \lain{<=}, \lain{>}, \lain{<},
\lain{==}$\}$ and \textit{bound} is an integer literal.

\pnum
Constrained types are used in function return types and parameter types to
declare value range invariants that the value range analysis
(\S\,\ref{sec:13-vra}) shall enforce.

\begin{example}
\begin{laincode}
func nonzero_divisor() int >= 1 {
    return 42   // VRA verifies: 42 >= 1 holds
}
\end{laincode}
\end{example}

\pnum
The constraint is part of the function signature, not the type system's
nominal identity.  Two constrained types with the same base type are
assignment-compatible.

\section{Type compatibility and equivalence}
\label{sec:types.compat}
\index{type equivalence}
\index{type compatibility}

\pnum
Lain uses \textbf{nominal type equivalence}: two types are the same if and
only if they have the same declared name.

\pnum
Two types $T_1$ and $T_2$ are \textit{compatible} for an expression context
if:
\begin{itemize}
  \item $T_1 = T_2$ (same nominal type), or
  \item $T_1$ is an integer type of lower rank than $T_2$ (implicit
        widening; \S\,\ref{sec:types.rank}).
\end{itemize}

\begin{constraint}
  An expression of type $T_1$ used in a context requiring type $T_2$, where
  $T_1$ and $T_2$ are not compatible, is \illformed{}.
\end{constraint}

\pnum
The operators \lain{==} and \lain{!=} are defined for:
\begin{itemize}
  \item primitive numeric types;
  \item \lain{bool};
  \item pointer types.
\end{itemize}
They are not defined for struct types, array types, or enum types.

% 08-expressions.tex — Clause 8: Expressions
\chapter{Expressions}
\label{sec:08-expressions}
\index{expression}

\pnum
An expression is a syntactic construct that evaluates to a value.  Every
expression has a type, determined statically at compile time.  The grammar
production for \gramref{expression} is in Annex~A.

\pnum
Expressions are evaluated strictly left-to-right: for any binary operator
$a\ \mathit{op}\ b$ or function call \lain{f(a, b)}, the left operand
(or earlier argument) is fully evaluated before the right operand (or
later argument).

\section{Operator precedence and associativity}
\label{sec:expr.precedence}
\index{operator precedence}

\pnum
The following table lists all binary operators in order of decreasing
precedence.  All operators listed on the same row have equal precedence.

\begin{longtable}{@{}llll@{}}
\toprule
\textbf{Prec.} & \textbf{Operators} & \textbf{Assoc.} & \textbf{Result type} \\
\midrule
\endfirsthead
\toprule
\textbf{Prec.} & \textbf{Operators} & \textbf{Assoc.} & \textbf{Result type} \\
\midrule
\endhead
\bottomrule
\endlastfoot
1 (high) & \lain{()} \lain{[]} \lain{.}              & left  & (context-dependent) \\
2        & unary \lain{-} \lain{!} \lain{not} \lain{mov} \lain{ref} \lain{mut} & right & (operand type) \\
3        & \lain{as}                                  & left  & target type \\
4        & \lain{*} \lain{/} \lain{\%}                & left  & wider operand type \\
5        & \lain{+} \lain{-}                          & left  & wider operand type \\
6        & \lain{<<} \lain{>>}                        & left  & left operand type \\
7        & \lain{\&}                                  & left  & wider operand type \\
8        & \texttt{\textasciicircum}                                & left  & wider operand type \\
9        & \lain{|}                                   & left  & wider operand type \\
10       & \lain{<} \lain{>} \lain{<=} \lain{>=} \lain{in} & left & \lain{bool} \\
11       & \lain{==} \lain{!=}                        & left  & \lain{bool} \\
12       & \lain{and}                                 & left  & \lain{bool} \\
13 (low) & \lain{or}                                  & left  & \lain{bool} \\
\end{longtable}

\pnum
Parentheses override precedence.

\section{Primary expressions}
\label{sec:expr.primary}
\index{primary expression}

\pnum
A primary expression is one of: an identifier, a literal, a parenthesized
expression, a block expression, an \lain{if} expression, a \lain{case}
expression, a constructor expression, or a \lain{comptime} expression.

\pnum
\textbf{Identifier expression}: evaluates to the current value of the named
variable.

\begin{constraint}
  The identifier shall be declared in scope
  (\diagcode{E013}).
\end{constraint}

\begin{constraint}
  The identifier shall have been initialized before use
  (\diagcode{E005}).
\end{constraint}

\section{Arithmetic operators}
\label{sec:expr.arith}
\index{arithmetic operator}

\pnum
The binary arithmetic operators are \lain{+} (addition), \lain{-}
(subtraction), \lain{*} (multiplication), \lain{/} (division), and
\lain{\%} (remainder).  The result type is the wider of the two operand
types (\S\,\ref{sec:types.rank}).

\pnum
Integer division truncates toward zero.  \lain{f32}/\lain{f64} division
follows IEEE~754.

\begin{constraint}
  Integer division or remainder where the divisor is statically provable to
  be zero (via value range analysis, \S\,\ref{sec:13-vra}) is
  \illformed{}~(\diagcode{E015}).
\end{constraint}

\begin{constraint}
  \lain{\%} shall not be applied to floating-point operands.
\end{constraint}

\pnum
The unary \lain{-} operator negates its operand.  For unsigned integer
types the result wraps modulo $2^N$.

\section{Comparison operators}
\label{sec:expr.compare}
\index{comparison operator}

\pnum
The comparison operators \lain{==}, \lain{!=}, \lain{<}, \lain{>},
\lain{<=}, \lain{>=} compare two values and yield \lain{bool}.

\begin{constraint}
  Operands of a comparison operator shall have compatible types
  (\S\,\ref{sec:types.compat}).
\end{constraint}

\begin{constraint}
  \lain{==} and \lain{!=} shall not be applied to struct types, array types,
  or enum types (\S\,\ref{sec:types.enum}).
\end{constraint}

\section{Logical operators}
\label{sec:expr.logical}
\index{logical operator}

\pnum
\lain{and} and \lain{or} are short-circuit operators: the right operand
is not evaluated if the result is determined by the left operand alone.
\lain{not} and \lain{!} are prefix unary negation.  All yield \lain{bool}.

\begin{constraint}
  Operands of \lain{and}, \lain{or}, \lain{not}, and \lain{!} shall have
  type \lain{bool}.
\end{constraint}

\section{Bitwise operators}
\label{sec:expr.bitwise}
\index{bitwise operator}

\pnum
The bitwise operators \lain{\&}, \lain{|}, \texttt{\textasciicircum}, \lain{<<}, \lain{>>}
operate on integer operands.  The result type is the wider operand type for
\lain{\&}, \lain{|}, \texttt{\textasciicircum}; for shift operators, the result type is
the left operand type.

\begin{constraint}
  Operands of bitwise operators shall be integer types.
\end{constraint}

\begin{constraint}
  The shift amount (right operand of \lain{<<} and \lain{>>}) shall
  satisfy $0 \leq \mathit{shift} < N$ where $N$ is the bit width of the
  left operand.  A shift amount outside this range is
  \illformed{}~(\diagcode{E100}).
\end{constraint}

\section{The \texttt{in} operator}
\label{sec:expr.in}
\index{in operator}
\index{in guard}

\pnum
The expression \lain{idx in arr} evaluates to \lain{true} if and only if
$0 \leq \mathit{idx} < \mathit{arr.len}$.  The operand \textit{arr} shall
be an array or slice type; \textit{idx} shall be an integer type.

\pnum
When \lain{idx in arr} appears as the condition of a \lain{while} or
\lain{if} statement (or as the leftmost operand of an \lain{and}-chain
condition), it establishes an \textbf{in guard}: inside the body of that
statement, the implementation shall accept \lain{arr[idx]} without
requiring further bounds verification.

\pnum
In-guards propagate through \lain{and}-chains: if the left operand of
\lain{and} establishes an in-guard, that guard is active when evaluating
the right operand.

\begin{example}
\begin{laincode}
// In-guard: arr[i] is safe without explicit VRA
while i in arr {
    var x = arr[i]   // OK: in-guarded
    i = i + 1
}

// And-chain propagation:
while l.pos in l.src and (l.src[l.pos] as int) != '"' {
    l.pos = l.pos + 1
}
\end{laincode}
\end{example}

\begin{constraint}
  An in-guard of the form \lain{idx in arr} covers only \lain{arr[idx]}
  exactly.  Accesses \lain{arr[idx + k]} for any $k \neq 0$ are not
  covered by this guard and require separate verification or an
  \lain{unsafe} block.
\end{constraint}

\section{Cast expression (\texttt{as})}
\label{sec:expr.cast}
\index{cast expression}
\index{as operator}

\pnum
The expression \lain{expr as T} explicitly converts the value of
\textit{expr} to type \lain{T}.  The following casts are permitted:

\begin{longtable}{@{}lll@{}}
\toprule
\textbf{From} & \textbf{To} & \textbf{Semantics} \\
\midrule
\endfirsthead
\toprule
\textbf{From} & \textbf{To} & \textbf{Semantics} \\
\midrule
\endhead
\bottomrule
\endlastfoot
integer & integer & Truncation (narrowing) or extension (widening) \\
integer & float   & Nearest representable value (IEEE~754) \\
float   & integer & Truncation toward zero \\
float   & float   & Precision change \\
pointer & pointer & Reinterpretation (\lain{unsafe} required) \\
integer & pointer & Reinterpretation (\lain{unsafe} required) \\
pointer & integer & Reinterpretation (\lain{unsafe} required) \\
\end{longtable}

\begin{constraint}
  Casts between pointer types, between a pointer and an integer, or between
  any type and a pointer shall occur only inside an \lain{unsafe} block.
  Such a cast outside \lain{unsafe} is \illformed{}~(\diagcode{E012}).
\end{constraint}

\begin{constraint}
  A cast between non-numeric, non-pointer types (e.g., struct to integer)
  is \illformed{}.
\end{constraint}

\section{Move expression (\texttt{mov})}
\label{sec:expr.mov}
\index{move expression}

\pnum
\lain{mov expr} transfers ownership of the value produced by \textit{expr}.
The source is placed in the Consumed state; any subsequent use of the
source variable is \illformed{}~(\diagcode{E001}).  See \S\,\ref{sec:11-ownership}.

\section{Borrow expressions (\texttt{ref}, \texttt{mut})}
\label{sec:expr.borrow}
\index{borrow expression}

\pnum
\lain{ref expr} creates a shared (read-only) borrow of \textit{expr}.
\lain{mut expr} creates a mutable borrow of \textit{expr}.
Both expressions are evaluated in argument contexts to pass values without
transferring ownership.  See \S\,\ref{sec:11-ownership}.

\section{Field access}
\label{sec:expr.field}
\index{field access}

\pnum
The expression \lain{expr.field} accesses the named field of a struct or
ADT value, or a property of an array or slice (\lain{.len}, \lain{.data}).

\begin{constraint}
  The field name shall exist on the type of \textit{expr}.
\end{constraint}

\begin{constraint}
  Writing to a field (\lain{expr.field = v}) requires \textit{expr} to be
  a mutable lvalue.
\end{constraint}

\section{Index expression}
\label{sec:expr.index}
\index{index expression}

\pnum
The expression \lain{arr[idx]} accesses the element of array or slice
\textit{arr} at integer index \textit{idx}.

\begin{constraint}
  The value range analysis shall prove $0 \leq \mathit{idx} < \mathit{arr.len}$
  at every index expression.  An index expression for which this cannot be
  proved is \illformed{}~(\diagcode{E014}).  An \lain{in} guard
  (\S\,\ref{sec:expr.in}) or an \lain{unsafe} block suppresses this
  requirement.
\end{constraint}

\section{Function call}
\label{sec:expr.call}
\index{function call}

\pnum
A function call expression \lain{f(args)} evaluates the arguments and
transfers control to the callee.  Arguments are evaluated left-to-right.

\pnum
Argument ownership modes shall be given explicitly at the call site:

\begin{longtable}{@{}ll@{}}
\toprule
\textbf{Syntax} & \textbf{Meaning} \\
\midrule
\endfirsthead
\toprule
\textbf{Syntax} & \textbf{Meaning} \\
\midrule
\endhead
\bottomrule
\endlastfoot
\lain{f(x)}     & shared borrow (default; parameter mode must be \lain{shared}) \\
\lain{f(var x)} & mutable borrow (parameter mode must be \lain{var}) \\
\lain{f(mov x)} & ownership transfer (parameter mode must be owned/\lain{mov}) \\
\end{longtable}

\begin{constraint}
  The mode annotation at the call site shall match the declared parameter
  mode.  A mismatch is \illformed{}.
\end{constraint}

\begin{constraint}
  The number of arguments shall equal the number of parameters (unless the
  callee is a variadic C function declared with \lain{...}).
\end{constraint}

\pnum
\textbf{Two-phase borrows}: during argument evaluation, borrows of the
receiver of a method call are placed in the RESERVED phase and promoted to
ACTIVE after all arguments are evaluated
(see \S\,\ref{sec:11-ownership}).

\section{Block expression}
\label{sec:expr.block}
\index{block expression}

\pnum
A block expression \texttt{\{ stmts... expr \}} evaluates each statement
in order, then evaluates the final expression, whose value is the value of
the block.  If the final expression is absent, the block has no value.

\section{If expression}
\label{sec:expr.if}
\index{if expression}

\pnum
An \lain{if} expression of the form
\lain{if cond \{ e1 \} else \{ e2 \}} evaluates \textit{cond}; if true,
evaluates $e_1$; otherwise, evaluates $e_2$.

\begin{constraint}
  The types of $e_1$ and $e_2$ shall be compatible.  If they are not, the
  expression is \illformed{}.
\end{constraint}

\begin{constraint}
  \textit{cond} shall have type \lain{bool}.
\end{constraint}

\section{Match expression (\texttt{case})}
\label{sec:expr.match}
\index{case expression}

\pnum
A \lain{case} expression (\gramref{match-expression}) selects one of
several alternatives based on pattern matching.  All arms shall produce
values of compatible types; the type of the \lain{case} expression is that
common type.  See \S\,\ref{sec:15-match} for full exhaustiveness rules.

\section{Constructor expression}
\label{sec:expr.constructor}
\index{constructor expression}

\pnum
A constructor expression builds a value of a struct or ADT type by
supplying all field values.  See \S\,\ref{sec:types.struct} and
\S\,\ref{sec:types.adt} for constructor rules.

\section{Compound assignment}
\label{sec:expr.compound-assign}
\index{compound assignment}

\pnum
Compound assignment \lain{x op= e} is syntactic sugar for
\lain{x = x op e}.  The left-hand side shall be a mutable lvalue.
All compound forms of the arithmetic and bitwise operators are supported
(see \S\,\ref{sec:lex.operators}).

% 09-statements.tex — Clause 9: Statements
\chapter{Statements}
\label{sec:09-statements}
\index{statement}

\pnum
A statement is a unit of execution.  Statements are terminated by line
feeds (\S\,\ref{sec:lex.whitespace}).  The grammar production for
\gramref{statement} is in Annex~A.

\section{Block statement}
\label{sec:stmt.block}
\index{block statement}

\pnum
A block statement \lain{\{ stmts... \}} creates a new scope
(\S\,\ref{sec:basics.scopes}) and executes each contained statement in
order.  Names declared inside a block are not accessible outside it.

\section{Immutable declaration}
\label{sec:stmt.decl}
\index{immutable declaration}

\pnum
\lain{name [Type] = expr} declares and initializes an immutable binding.
The type may be omitted (inferred from the initializer).  The binding
cannot be reassigned after initialization.

\begin{laincode}
x i32 = 42          // explicit type
y = 10              // type inferred (i32)
p Point = Point(1, 2)
\end{laincode}

\begin{constraint}
  An immutable declaration requires an initializer expression.
  The initializer shall not be \lain{undefined} (immutable bindings must
  have a definite value).
\end{constraint}

\begin{constraint}
  When a type annotation is present, the type of the initializer shall be
  compatible with the declared type (\S\,\ref{sec:types.compat}).
\end{constraint}

\begin{note}
  When the type is omitted, the distinction between declaration and
  reassignment is resolved by scope: \lain{x = 5} where \lain{x} is not
  yet declared creates a new immutable binding; where \lain{x} is already
  declared as \lain{var}, it is a reassignment.
\end{note}

\section{Mutable variable declaration}
\label{sec:stmt.var}
\index{variable declaration}

\pnum
\lain{var name Type = expr} declares a mutable variable.  Mutable variables
can be reassigned.  The type may be omitted if it can be inferred from the
initializer.

\begin{laincode}
var counter i32 = 0
counter = counter + 1       // OK: var is mutable

var buf u8[256]              // OK: type present, defaults to undefined
\end{laincode}

\begin{constraint}
  A mutable declaration shall include either an explicit type annotation
  or an initializer expression.  A declaration with neither is
  \illformed{}~(\diagcode{E100}).
\end{constraint}

\begin{constraint}
  When both a type annotation and an initializer are present, the type of
  the initializer shall be compatible with the declared type
  (\S\,\ref{sec:types.compat}).
\end{constraint}

\section{Assignment statement}
\label{sec:stmt.assign}
\index{assignment statement}

\pnum
An assignment statement \lain{lvalue = expr} stores the value of
\textit{expr} into the storage designated by \textit{lvalue}.
An assignment target without a type annotation is always a reassignment,
never a declaration.

\begin{constraint}
  The left-hand side shall be a mutable lvalue (\S\,\ref{sec:basics.lvalues}).
  Assigning to an immutable binding is \illformed{}~(\diagcode{E009}).
\end{constraint}

\begin{constraint}
  Overwriting a linear variable whose current value has not been consumed
  would leak that value.  Such an assignment is
  \illformed{}~(\diagcode{E003}).
\end{constraint}

\begin{note}
  Assigning to a linear variable after its previous value has been consumed
  (e.g., loop reassign-and-consume pattern) is permitted.  The implementation
  resets the variable's linear state to Unconsumed after a legal overwrite.
\end{note}

\begin{example}
\begin{laincode}
import std.fs

proc main() int {
    var f = open_file("data.txt", "r")
    close_file(mov f)               // f is now Consumed
    f = open_file("data.txt", "r")  // OK: overwriting Consumed state
    close_file(mov f)
    return 0
}
\end{laincode}
\end{example}

\section{Expression statement}
\label{sec:stmt.expr}
\index{expression statement}

\pnum
An expression followed by a statement terminator is an expression
statement.  The value of the expression is discarded.

\section{If statement}
\label{sec:stmt.if}
\index{if statement}

\pnum
\lain{if cond \{ body \} [else if cond \{ body \}]... [else \{ body \}]}
evaluates \textit{cond}; if true, executes the then-body; otherwise,
tests the \lain{else if} clauses in order; if none matches, executes the
\lain{else} body (if present).

\begin{constraint}
  \textit{cond} shall have type \lain{bool}.
\end{constraint}

\pnum
\textbf{Range refinement}: inside the then-branch of \lain{if x < N \{ \}},
the value range analysis (\S\,\ref{sec:13-vra}) knows that
$x \in (-\infty, N-1]$; inside the else-branch, $x \in [N, +\infty)$.

\pnum
\textbf{In-guard}: when \textit{cond} contains or is an \lain{in}
expression (\S\,\ref{sec:expr.in}), the in-guard is active inside the
then-branch only.

\pnum
\textbf{Branch consistency}: if one branch consumes a linear variable,
all branches shall consume it.  If branches disagree on the linear state
of a variable from the enclosing scope, the program is
\illformed{}~(\diagcode{E003}).

\section{While statement}
\label{sec:stmt.while}
\index{while statement}

\pnum
\lain{while cond \{ body \}} repeatedly evaluates \textit{cond} and, if
true, executes \textit{body}.

\begin{constraint}
  An unbounded \lain{while} loop (one without a \lain{decreasing} measure)
  shall not appear inside a \lain{func}.  Such use is
  \illformed{}~(\diagcode{E011}).
\end{constraint}

\begin{constraint}
  \textit{cond} shall have type \lain{bool}.
\end{constraint}

\pnum
\textbf{In-guard}: when \textit{cond} contains or is an \lain{in}
expression, the in-guard is active inside the loop body for each
iteration.

\subsection{Bounded while statement}
\label{sec:stmt.while.bounded}
\index{bounded while loop}
\index{decreasing measure}

\pnum
\lain{while cond decreasing M \{ body \}} is a bounded loop.  In addition
to the while semantics above, the implementation shall statically verify:

\begin{enumerate}
  \item \textbf{Non-negativity} (\diagcode{E080}): the implementation shall
        prove that $M \geq 0$ whenever \textit{cond} is true.
  \item \textbf{Variable dependency} (\diagcode{E081}): $M$ shall
        reference at least one variable that is assigned in \textit{body}.
        An $M$ that references no such variable would never decrease.
  \item \textbf{Strict decrease} (\diagcode{E082}): the implementation
        shall prove that every execution of \textit{body} strictly
        decreases the value of $M$.
\end{enumerate}

\pnum
A bounded \lain{while} loop is permitted inside a \lain{func} because the
implementation has statically verified termination.

\pnum
The following patterns are accepted by the standard verification algorithm:

\begin{longtable}{@{}lll@{}}
\toprule
\textbf{Condition} & \textbf{Measure} & \textbf{Why} \\
\midrule
\endfirsthead
\toprule
\textbf{Condition} & \textbf{Measure} & \textbf{Why} \\
\midrule
\endhead
\bottomrule
\endlastfoot
\lain{i < n}    & \lain{n - i}      & $i$ increases $\Rightarrow$ measure decreases \\
\lain{n > 0}    & \lain{n}          & $n$ decreases directly \\
\lain{a < b}    & \lain{b - a}      & either $a$ increases or $b$ decreases \\
\lain{i in arr} & \lain{arr.len - i}& \lain{in} implies $i < \texttt{arr.len}$ \\
\end{longtable}

\begin{example}
\begin{laincode}
func count(n int) int {
    var i = 0
    while i < n decreasing n - i {
        i = i + 1
    }
    return i
}
\end{laincode}
\end{example}

\begin{rationale}
  The \lain{decreasing} keyword bridges the gap between pure functions
  (which must terminate) and the common pattern of finite-input state
  machines (lexers, parsers) that are provably terminating.  The programmer
  states the termination argument; the compiler verifies it.
\end{rationale}

\section{Return statement}
\label{sec:stmt.return}
\index{return statement}

\pnum
\lain{return [expr]} terminates execution of the current function or
procedure and transfers the value of \textit{expr} (if present) to the
caller.

\begin{constraint}
  At the point of a \lain{return}, all linear variables in scope shall
  have been consumed.  A live (Unconsumed) linear variable at \lain{return}
  is \illformed{}~(\diagcode{E003}).
\end{constraint}

\begin{constraint}
  \textbf{Q-009.}
  A return statement shall not be prefixed with \lain{mov} or \lain{var}.
  The forms \lain{return mov x} and \lain{return var x} are
  syntactically \illformed{}~(\diagcode{E100}).  Ownership transfer in
  return position is implicit from the function's declared return type.
\end{constraint}

\section{Defer statement}
\label{sec:stmt.defer}
\index{defer statement}

\pnum
\lain{defer \{ body \}} schedules \textit{body} for execution immediately
before the enclosing scope exits, on all exit paths (\lain{return},
\lain{break}, \lain{continue}, or fall-through).  Multiple \lain{defer}
statements execute in LIFO order.

\pnum
\textbf{Interaction with linear types}: if a linear variable is consumed
inside a \lain{defer} body, the variable is marked as defer-consumed.  It
is treated as unconsumed for the purposes of the main body (it may be used
before the scope exit) and as consumed at scope exit.

\begin{constraint}
  A \lain{defer} body shall not contain \lain{return}, \lain{break}, or
  \lain{continue} statements that would transfer control out of the deferred
  body.
\end{constraint}

\begin{example}
\begin{laincode}
import std.fs

proc process() int {
    var f = open_file("data.txt", "r")
    defer { close_file(mov f) }   // f consumed at scope exit
    // ... use f ...
    return 0
}
\end{laincode}
\end{example}

\section{Unsafe block}
\label{sec:stmt.unsafe}
\index{unsafe block}

\pnum
\lain{unsafe \{ body \}} is an unsafe block (see \S\,\ref{sec:18-unsafe}).
It creates a scope in which certain static checks are relaxed.

\section{Comptime if statement}
\label{sec:stmt.comptime-if}
\index{comptime if}

\pnum
\lain{comptime if cond \{ then \} [else \{ else \}]} evaluates
\textit{cond} at compile time.  Only the taken branch is type-checked and
emitted; the other branch is discarded.  \textit{cond} shall be a
compile-time constant expression (\S\,\ref{sec:14-generics}).

\begin{note}
  \lain{comptime if} is the Lain equivalent of C preprocessor
  \texttt{\#if}: it enables platform-specific or type-specific code
  selection without dead-code warnings.
\end{note}

\section{For statement}
\label{sec:stmt.for}
\index{for statement}

\pnum
\lain{for i in start..end \{ body \}} iterates with \textit{i} ranging
from \textit{start} (inclusive) to \textit{end} (exclusive).  The loop
variable \textit{i} is immutable; its value range is statically known to
be $[\textit{start}, \textit{end}-1]$ throughout the body, enabling safe
array indexing without runtime bounds checks.

\pnum
\lain{for} loops are permitted in both \lain{func} and \lain{proc}.

\begin{constraint}
  \textit{start} and \textit{end} shall have integer types.
\end{constraint}

% 10-declarations.tex — Clause 10: Declarations
\chapter{Declarations}
\label{sec:10-declarations}
\index{declaration}

\pnum
This clause specifies top-level declarations.  Local variable declarations
are specified in \S\,\ref{sec:stmt.var}.  The grammar productions are in
Annex~A (\gramref{function-declaration}, \gramref{struct-declaration},
\gramref{enum-declaration}).

\section{Declaration attributes}
\label{sec:decl.attributes}
\index{attribute}

\pnum
A declaration may be preceded by zero or more attributes.  An attribute
has the form \tok{[name]} or \tok{[name(args)]}, where \textit{name} is
a recognized attribute identifier from the closed set in
Table~\ref{tab:attributes}, and \textit{args} (if present) is a
comma-separated list of expressions.  Multiple attributes on a single
declaration are written one per line.

\begin{laincode}
[fast_math]
func dot(a f64[], b f64[], n int) f64 { ... }

[private]
func internal_helper() int { ... }

[fast_math]
[private]
func hot_path() f64 { ... }
\end{laincode}

\pnum
Recognized attribute names (closed set, version 1.1):

\begin{longtable}{@{}lll@{}}
\caption{Recognized attribute names}\label{tab:attributes}\\
\toprule
\textbf{Attribute} & \textbf{Applies to} & \textbf{Purpose} \\
\midrule
\endfirsthead
\toprule
\textbf{Attribute} & \textbf{Applies to} & \textbf{Purpose} \\
\midrule
\endhead
\bottomrule
\endlastfoot
\tok{[fast\_math]} & any function & Enable FMA fusion (P4 opt-out, local). \\
\tok{[private]}    & any declaration         & Module-private visibility. \\
\end{longtable}

\begin{constraint}
  \textbf{Q-017.}
  An attribute name not listed in Table~\ref{tab:attributes} is
  \illformed{}~(\diagcode{E103}).  An attribute that omits the
  identifier (e.g.\ \texttt{[ ]} or \texttt{[(arg)]}) is
  \illformed{}~(\diagcode{E102}).
\end{constraint}

\begin{note}
  Attributes are recognized at declaration position only.  At expression
  position, \tok{[} and \tok{]} continue to delimit array and slice
  subscripts and array literals; the parser disambiguates by context.
\end{note}

\section{Struct declarations}
\label{sec:decl.struct}
\index{struct declaration}

\pnum
A struct declaration introduces a named product type with one or more
named fields.  Fields are listed in declaration order; order is significant
for construction and layout.

\begin{laincode}
type Point {
    x i32
    y i32
}
\end{laincode}

\pnum
A field may be annotated with \lain{mov}, making the field (and therefore
the containing struct) a linear type:

\begin{laincode}
type File {
    mov handle *FILE
}
\end{laincode}

\begin{note}
  A struct with zero fields (a \textit{unit} struct, \lain{type Empty \{\}})
  is permitted: it is a zero-sized marker type, constructed as \lain{Empty()}.
\end{note}

\section{Enum and ADT declarations}
\label{sec:decl.enum}
\index{enum declaration}
\index{ADT declaration}

\pnum
An enum declaration introduces a type whose values are drawn from a fixed,
named set of variants.  If every variant has no associated fields, the type
is a simple enum (\S\,\ref{sec:types.enum}).  If at least one variant has
associated fields, the type is an algebraic data type (ADT)
(\S\,\ref{sec:types.adt}).

\begin{laincode}
type Direction {              // simple enum (one variant per line)
    North
    South
    East
    West
}

type Shape {                                  // ADT
    Circle     { radius i32 }
    Rectangle  { width i32, height i32 }
    Point
}
\end{laincode}

\section{Function declarations}
\label{sec:decl.func}
\index{function declaration}

\pnum
A function declaration introduces a \lain{func}: a pure, terminating,
side-effect-free callable.  Its grammar production is
\gramref{function-declaration}.

\pnum
The signature consists of: the keyword \lain{func}, a name, an optional
list of type parameters (\S\,\ref{sec:14-generics}), a parameter list, and a
return type annotation.

\subsection{Parameters}
\label{sec:decl.func.params}
\index{parameter}

\pnum
Each parameter has a name, a type, and an ownership mode.  The ownership
mode determines how the argument is passed:

\begin{longtable}{@{}lll@{}}
\toprule
\textbf{Syntax} & \textbf{Mode} & \textbf{ABI (struct T)} \\
\midrule
\endfirsthead
\toprule
\textbf{Syntax} & \textbf{Mode} & \textbf{ABI (struct T)} \\
\midrule
\endhead
\bottomrule
\endlastfoot
\lain{name T}        & shared (read-only) & \texttt{const T*} \\
\lain{var name T}    & mutable            & \texttt{T*}       \\
\lain{mov name T}    & owned (transfer)   & \texttt{T} (by value) \\
\end{longtable}

\begin{note}
  For primitive types (integer, float, bool), the ABI passes the value
  directly regardless of ownership mode, as for primitive C arguments.
\end{note}

\pnum
A parameter whose declared type is the built-in \lain{type} --- written
\lain{name type}, with no \lain{comptime} prefix --- is a compile-time type
parameter; it is bound to a concrete type at every call site, by inference
from a value argument or explicitly (\S\,\ref{sec:14-generics}).

\subsection{Return type annotation}
\label{sec:decl.func.return}
\index{return type}

\pnum
The return type is written after the parameter list.  It may carry a value
range constraint:

\begin{laincode}
func abs(x int) int >= 0 { ... }
\end{laincode}

\begin{constraint}
  All return paths of a \lain{func} shall produce a value whose type and
  value range satisfy the declared return type and constraint.
\end{constraint}

\subsection{Purity constraints}
\label{sec:decl.func.purity}
\index{func purity constraints}

\begin{constraint}
  A function's \lain{effects} row is an upper bound on what its body does, and its
  default is $\emptyset$.  An effect the body has and the row does not name is
  \illformed{}~(\diagcode{E011} where the bound is $\emptyset$,
  \diagcode{E130} where a non-empty row understates it).  The three
  constraints below are consequences of that one rule, not separate rules.
\end{constraint}

\begin{constraint}
  A function shall not contain an unbounded \lain{while} loop --- one whose
  measure is neither given by \lain{decreasing} nor inferable --- unless its row
  names \lain{diverge}~(\diagcode{E011}/\diagcode{E082}).
\end{constraint}

\begin{constraint}
  A function shall not call one whose row names an effect its own row does
  not~(\diagcode{E011}).  The callee's \emph{introducer} is irrelevant: what is
  refused is the effect, not the keyword that declared it.
\end{constraint}

\begin{constraint}
  A function shall not participate in mutual recursion unless its row names
  \lain{diverge}~(\diagcode{E011}); termination cannot be established for a cycle
  no measure covers.  Direct self-recursion \emph{with} an inferred or given
  measure is permitted.
\end{constraint}

\section{The entry point}
\label{sec:decl.proc}
\index{entry point}

\pnum
\textbf{\lain{proc} has been removed.}  It was an introducer that granted effects
instead of stating them, and every effect it was needed for is now named by an
\lain{effects} row on an ordinary \lain{func} --- including \lain{diverge}, which
admits an unbounded loop.  The keyword remains \emph{reserved}, and using it is
\diagcode{E100} with the replacement spelled out, so old source fails with an
explanation rather than with "expected declaration".  The same applies to the
function-pointer type \lain{*proc(...)}, replaced by \lain{*func(...)} with a row
(\S\,\ref{sec:fn.distinction}).

\pnum
The entry point of every Lain program shall be declared as:

\begin{laincode}
func main() int { ... }          // with `effects io` if it performs any
\end{laincode}

\begin{constraint}
  A complete \emph{program} shall define exactly one \lain{main} returning \lain{int} as
  its entry point.  An individual module (translation unit) compiled on its own
  --- e.g.\ a library --- need not define \lain{main}: the reference compiler
  accepts a \lain{main}-less source file, and a program missing an entry point
  is diagnosed when the translation outputs are linked.
\end{constraint}

\section{Extern declarations}
\label{sec:decl.extern}
\index{extern declaration}

\pnum
An extern declaration introduces a C function into the Lain name space
without providing a body.  See \S\,\ref{sec:17-interop} for full semantics.

\begin{laincode}
extern func printf(fmt *u8, ...) int effects io
extern func strlen(s *u8) usize effects
\end{laincode}

\begin{constraint}
  Variadic parameters (\lain{...}) shall only appear in \lain{extern}
  declarations.  User-defined variadic functions are \illformed{}.
\end{constraint}

\section{Comptime declarations}
\label{sec:decl.comptime}
\index{comptime declaration}

\begin{note}
  Top-level \lain{comptime \{ ... \}} blocks --- evaluated by the compiler
  before any runtime code is emitted, to define compile-time constants or
  trigger \lain{compileError} assertions --- are a \emph{planned} feature not
  yet accepted by the reference implementation (a top-level \lain{comptime}
  block is presently \diagcode{E100}).  The \lain{comptime} keyword is
  nonetheless reserved, and \lain{comptime if} (\S\,\ref{sec:stmt.comptime-if})
  is available.
\end{note}

% 11-ownership.tex — Clause 11: Ownership and linearity
\chapter{Ownership and linearity}
\label{sec:11-ownership}
\index{ownership}
\index{linearity}

\pnum
This clause specifies Lain's linear type system, ownership modes, borrow
checker, and non-lexical lifetime rules.  All rules in this clause are
enforced at compile time with zero runtime overhead.

\section{Ownership modes}
\label{sec:own.modes}
\index{ownership mode}

\pnum
Every parameter, return value, and struct field has one of three ownership
modes:

\begin{longtable}{@{}llll@{}}
\toprule
\textbf{Mode} & \textbf{Param syntax} & \textbf{Call-site} & \textbf{Semantics} \\
\midrule
\endfirsthead
\toprule
\textbf{Mode} & \textbf{Param syntax} & \textbf{Call-site} & \textbf{Semantics} \\
\midrule
\endhead
\bottomrule
\endlastfoot
Shared  & \lain{name T}     & \lain{f(x)}     & read-only; multiple concurrent borrows allowed \\
Mutable & \lain{var name T} & \lain{f(var x)} & exclusive read-write; no other borrow concurrent \\
Owned   & \lain{mov name T} & \lain{f(mov x)} & transfers ownership; source becomes Consumed \\
\end{longtable}

\pnum
Local variables declared with \lain{var name = expr} have mutable storage
accessible within their scope; their ownership mode in the type system is
Owned (they hold linear values that must be consumed at scope exit).

\section{Linear states}
\label{sec:own.states}
\index{linear state}

\pnum
Every variable of a linear type has a \textit{linear state} at each program
point.  The three states are:

\begin{description}[leftmargin=2em]
  \item[Uninit]
    The variable has been declared but no value has been assigned.
  \item[Unconsumed]
    The variable holds a value that has not yet been transferred.
  \item[Consumed]
    The variable's value has been transferred via \lain{mov} or a consuming
    function call.
\end{description}

\pnum
The state machine governing transitions is:

\begin{center}
\begin{tabular}{lll}
\toprule
\textbf{State before} & \textbf{Operation} & \textbf{State after} \\
\midrule
Uninit      & first assignment              & Unconsumed \\
Unconsumed  & \lain{mov x} (transfer)       & Consumed \\
Unconsumed  & borrow (\lain{var} or shared) & Unconsumed (borrow active) \\
Consumed    & assignment (overwrite)        & Unconsumed \\
Consumed    & read or \lain{mov}            & \illformed{} \\
Uninit      & read or \lain{mov}            & \illformed{} \\
\bottomrule
\end{tabular}
\end{center}

\begin{note}
  The Consumed~$\to$~Unconsumed transition via overwrite enables the
  loop-reassign-then-consume pattern: a linear variable declared outside a
  loop may be consumed and reassigned on each iteration without leaking.
\end{note}

\section{Diagnostics}
\label{sec:own.diag}
\index{ownership diagnostics}

\begin{constraint}
  Using a variable that is in the Consumed or Uninit state (reading its
  value or transferring it with \lain{mov}) is \illformed{}~(\diagcode{E001}).
\end{constraint}

\begin{constraint}
  Transferring ownership of a variable that has already been transferred
  (double-move) is \illformed{}~(\diagcode{E002}).
\end{constraint}

\begin{constraint}
  A linear variable that is in the Unconsumed state when its scope exits
  (resource leak) is \illformed{}~(\diagcode{E003}).
\end{constraint}

\begin{constraint}
  An aliasing violation — attempting a mutable borrow when a shared borrow
  is active, or a shared borrow when a mutable borrow is active, or a
  second mutable borrow of the same variable — is
  \illformed{}~(\diagcode{E004}).
\end{constraint}

\begin{constraint}
  Attempting to move a variable that has an active borrow is
  \illformed{}~(\diagcode{E005}).
\end{constraint}

\begin{constraint}
  Moving a variable inside a loop body where the loop may execute more than
  once is \illformed{}~(\diagcode{E006}).
\end{constraint}

\begin{constraint}
  If a linear variable is consumed in one branch of an \lain{if}/\lain{else}
  chain but not in all branches, the program is
  \illformed{}~(\diagcode{E007}).
\end{constraint}

\begin{constraint}
  Moving a value that is currently borrowed (\lain{var} or shared), or
  moving a struct after some of its fields have already been consumed,
  is \illformed{}~(\diagcode{E008}).
\end{constraint}

\begin{constraint}
  If a linear value is in state \textsc{Consumed} on some control-flow
  paths through an \lain{if}/\lain{case} statement and in state
  \textsc{Unconsumed} on others, the merged state at the join point
  is undefined and the program is \illformed{}~(\diagcode{E016}). The
  same diagnostic applies at field level: if a linear field of a
  struct is consumed on some paths but not on others, \diagcode{E016}
  is issued for that field.
\end{constraint}

\section{Mutable parameters for primitive types (\texttt{var T})}
\label{sec:own.var-primitive}
\index{var parameter!primitive type}
\index{write-through}

\pnum
A parameter declared \lain{var name T}, where \lain{T} is a primitive type
(\lain{iN}, \lain{uN}, \lain{f32}, \lain{f64}, \lain{bool}), is passed
\textbf{by pointer} in the C ABI: the parameter's C type is \texttt{T*}.

\pnum
\textbf{Write-through semantics}: an assignment to such a parameter inside
the function body writes through the pointer and is visible to the caller.
Reading the parameter's value automatically dereferences the pointer.

\begin{example}
\begin{laincode}
proc increment(var n i32) {
    n = n + 1    // writes through: caller's variable is updated
}

proc main() i32 {
    var x = 0
    increment(var x)
    return x     // x == 1
}
\end{laincode}
\end{example}

\pnum
\textbf{Forwarding}: when a \lain{var T} primitive parameter is passed as
argument to another \lain{var T} parameter, the pointer is forwarded
directly without dereferencing and re-taking the address.

\pnum
This behavior is identical to mutable parameters of struct type, where the
by-pointer ABI was already in effect.  The distinction is that for primitive
types the codegen previously did not dereference reads; this was a bug that
has been corrected.

\begin{note}
  The call site shall annotate the argument with \lain{var}:
  \lain{f(var x)}.  Omitting \lain{var} passes a shared (read-only)
  borrow, which does not permit writes through the callee.
\end{note}

\section{The read-write lock invariant}
\label{sec:own.rwlock}
\index{read-write lock invariant}

\pnum
\textbf{Invariant}: At every program point, for every variable~$V$, exactly
one of the following holds:
\begin{enumerate}
  \item $V$ has $N \geq 0$ active shared borrows and zero active mutable borrows.
  \item $V$ has exactly one active mutable borrow and zero shared borrows.
\end{enumerate}

\pnum
This invariant is the compile-time analogue of a reader-writer lock.  A
conforming implementation shall enforce it at every program point.

\section{Non-lexical lifetimes (NLL)}
\label{sec:own.nll}
\index{non-lexical lifetime}

\pnum
A borrow's lifetime ends at the borrow's \textit{last use}, not at the end
of the variable's lexical scope.  The implementation computes last-use
points via a use-analysis pre-pass over the AST.

\begin{example}
\begin{laincode}
var data = Data(42)
var r = get_ref(var data)  // mutable borrow starts
use_ref(var r)             // last use of r: borrow ends HERE (NLL)
var y = read(data)         // OK: borrow has expired
\end{laincode}
\end{example}

\section{Two-phase borrows}
\label{sec:own.two-phase}
\index{two-phase borrow}

\pnum
During argument evaluation for a function call, borrows of the receiver
object are placed in the \textit{RESERVED} phase.  A RESERVED borrow
allows concurrent shared reads of the same owner but blocks mutable writes
and moves.  After all arguments have been evaluated, RESERVED borrows are
promoted to the \textit{ACTIVE} phase.

\pnum
This mechanism enables UFCS calls of the form \lain{x.f(x.field)} where
the first argument (the receiver) requires a mutable borrow of \lain{x}
while the second argument reads a field of \lain{x}.

\begin{longtable}{@{}lll@{}}
\toprule
\textbf{Phase} & \textbf{Allows on owner} & \textbf{Blocks on owner} \\
\midrule
\endfirsthead
\toprule
\textbf{Phase} & \textbf{Allows on owner} & \textbf{Blocks on owner} \\
\midrule
\endhead
\bottomrule
\endlastfoot
RESERVED & shared reads & mutable writes, moves \\
ACTIVE   & (nothing)    & all other access \\
\end{longtable}

\begin{example}
\begin{laincode}
proc push_n(var v Vec, n int) { v.data = v.data + n }

var v = Vec(0, 10)
v.push_n(v.cap)   // desugars to push_n(var v, v.cap)
                  // var v enters RESERVED phase; v.cap read is allowed
\end{laincode}
\end{example}

\section{Move semantics}
\label{sec:own.move}
\index{move semantics}

\pnum
The \lain{mov} operator transfers ownership.  The source variable
transitions to Consumed state; any subsequent read or move of the source
is \illformed{}~(\diagcode{E001}).

\pnum
Ownership transfer is the only way to pass a linear value to a consuming
parameter.  The \lain{mov} keyword shall be present at the call site:
\lain{f(mov x)}.

\section{Borrowed match}
\label{sec:own.borrowed-match}
\index{borrowed match}

\pnum
The \lain{case \&expr \{ ... \}} form (see \S\,\ref{sec:15-match}) performs
pattern matching without consuming the scrutinee.  A temporary shared
borrow of the scrutinee is registered for the duration of the match body
and released when the match completes.

\begin{constraint}
  Inside a \lain{case \&expr} body, the scrutinee shall not be mutated
  or moved.  Attempting to do so violates the active shared
  borrow~(\diagcode{E004}).
\end{constraint}

\section{Loop and ownership}
\label{sec:own.loop}
\index{ownership in loops}

\pnum
A linear variable declared outside a loop may be consumed inside the loop
body if and only if it is also reassigned (overwritten) on every path
through the loop body before the next iteration begins.

\begin{example}
\begin{laincode}
import std.fs

proc main() int {
    var f = open_file("data.txt", "r")
    close_file(mov f)        // consumed before loop
    var i = 0
    while i < 3 decreasing 3 - i {
        f = open_file("data.txt", "r")  // reassign (Consumed -> Unconsumed)
        close_file(mov f)               // consume
        i = i + 1
    }
    return 0
}
\end{laincode}
\end{example}

\section{Branch consistency}
\label{sec:own.branches}
\index{branch consistency}

\pnum
After an \lain{if}/\lain{else} chain or \lain{case} match, the linear
state of each variable from the enclosing scope is the
\textit{intersection} of its states across all branches:

\begin{itemize}
  \item If consumed in all branches: the variable is Consumed after the chain.
  \item If consumed in no branches: the variable is Unconsumed after the chain.
  \item If consumed in some but not all branches: \illformed{}~(\diagcode{E016}).
\end{itemize}

\pnum
The same rule applies at field level: if a linear field of a struct is
consumed on some branches but not on others, \diagcode{E016} is emitted
for that field.

\begin{example}
\begin{laincode}
type Handle { mov ptr *i32 }

func maybe_consume(mov h Handle, flag bool) {
    if flag {
        drop(mov h)   // consumed in then-branch
    }
    // [E016]: h consumed in then but not in else (implicit else)
}
\end{laincode}
\end{example}

\begin{note}
  \diagcode{E007} (generic ownership mismatch) was the diagnostic emitted
  before \diagcode{E016} was introduced.  \diagcode{E016} is now the
  dedicated code for branch-inconsistency on linear values.
\end{note}

\section{Field-level tracking}
\label{sec:own.fields}
\index{field-level borrow tracking}

\pnum
The borrow checker tracks borrows at the field level.  Borrowing
non-overlapping fields of the same struct simultaneously is permitted:

\begin{laincode}
type Pair { a i32, b i32 }
var p = Pair(1, 2)
var ra = read_a(p.a)   // shared borrow of p.a (implicit)
var rb = read_b(p.b)   // shared borrow of p.b -- OK: non-overlapping
\end{laincode}

\begin{note}
  The reference implementation tracks field borrows at one level of depth.
  Nested field borrows (e.g., \lain{p.a.x} vs \lain{p.a.y}) are conservatively
  treated as overlapping.
\end{note}

% 12-functions.tex — Clause 12: Functions and procedures
\chapter{Functions and procedures}
\label{sec:12-functions}
\index{function}
\index{procedure}

\section{The func/proc distinction}
\label{sec:fn.distinction}
\index{func vs proc}

\pnum
Lain distinguishes between pure functions (\lain{func}) and side-effecting
procedures (\lain{proc}).  This distinction is enforced statically and
reflected in the type system.

\begin{longtable}{@{}lll@{}}
\toprule
\textbf{Property} & \textbf{\lain{func}} & \textbf{\lain{proc}} \\
\midrule
\endfirsthead
\toprule
\textbf{Property} & \textbf{\lain{func}} & \textbf{\lain{proc}} \\
\midrule
\endhead
\bottomrule
\endlastfoot
Side effects                      & forbidden          & allowed \\
Unbounded \lain{while}            & forbidden          & allowed \\
Bounded \lain{while decreasing}   & allowed            & allowed \\
\lain{for} loops                  & allowed            & allowed \\
Direct/indirect recursion         & forbidden          & allowed \\
Calling \lain{proc}               & forbidden          & allowed \\
Calling \lain{func}               & allowed            & allowed \\
Global mutable state              & no access          & full access \\
Termination guarantee             & yes                & no \\
\end{longtable}

\pnum
A \lain{func} is referentially transparent: calling it with the same
arguments always produces the same result.  Its evaluation has no
observable effect on any state outside its own stack frame.

\pnum
Local mutation through a \lain{var T} parameter (exclusive mutable
borrow) is permitted in a \lain{func}. The borrow checker guarantees
no aliasing, so the mutation is observable only through the
parameter binding itself; the function remains deterministic and
referentially transparent (same input borrow $\Rightarrow$ same
mutation pattern).

\begin{example}
\begin{laincode}
type Counter { val i32 }

func increment(var c Counter) {  // legal: mutation is local
    c.val = c.val + 1
}
\end{laincode}
\end{example}

\pnum
The compiler emits \diagcode{W130} when a \lain{proc} has no
observable side effect — no calls to other \lain{proc} or
\lain{extern proc}, no \lain{panic}, no mutable global access, no
unbounded \lain{while}, no self-recursion — to suggest that it could
be declared \lain{func}. This is a non-blocking warning; the
programmer may keep it as \lain{proc} when future evolution will
introduce effects.

\section{Purity enforcement}
\label{sec:fn.purity}
\index{purity enforcement}

\begin{constraint}
  A \lain{func} that calls a \lain{proc} or \lain{extern proc} is
  \illformed{}~(\diagcode{E011}).
\end{constraint}

\begin{constraint}
  A \lain{func} that contains an unbounded \lain{while} loop (one without
  a \lain{decreasing} measure) is \illformed{}~(\diagcode{E011}).
\end{constraint}

\begin{constraint}
  A \lain{func} that calls itself directly, or that participates in a
  mutual-recursion cycle, is \illformed{}~(\diagcode{E011}).  The
  implementation shall detect mutual recursion by performing a depth-first
  search on the call graph.
\end{constraint}

\begin{note}
  Generic functions (\lain{func} with \lain{type} parameters) are
  excluded from mutual-recursion detection until they are monomorphized.
\end{note}

\section{Termination analysis}
\label{sec:fn.termination}
\index{termination analysis}

\pnum
A \lain{func} is guaranteed to terminate because:
\begin{enumerate}
  \item all \lain{while} loops carry a verified termination measure;
  \item \lain{for} loops iterate over finite ranges;
  \item recursion is forbidden.
\end{enumerate}

\pnum
No termination analysis is performed for \lain{proc}s.

\section{Universal Function Call Syntax (UFCS)}
\label{sec:fn.ufcs}
\index{UFCS}

\pnum
An expression of the form \lain{x.f(args)} is desugared to \lain{f(x, args)}
before name resolution, provided \lain{f} is not a declared field of
\lain{x}'s type.  This is called Universal Function Call Syntax (UFCS).

\pnum
\textbf{UFCS ownership}: the ownership mode of the first argument in the
desugared call is determined by the declared mode of the first parameter of
\lain{f}:

\begin{itemize}
  \item If the first parameter is \lain{var}, the receiver is automatically
        mutably borrowed: \lain{x.f()} desugars to \lain{f(var x)}.
  \item If the first parameter is \lain{mov}, the receiver is automatically
        moved: \lain{x.f()} desugars to \lain{f(mov x)}.
  \item If the first parameter is shared (default), the receiver is shared:
        \lain{x.f()} desugars to \lain{f(x)}.
\end{itemize}

\pnum
\textbf{UFCS resolution algorithm}:
\begin{enumerate}
  \item Check if \lain{f} is a declared field of \lain{x}'s type.  If so,
        this is a field access, not UFCS.
  \item Search for a function named \lain{f} whose first parameter type is
        compatible with the type of \lain{x}.
  \item If found, desugar the call.
  \item If no matching function is found, the call is unresolved and the
        program is \illformed{} (an unresolved receiver or callee is reported;
        a bare undeclared identifier is \diagcode{E106},
        \S\,\ref{sec:basics.resolution}).
\end{enumerate}

\begin{example}
\begin{laincode}
func is_empty(s u8[]) bool { return s.len == 0 }

var buf = "hello"
var empty = buf.is_empty()   // desugars to: is_empty(buf)
\end{laincode}
\end{example}

\section{Calling conventions}
\label{sec:fn.abi}
\index{calling convention}

\pnum
The calling convention for Lain functions corresponds to the C99 calling
convention of the target platform, as determined by the C compiler used for
the translation output.  The parameter passing rules are:

\begin{longtable}{@{}lll@{}}
\toprule
\textbf{Lain parameter} & \textbf{C type} & \textbf{Semantics} \\
\midrule
\endfirsthead
\toprule
\textbf{Lain parameter} & \textbf{C type} & \textbf{Semantics} \\
\midrule
\endhead
\bottomrule
\endlastfoot
shared struct \lain{T} & \texttt{const T*} & pointer to read-only copy \\
\lain{var} struct \lain{T} & \texttt{T*}       & pointer, caller's copy \\
\lain{mov} struct \lain{T} & \texttt{T}        & by value (bitwise copy) \\
primitive any mode      & \texttt{T}        & by value \\
\end{longtable}

\section{Return values}
\label{sec:fn.return}
\index{return value}

\pnum
A return statement has the form \lain{return <expr>}.  The mode of the
returned value is determined by the function's declared return type and
follows the same rules as call-site argument passing.

\pnum
A \lain{return} may be prefixed with \lain{mov} or \lain{var}:
\lain{return mov x} makes the ownership transfer explicit, and
\lain{return var x} returns a mutable borrow (the pattern for a \lain{func}
whose declared return type is \lain{var T}).  When the return type is owned,
a bare \lain{return expr} already moves \lain{expr}, so \lain{mov} is optional
there.

\begin{note}
  A \lain{return var x} / \lain{return mov x} that would escape a borrow of a
  \emph{local} variable is a dangling reference (\diagcode{E010},
  \S\,\ref{sec:stmt.return}).  This supersedes an earlier design (Q-009) in
  which explicit \lain{mov}/\lain{var} in return position was rejected.
\end{note}


\section{The \texttt{panic} builtin}
\label{sec:fn.panic}
\index{panic}

\pnum
\lain{panic(msg)} is a builtin expression that terminates program
execution immediately.  Its type is \textit{Never} (the bottom type:
a subtype of every type).  Because it does not return, \lain{panic}
may appear in any expression position---including as the body of a
match arm whose declared type is otherwise unrelated.

\pnum
\lain{panic} may be called from both \lain{func} and \lain{proc}
contexts.  The call is deterministic (same input yields the same
termination), so it does not violate the purity of an enclosing
\lain{func}.

\begin{constraint}
  \textbf{Q-014, S16.}
  When \lain{panic} is executed, the program terminates immediately
  via the platform abort primitive.  No \lain{defer} block is
  executed.  No stack unwinding occurs.
\end{constraint}

\begin{note}
  The reference implementation emits, for \lain{panic(msg)}, the
  equivalent of
  \texttt{fprintf(stderr, "panic: \%.*s\textbackslash{}n", ...) ; abort()}.
  The message is written to the standard error stream before abort.
\end{note}

\begin{example}
\begin{laincode}
func head(s u8[]) u8 {
    if s.len == 0 {
        panic("empty slice")    // Never; satisfies the return type
    }
    return s[0]
}
\end{laincode}
\end{example}

\section{Function attributes}
\label{sec:fn.attributes}
\index{attribute}

\pnum
A function declaration may be prefixed with one or more attributes
(\S\,\ref{sec:decl.attributes}).  The attributes applicable to
functions in version~1.0 are:

\begin{itemize}
  \item \tok{[fast\_math]} --- enable fast floating-point optimizations
        (FMA fusion, contraction) for the body of this function.  This
        is an opt-out from the determinism guarantee (P4) and applies
        only locally; functions without \tok{[fast\_math]} remain
        bit-deterministic within their target platform.
  \item \tok{[private]} --- module-private visibility
        (\S\,\ref{sec:mod.visibility}).
\end{itemize}

\begin{note}
  The reference implementation translates \tok{[fast\_math]} on a
  function~$f$ to a pair of pragmas surrounding the body of $f$ in the
  emitted C source, enabling \texttt{fp-contract=fast} for $f$ only.
  The default (i.e., absence of \tok{[fast\_math]}) emits
  \texttt{-ffp-contract=off} flags, providing deterministic floating
  point within a single target platform and C compiler.
\end{note}

% 13-vra.tex — Clause 13: Value range constraints
\chapter{Value range constraints}
\label{sec:13-vra}
\index{value range analysis (VRA)}
\index{VRA}

\section{Overview}
\label{sec:vra.overview}

\pnum
Value Range Analysis (VRA) is a static analysis that tracks integer value
ranges through a program using interval arithmetic.  It is used to:
\begin{itemize}
  \item eliminate array bounds checks (\S\,\ref{sec:types.array});
  \item prove division-by-zero absence (\S\,\ref{sec:expr.arith});
  \item verify parameter and return constraints (\S\,\ref{sec:vra.constraints});
  \item prove in-guard safety (\S\,\ref{sec:expr.in}).
\end{itemize}

\pnum
VRA is decidable and polynomial in program size.  No SMT solver is required.
All analyses performed by a conforming implementation shall terminate.

\section{Range representation}
\label{sec:vra.range}
\index{range representation}

\pnum
A value range is an interval $[\ell,\, u]$ where $\ell \leq u$.  The
sentinels $-\infty$ and $+\infty$ represent the absence of a lower or upper
bound, respectively.  A range may additionally be \textit{unknown} (no
information available), in which case the interval is $(-\infty, +\infty)$.

\section{Range propagation rules}
\label{sec:vra.propagation}
\index{range propagation}

\pnum
The following table defines the result range for binary arithmetic
operations on intervals $[a,b]$ and $[c,d]$.  All arithmetic is
saturating: overflow wraps to $\pm\infty$.

\begin{longtable}{@{}lll@{}}
\toprule
\textbf{Operation} & \textbf{Result range} & \textbf{Notes} \\
\midrule
\endfirsthead
\toprule
\textbf{Operation} & \textbf{Result range} & \textbf{Notes} \\
\midrule
\endhead
\bottomrule
\endlastfoot
$[a,b] + [c,d]$ & $[a+c,\; b+d]$ & saturating \\
$[a,b] - [c,d]$ & $[a-d,\; b-c]$ & saturating \\
$[a,b] \times [c,d]$ & $[\min(ac,ad,bc,bd),\;\max(ac,ad,bc,bd)]$ & saturating \\
$[a,b] \div [c,d]$, $c > 0$ & $[a \div d,\; b \div c]$ & integer truncation \\
$[a,b] \div [c,d]$, $0 \in [c,d]$ & \diagcode{E015} & division by zero \\
$[a,b] \mathbin{\%} [c,d]$, $a \geq 0$, $c > 0$ & $[0,\; d-1]$ & \\
$[a,b] \mathbin{\%} [c,d]$, otherwise & unknown & conservative \\
\end{longtable}

\section{Conditional refinement}
\label{sec:vra.refinement}
\index{conditional refinement}

\pnum
After an \lain{if x < N} test, the VRA refines the range of \lain{x}
in the then-branch to $[\ell,\, N-1]$ and in the else-branch to $[N,\, u]$
(where $[\ell, u]$ is the range before the test).

\pnum
At the join point after an \lain{if}/\lain{else}, the range is the union
of the ranges from both branches.

\pnum
When conditions are chained with \lain{and}, refinements from the
left operand are active when evaluating the right operand
(\S\,\ref{sec:expr.in}).

\section{In-guard semantics}
\label{sec:vra.in-guard}
\index{in guard!VRA semantics}

\pnum
An in-guard \lain{idx in arr} (\S\,\ref{sec:expr.in}) introduces the
constraint $0 \leq \mathit{idx} < \mathit{arr.len}$ into the VRA's
range environment.  Inside the guarded body, \lain{arr[idx]} is accepted
without further proof.

\pnum
When \lain{idx in arr} is the condition of a \lain{while decreasing M}
loop, the VRA infers that $\mathit{arr.len} - \mathit{idx} \geq 0$ (a
valid non-negativity proof for the measure $\mathit{arr.len} - \mathit{idx}$).

\section{Constraint annotations}
\label{sec:vra.constraints}
\index{constraint annotation}

\pnum
A constraint annotation of the form \lain{T op bound} on a parameter or
return type declares an invariant that the VRA shall enforce:

\begin{itemize}
  \item \textbf{Parameter constraint}: at every call site, the VRA shall
        prove that the argument satisfies the constraint.
  \item \textbf{Return constraint}: at every \lain{return} statement, the
        VRA shall prove that the return expression satisfies the constraint.
\end{itemize}

\begin{constraint}
  A function call where the argument range cannot be proven to satisfy
  a parameter constraint is \illformed{}~(\diagcode{E012}).
\end{constraint}

\begin{constraint}
  A \lain{return} expression whose range cannot be proven to satisfy the
  declared return constraint is \illformed{}~(\diagcode{E086}) --- the same
  boundary mechanism (\S\,\ref{sec:vra.boundary}) that checks assignments and
  narrowing conversions.  (A \emph{parameter}-constraint violation at a call
  site is the distinct code \diagcode{E012}, above.)
\end{constraint}

\section{Interprocedural range propagation}
\label{sec:vra.interprocedural}
\index{interprocedural VRA}

\pnum
When a call expression appears in a context where the VRA needs a range,
it uses the callee's declared return constraint (if any) as a conservative
approximation.

\begin{example}
\begin{laincode}
func nonzero() int >= 1 { return 42 }

func divide(a int) int {
    return a / nonzero()   // divisor range known >= 1: no [E015]
}
\end{laincode}
\end{example}

\pnum
When the callee has no declared return constraint, the call expression has
an unknown range.

\begin{note}
  Interprocedural VRA is conservative: only declared constraints are used,
  not body analysis.  This preserves decidability and ensures that changing
  a function's body cannot silently affect its callers' safety proofs.
\end{note}

\section{Limits and known gaps}
\label{sec:vra.limits}

\pnum
The following limitations are known and accepted:

\begin{enumerate}
  \item VRA does not track general integer value ranges of struct fields: a
        field's value range is unknown (a field read through a pointer always
        is).  Two narrower field properties \emph{are} tracked --- per-field
        \emph{initialization} (definite assignment, \S\,\ref{sec:own.diag}) and
        a struct field's \emph{slice length} within an \lain{and}-scope (to
        prove indexing of a field slice).
  \item VRA does not track ranges across loop-carried dependences with
        non-linear measures.
  \item Integer overflow in VRA arithmetic uses saturating semantics; the
        actual Lain program uses wrapping arithmetic.  Programs whose
        correctness depends on integer overflow are not analyzed correctly
        by VRA.
  \item Floating-point values are not tracked by VRA.
\end{enumerate}


\section{Loop body affine recap}
\label{sec:vra.affine}
\index{affine loop body}

\pnum
For \lain{for V in start..end} loops whose body contains assignments of
the form \lain{x = x + c} or \lain{x = x - c} with constant step~\(c\),
the implementation computes the post-loop range of~\(x\) precisely:

\begin{align*}
  x_{\text{post}} \in [x_{\text{init}} + c \cdot n_{\min},\; x_{\text{init}} + c \cdot n_{\max}]
\end{align*}

where \(n_{\min}\), \(n_{\max}\) is the iteration count derived from the
loop's range. This avoids the conservative widening that would otherwise
mark \(x\) as unknown after the loop.

\begin{example}
\begin{laincode}
var x = 0
for i in 0..10 {
    x = x + 1
}
// Post-loop: x is exactly 10 (range [10, 10]).
\end{laincode}
\end{example}

\pnum
Patterns not matching the affine form (mixed updates in branches,
non-constant step, multiple variables updated together) fall back to
the conservative widening (unknown post-loop range).

\section{Overflow-at-boundary check}
\label{sec:vra.boundary}
\index{overflow check}

\pnum
The Phase 5 boundary check applies at every integer-type boundary
where a value crosses into a narrower or different-domain target.
The check sites are:
\begin{itemize}
  \item Variable declaration (\lain{var x T = expr}).
  \item Assignment to identifier, field, or indexed element
        (\lain{x = expr}, \lain{p.f = expr}, \lain{arr[i] = expr}).
  \item Return value vs the function's declared return type.
  \item Call argument vs the parameter's declared type.
  \item Struct constructor field initialization.
\end{itemize}

\pnum
At each site, if the source value's VRA range exceeds the target type's
natural range, the program is \illformed{}~(\diagcode{E086}). The
diagnostic offers these resolutions: widen the target; choose an explicit
overflow policy via a cast tier or operator --- checked (\lain{as?},
\lain{+?}), wrapping (\lain{as\%}, \lain{+\%}), or saturating (\lain{as|},
\lain{+|}); or tighten input constraints so VRA can prove safety
(\S\,\ref{sec:types.overflow}).

\pnum
The check is skipped:
\begin{itemize}
  \item Inside an \lain{unsafe} block;
  \item When the source range is unbounded (the analysis cannot rule out
        overflow without further information; the user must annotate);
  \item For integer-literal sources (literals are polymorphic across
        \lain{iN}/\lain{uN} and trusted to fit when written explicitly).
\end{itemize}

\pnum
The wrapping/saturating operators (\tok{+\%}, \tok{*\%}, \tok{+|}, \tok{*|})
produce a result whose range is automatically clamped to the LHS type's
natural range, so they pass the boundary check unconditionally.

\begin{example}
\begin{laincode}
var a u8 = 200
var b u8 = 100
// var bad u8 = a + b      // [E086]: range [300, 300] exceeds u8
var w u8 = a +% b           // OK: wrapping clamps to [0, 255]
var s u8 = a +| b           // OK: saturating clamps to [0, 255]
\end{laincode}
\end{example}

\section{Sub-slice bounds}
\label{sec:vra.subslice}
\index{sub-slice}
\index{slicing}

\pnum
A sub-slice expression \lain{arr[a..b]} produces a slice covering the
half-open range $[a, b)$ of \lain{arr}.  The inferred type of the
result is \lain{T[b-a]}, a sized slice
(\S\,\ref{sec:types.sized-slice}) whose length is statically known.
When both endpoints are compile-time integer literals, the compiler
verifies $a \leq b$; if the check fails the program is
\illformed{}~(\diagcode{E087}).

\pnum
When the endpoints are variables, each is range-checked against the
container length (\diagcode{E085}, the standard bounds check). A
runtime guarantee on the ordering of variable endpoints is not
emitted at compile time; the programmer must establish it through
refinement constraints or by limiting the range expression to
literal bounds.

\begin{example}
\begin{laincode}
var arr i32[10] = [0,1,2,3,4,5,6,7,8,9]
var ok  = arr[2..5]    // type i32[3], half-open, length 3
// var bad = arr[5..2] // [E087]: start 5 > end 2
var empty = arr[3..3]  // OK: empty slice of length 0
\end{laincode}
\end{example}

\section{Sized slice call-site verification}
\label{sec:vra.sized-slice-callsite}
\index{sized slice!call-site check}
\index{E087}

\pnum
When a function parameter is declared with a sized slice type
(\S\,\ref{sec:types.sized-slice}), the compiler verifies at every
call site that the supplied argument satisfies the parameter's
length constraint.  A proved violation is \illformed{}~(\diagcode{E087}).
The check is conservative: it fires only when VRA can \emph{prove}
the constraint is not met; it is silent when the constraint cannot
be determined statically.

\pnum
The cases checked are:

\begin{itemize}
  \item \lain{T[>= k]} or \lain{T[> k]} with integer literal \lain{k}:
        fires if the argument length's VRA upper bound is less than
        the required minimum (\lain{>= k} requires $\mathit{len} \geq k$;
        \lain{> k} requires $\mathit{len} \geq k+1$).
  \item \lain{T[k]} (exact length, literal \lain{k}): fires if the
        argument length's VRA range does not contain \lain{k}.
  \item \lain{T[ref.len]} (cross-parameter equality): fires if both
        sides have a concrete singleton VRA range and they differ.
\end{itemize}

\begin{example}
\begin{laincode}
func pair_sum(arr i32[>= 2]) i32 { return arr[0] + arr[1] }

proc main() i32 {
    var one i32[1] = [42]
    // pair_sum(one)        // [E087]: len 1 < required 2
    var two i32[2] = [1, 2]
    return pair_sum(two)    // OK: len 2 satisfies >= 2
}
\end{laincode}
\end{example}

\pnum
\textbf{Sized slice VRA entry.}
At function entry, for each parameter \lain{name T[expr]} the compiler
seeds the range table with the constraint derived from the size
expression:

\begin{itemize}
  \item \lain{T[k]} (literal): \texttt{\_\_len\_name} $\in [k, k]$.
  \item \lain{T[>= k]} (literal): \texttt{\_\_len\_name} $\in [k, +\infty)$.
  \item \lain{T[> k]} (literal): \texttt{\_\_len\_name} $\in [k{+}1, +\infty)$.
  \item \lain{T[a.len]} (member reference): difference constraint
        \texttt{\_\_len\_name} $=$ \texttt{\_\_len\_a}.
  \item \lain{T[a.len + b.len]} (arithmetic): constraints
        \texttt{\_\_len\_a} $\leq$ \texttt{\_\_len\_name} and
        \texttt{\_\_len\_b} $\leq$ \texttt{\_\_len\_name}.
\end{itemize}

These constraints are available to the VRA throughout the function body,
enabling the Omega Test (\S\,\ref{sec:vra.omega}) to discharge complex
indexing patterns.

\section{VRA analysis levels}
\label{sec:vra.levels}
\index{VRA!analysis levels}

\pnum
Lain's VRA applies four escalating analysis levels at each bounds-check
site.  The levels are tried in order; if one succeeds (the check is proved
safe), the remaining levels are skipped.

\begin{description}[leftmargin=2em]
  \item[L1 — Interval arithmetic.]
    Each variable is represented by an interval $[\ell, u]$.
    The check passes if the interval of the index falls within
    $[0, \mathit{arr.len} - 1]$ directly.

  \item[L2 — Difference constraints.]
    A system of difference constraints of the form
    $x - y \leq c$ is maintained.  The check passes if the constraint
    system can derive the required bound in one step.

  \item[L3 — One-step bridge (transitivity).]
    Two-hop transitivity: if $v_1 - X \leq a$ and $X - v_2 \leq b$
    are both known, the checker derives $v_1 - v_2 \leq a + b$ without
    running the full Fourier-Motzkin algorithm.  Handles patterns such
    as \lain{i < src.len - 1} via the chain
    $i < \mathit{src.len}$ and $\mathit{src.len} - 1 < \mathit{src.len}$.

  \item[L4 — Fourier-Motzkin linear arithmetic (Omega Test).]
    Full linear arithmetic elimination over all VRA variables.
    See \S\,\ref{sec:vra.omega}.
\end{description}

\pnum
All four levels are polynomial time.  No SMT solver is used.

\section{Omega Test — Fourier-Motzkin prover (VRA level 4)}
\label{sec:vra.omega}
\index{Omega Test}
\index{Fourier-Motzkin}
\index{VRA!level 4}

\pnum
When L1–L3 fail to discharge a bounds check, the implementation invokes
the \textit{Omega Test}: a self-contained Fourier-Motzkin linear
arithmetic prover (roughly 410 lines in \texttt{src/sema/omega.h}, no external
dependencies).

\pnum
\textbf{Proof method (proof by contradiction).}
To prove $\mathit{index} < \mathit{bound}$:

\begin{enumerate}
  \item Decompose both expressions into linear forms
        $\sum_i c_i \cdot x_i + k$ over named variables.
        References to \lain{x.len} are mapped to the synthetic VRA
        variable \texttt{\_\_len\_x}.
  \item Form the negated query: $\mathit{bound} - \mathit{index} \leq 0$.
  \item Extract all VRA range and difference constraints that mention
        the variables appearing in the query.
  \item Run Fourier-Motzkin variable elimination.  Each step crosses
        upper-bound inequalities with lower-bound inequalities to
        produce a system with one fewer variable.
  \item If the resulting constant system contains an inequality
        $0 \leq c$ with $c < 0$, the system is UNSAT: the negation
        is impossible, so $\mathit{index} < \mathit{bound}$ always holds.
\end{enumerate}

\pnum
The prover also supports non-negativity queries: to prove
$\mathit{expr} \geq 0$, it shows that $\mathit{expr} \leq -1$, combined
with VRA constraints, is UNSAT.

\pnum
\textbf{Example patterns proved by the Omega Test.}

\begin{example}
\begin{laincode}
// concat: out[j + a.len] with j < b.len and out.len == a.len + b.len
// FM eliminates j; the cross-pair (j < b.len, j >= b.len) gives 0 <= -1 (UNSAT).
func concat(a i32[], b i32[], out i32[a.len + b.len]) {
    for j in 0..b.len {
        out[j + a.len] = b[j]   // proved safe by Omega
    }
}

// reverse: out[n - i - 1] with n = arr.len and i < n
// FM eliminates n and __len_arr; reduces to (i <= -1) AND (i >= 0) (UNSAT).
func reverse(arr i32[], out i32[arr.len]) {
    var n = arr.len
    for i in 0..n {
        out[n - i - 1] = arr[i]   // proved safe by Omega
    }
}
\end{laincode}
\end{example}

\pnum
\textbf{Equality propagation for local length aliases.}
When the body contains \lain{var n = x.len}, the compiler registers the
equality $n - \mathtt{\_\_len\_x} = 0$ (expressed as two difference
constraints) so the Omega Test can cancel \lain{n} and
\texttt{\_\_len\_x} in the same linear form.

\pnum
\textbf{Conservatism.}
If the FM buffer overflows (more than 80 inequalities in the working
system), the prover gives up and returns \textit{unknown}; the check
then falls through to E085 as if L4 had not fired.  This bound is
unreachable for the small systems (2--6 variables, 5--15 inequalities)
typical of slice manipulation code.

% 14-generics.tex — Clause 14: Generics and compile-time evaluation
\chapter{Generics and compile-time evaluation}
\label{sec:14-generics}
\index{generics}
\index{compile-time evaluation}
\index{comptime}

\section{Overview}
\label{sec:gen.overview}

\pnum
Lain implements genericity through compile-time function evaluation (CTFE).
Generic types and functions accept \textit{type parameters} (declared as
\lain{name type}) or \textit{compile-time value parameters} (declared as
\lain{comptime name non-type}).  The compiler evaluates them at compile time
and monomorphizes the result.  There is no type erasure and no runtime
polymorphism.

\section{Type parameters}
\label{sec:gen.type-params}
\index{type parameter}

\pnum
A parameter declared \lain{name type} is a type parameter.  Types are
implicitly compile-time — the \lain{comptime} keyword shall \emph{not}
precede a type parameter.  At every call site, the argument shall be a type
expression known at compile time.

\begin{laincode}
func identity(T type, x T) T {
    return x
}

var n = identity(int, 42)       // T = int, monomorphized
var s = identity(u8[:0], "hi")  // T = u8[:0], separate monomorphization
\end{laincode}

\pnum
Type parameters are evaluated entirely at compile time; they produce
no runtime code.

\begin{constraint}
  The \lain{comptime} keyword shall not appear before a parameter whose
  type is \lain{type}.  \lain{func f(comptime T type)} is ill-formed;
  write \lain{func f(T type)} instead.
\end{constraint}

\section{Compile-time value parameters}
\label{sec:gen.comptime-params}
\index{comptime parameter}

\pnum
A parameter declared \lain{comptime name non-type} is a compile-time value
parameter.  The argument shall be a constant expression of the declared type
known at compile time.

\begin{laincode}
type RingBuffer(T type, comptime N int) = type {
    items T[N]
    head  usize
    tail  usize
}
\end{laincode}

\begin{constraint}
  A compile-time evaluation context (CTFE) shall not call a \lain{proc}
  or an \lain{extern proc}; only \lain{func} calls are admissible. A
  violation is \illformed{}~(\diagcode{E101}).
\end{constraint}

\section{Type as value}
\label{sec:gen.type-value}
\index{type as value}

\pnum
A \lain{func} with return type \lain{type} is a type constructor.  Its
body shall consist only of expressions and statements that are evaluable
at compile time.  The return value is a type that the caller may use in
type annotations and variable declarations.

\begin{laincode}
type Option(T type) = type {
    Some { value T }
    None
}

var x = Option(int).Some(42)
\end{laincode}

\pnum
Type constructors may also be written as \lain{func} declarations returning
\lain{type}:

\begin{laincode}
func Option(T type) type {
    return type {
        Some { value T }
        None
    }
}
\end{laincode}

\begin{constraint}
  A function returning \lain{type} shall be declared as \lain{func} (not
  \lain{proc}).
\end{constraint}

\section{Monomorphization}
\label{sec:gen.monomorphization}
\index{monomorphization}

\pnum
For each unique combination of comptime arguments at call sites, the
implementation produces a separate instantiation of the function.  The
name of each instantiation in the translation output is
\idbehav{determined by the implementation}.

\pnum
Monomorphization is a compile-time-only operation; at runtime, each
call site calls a specific, non-generic function.

\section{Comptime expression blocks}
\label{sec:gen.comptime-block}
\index{comptime block}

\pnum
A \lain{comptime \{ expr \}} expression evaluates \textit{expr} at compile
time.  The result is a compile-time constant that may be used as a type
annotation, array size, or parameter constraint value.

\begin{constraint}
  A \lain{comptime} expression shall be pure and terminating (i.e., it may
  only use \lain{func} calls and the constructs permitted in \lain{func}
  bodies).
\end{constraint}

\section{Comptime if}
\label{sec:gen.comptime-if}
\index{comptime if}

\pnum
\lain{comptime if cond \{ ... \} else \{ ... \}} is a compile-time
conditional.  Only the taken branch is type-checked and emitted; the other
branch is entirely discarded.  \textit{cond} shall evaluate to a
compile-time boolean constant.

\begin{note}
  \lain{comptime if} enables platform-specific code selection without
  dead-code warnings.  It is the preferred way to write code that
  varies based on compile-time parameters.
\end{note}

\section{Comptime recursion depth}
\label{sec:gen.recursion-depth}
\index{comptime recursion depth}

\pnum
The maximum depth of comptime function evaluation is
\idbehav{the implementation's comptime evaluation depth limit; the
reference implementation enforces a limit of 256}.  A comptime evaluation
that exceeds this depth is \illformed{}~(\diagcode{E014}).

\section{Mode preservation in generics}
\label{sec:gen.mode-preserve}
\index{mode preservation}

\pnum
When a type parameter \lain{T} is substituted with a concrete type, the
ownership mode of the parameter is preserved.  If a parameter was declared
\lain{var T}, and \lain{T} is substituted with a type that has mode
\lain{shared} in its declaration, the substituted parameter retains the
\lain{var} mode.

\begin{note}
  This rule ensures that mode annotations in generic functions behave
  consistently regardless of the concrete type argument.
\end{note}

% 15-match.tex — placeholder
\chapter{Match}
\label{sec:15-match}

% TODO: write this chapter.

% 16-modules.tex — Clause 16: Module system
\chapter{Module system}
\label{sec:16-modules}
\index{module system}

\section{Overview}
\label{sec:mod.overview}

\pnum
Each Lain source file defines one module.  Modules are the unit of code
organization and compilation.  All names declared at the top level of a
source file are in that module's scope and are accessible from other
modules that import the module.

\section{Module identity}
\label{sec:mod.declaration}
\index{module declaration}

\pnum
A module's identity is \idbehav{derived from its file path by the
implementation}; there is \emph{no} \lain{module} declaration keyword (a
leading \lain{module}~\dots{} line is \diagcode{E100}).  The reference
implementation names a module by its path relative to the working directory,
with directory separators mapped to dots --- so \texttt{std/io.ln} is the
module \lain{std.io}.

\section{Import declarations}
\label{sec:mod.import}
\index{import declaration}

\pnum
\lain{import module.path} makes the declarations of the named module
available in the current source file.  The grammar production is
\gramref{import-declaration} (Annex~A).

\pnum
\textbf{Selective import.}  \lain{import module.path.\{name1, name2, ...\}}
brings only the listed names into scope rather than the whole module.  This is
the form used throughout the standard library, e.g.\
\lain{import std.mem.\{mem\_alloc, mem\_free\}} or
\lain{import std.fs.\{File, open\_file, close\_file\}}.

\pnum
The \lain{use} statement (\S\,\ref{sec:lex.keywords}) provides a
\emph{local-scope} import/alias inside a function body.  \idbehav{Its precise
semantics are implementation-defined in this revision}; the standard library
uses top-level \lain{import} rather than \lain{use}.

\pnum
\textbf{Path resolution}: a module path is converted to a file system path
by replacing dots with directory separators and appending \texttt{.ln}.
The search is relative to the compiler's working directory.
\idbehav{The exact search algorithm, including support for search paths,
is implementation-defined}.

\begin{longtable}{@{}ll@{}}
\toprule
\textbf{Module path} & \textbf{File path} \\
\midrule
\endfirsthead
\toprule
\textbf{Module path} & \textbf{File path} \\
\midrule
\endhead
\bottomrule
\endlastfoot
\lain{std.c}        & \texttt{std/c.ln}        \\
\lain{std.io}       & \texttt{std/io.ln}        \\
\lain{std.fs}       & \texttt{std/fs.ln}        \\
\lain{std.math}     & \texttt{std/math.ln}      \\
\lain{mylib.utils}  & \texttt{mylib/utils.ln}   \\
\end{longtable}

\section{Module loading}
\label{sec:mod.loading}
\index{module loading}

\pnum
When the compiler encounters an \lain{import} declaration:
\begin{enumerate}
  \item The module path is converted to a file path.
  \item If the module has already been loaded, the import is a no-op
        (deduplication).
  \item Otherwise, the file is read, lexed, parsed, and recursively
        resolved.
  \item The loaded module's declarations are merged into the importing
        module's declaration list.
\end{enumerate}

\pnum
Each module is loaded at most once.  Circular imports are therefore
implicitly handled: if module~A imports module~B which imports module~A,
the second import of~A is a no-op because A is already loaded.

\begin{note}
  The reference implementation maintains a linked list of loaded modules
  keyed by their dot-separated name.
\end{note}

\section{Name visibility}
\label{sec:mod.visibility}
\index{name visibility}

\pnum
By default, every name declared at module scope is \textit{public}:
visible to any module that imports the declaring module.  This default
matches the typical usage where module declarations form an API surface.

\pnum
A declaration prefixed with the attribute \tok{[private]}
(\S\,\ref{sec:10-declarations}) is \textit{module-private}: it is
visible only within the defining module and may not be referenced from
any importing module.

\begin{laincode}
// In module foo.bar:
func public_api() int { ... }    // public by default

[private]
func internal_helper() int { ... }   // only foo.bar may reference this
\end{laincode}

\begin{constraint}
  \textbf{Q-018.}
  A reference to a \tok{[private]} declaration from a module other than
  the one where it is defined is \illformed{}~(\diagcode{E084}).  The
  reference may appear in any expression position---identifier lookup,
  member access, type instantiation, etc.
\end{constraint}

\begin{note}
  Compilation under the reference implementation translates each
  \tok{[private]} declaration to a C declaration with internal linkage
  (\texttt{static} storage class), and each public declaration to one
  with external linkage.
\end{note}

\section{Standard library module paths}
\label{sec:mod.stdlib}
\index{standard library!module paths}

\pnum
The following module paths are defined by this specification for the
standard library:

\begin{longtable}{@{}ll@{}}
\toprule
\textbf{Module path} & \textbf{Description} \\
\midrule
\endfirsthead
\toprule
\textbf{Module path} & \textbf{Description} \\
\midrule
\endhead
\bottomrule
\endlastfoot
\lain{std.c}    & C standard library bindings \\
\lain{std.io}   & Console I/O \\
\lain{std.fs}   & File system operations \\
\lain{std.math} & Mathematical functions \\
\end{longtable}

\section{Name mangling}
\label{sec:mod.mangling}
\index{name mangling}

\pnum
The name of a Lain declaration in the translation output is
\idbehav{determined by the implementation}.  The reference implementation
uses the function name prefixed with the module path, with dots replaced
by underscores.  Name collisions between modules are
\idbehav{detected by the implementation; the reference implementation
appends a numeric suffix to resolve conflicts}.

% 17-interop.tex — Clause 17: C interoperability
\chapter{C interoperability}
\label{sec:17-interop}
\index{C interoperability}

\section{Overview}
\label{sec:interop.overview}

\pnum
Lain compiles to C99 and provides direct, zero-overhead interoperability
with C libraries through three mechanisms:
\begin{enumerate}
  \item \lain{c\_include} --- injects C \texttt{\#include} directives;
  \item \lain{extern} declarations --- declares C functions for use in Lain;
  \item type mapping --- Lain types map to their C99 counterparts.
\end{enumerate}

\section{The \texttt{c\_include} declaration}
\label{sec:interop.cinclude}
\index{c\_include}

\pnum
\lain{c\_include "header.h"} causes the compiler to emit a
\texttt{\#include} directive in the generated C code.  The string content
is emitted verbatim.

\begin{laincode}
c_include "<stdio.h>"
c_include "<stdlib.h>"
\end{laincode}

\begin{constraint}
  \lain{c\_include} shall appear only at module top level.
\end{constraint}

\section{Extern function declarations}
\label{sec:interop.extern}
\index{extern function}

\pnum
\lain{extern func name(params) return} declares a C function.  Because an
extern has no body, its effect row cannot be inferred; the row is therefore
\emph{believed}, and its default is the set of \emph{all} effects
(\lain{io}, \lain{diverge}, \lain{raises}, \lain{alloc}).  An
\lain{effects} clause on an extern \emph{narrows} that default.

\pnum
This is the inverse of a declaration with a body, whose default row is
$\emptyset$ and whose clause \emph{widens} it~(\S\,\ref{sec:decl.func}).
The direction differs for one stated reason: inference is available in the
one case and impossible in the other.  A caller of an extern acknowledges
whatever the extern's row names, by the ordinary rule.

\begin{constraint}
  An \lain{extern} with no \lain{effects} clause has the row
  \{\lain{io}, \lain{diverge}, \lain{raises}, \lain{alloc}\}.  A caller
  whose own bound does not cover it is \illformed{}~(\diagcode{E011}).
\end{constraint}

\begin{note}
  \textbf{Trust boundary.}
  An extern's \lain{effects} clause is a \textit{programmer assertion} that
  the referenced C function does no more than the row names.  The compiler
  cannot verify it and believes it, so a wrong narrowing---declaring
  \lain{libc\_rand} with an empty row, say---may invalidate the guarantees
  that depend on purity (notably P4 determinism for any function that
  transitively calls it).  What has changed is the cost of \emph{silence}:
  omitting the clause once meant "trusted pure" (the rule once labelled I-016), so
  forgetting it on a C function that performs I/O was a silent fail-open.
  Omitting it now claims everything, so forgetting costs precision instead.
  The empty row, \lain{effects} with no names, is how purity is asserted
  deliberately.
\end{note}

\pnum
Ownership annotations on extern parameters work identically to those on
native Lain functions; the borrow checker tracks ownership of resources
obtained via extern calls.

\section{Type mapping}
\label{sec:interop.types}
\index{type mapping}

\pnum
The following table defines the normative mapping between Lain types and
their C99 equivalents.  The mapping is the ABI used by the translation
output.

\begin{longtable}{@{}lll@{}}
\toprule
\textbf{Lain type} & \textbf{C99 type} & \textbf{Notes} \\
\midrule
\endfirsthead
\toprule
\textbf{Lain type} & \textbf{C99 type} & \textbf{Notes} \\
\midrule
\endhead
\bottomrule
\endlastfoot
\lain{u8}          & \texttt{uint8\_t}   & \\
\lain{u16}         & \texttt{uint16\_t}  & \\
\lain{u32}         & \texttt{uint32\_t}  & \\
\lain{u64}         & \texttt{uint64\_t}  & \\
\lain{usize}       & \texttt{size\_t}    & pointer-sized \\
\lain{i8}          & \texttt{int8\_t}    & \\
\lain{i16}         & \texttt{int16\_t}   & \\
\lain{i32}         & \texttt{int32\_t}   & \\
\lain{i64}         & \texttt{int64\_t}   & \\
\lain{isize}       & \texttt{ptrdiff\_t} & pointer-sized \\
\lain{int}         & \texttt{int}        & at least 32-bit \\
\lain{f32}         & \texttt{float}      & IEEE 754 \\
\lain{f64}         & \texttt{double}     & IEEE 754 \\
\lain{bool}        & \texttt{\_Bool}     & C99 bool (1 byte); \idbehav{0 or 1} \\
shared struct \lain{T} & \texttt{const T*} & \\
\lain{var} struct \lain{T} & \texttt{T*}  & \\
\lain{mov} struct \lain{T} & \texttt{T}   & by value \\
\lain{T[]}         & \texttt{struct \{ size\_t len; T* data; \}} & fat slice \\
\lain{*T}          & \texttt{const T*}  & unsafe only \\
\lain{var *T}      & \texttt{T*}        & unsafe only \\
\end{longtable}

\section{ABI for function parameters}
\label{sec:interop.abi}
\index{ABI}

\pnum
This table is normative and applies to all Lain function calls, including
those to \lain{extern} functions.

\begin{longtable}{@{}llll@{}}
\toprule
\textbf{Lain param} & \textbf{Type} & \textbf{C param} & \textbf{Direction} \\
\midrule
\endfirsthead
\toprule
\textbf{Lain param} & \textbf{Type} & \textbf{C param} & \textbf{Direction} \\
\midrule
\endhead
\bottomrule
\endlastfoot
\lain{name T}     & struct & \texttt{const T*}  & in \\
\lain{var name T} & struct & \texttt{T*}        & in/out \\
\lain{mov name T} & struct & \texttt{T}         & in (by value) \\
any mode          & primitive & \texttt{T}      & in \\
\end{longtable}

\section{Pointer operations}
\label{sec:interop.pointers}
\index{pointer operations}

\pnum
Raw pointer operations (\lain{*ptr} dereference, \lain{\&x} address-of)
require an \lain{unsafe} block (see \S\,\ref{sec:18-unsafe}).

\begin{constraint}
  Pointer-to-pointer or integer-to-pointer casts shall occur only inside
  an \lain{unsafe} block~(\diagcode{E012}).
\end{constraint}

\section{String interop}
\label{sec:interop.strings}
\index{string interop}

\pnum
Lain strings (\lain{u8[]}) are byte slices, not null-terminated C strings.
To pass a Lain string to a C function expecting \texttt{char*}, the
\lain{.data} pointer must be extracted explicitly:

\begin{laincode}
c_include "<stdio.h>"
extern func puts(s *u8) int effects io

var s = "hello"
puts(s.data)   // OK: .data is *u8 = const char*
\end{laincode}

\pnum
String literals are null-terminated (\S\,\ref{sec:lex.literals.string}),
so \lain{literal.data} is a valid \texttt{const char*} for C functions
that expect a C string.

\section{Opaque C types}
\label{sec:interop.opaque}
\index{opaque type}

\pnum
An opaque C type is declared with \lain{extern type Name}.  The type's
layout is unknown to Lain; it may only be used behind a pointer.

\begin{laincode}
extern type FILE
extern func fopen(path *u8, mode *u8) mov *FILE effects io, alloc
\end{laincode}

\begin{constraint}
  An opaque type shall not be used as a value type (only as a pointer
  target).  \lain{var f FILE} is \illformed{}.
\end{constraint}

% 18-unsafe.tex — Clause 18: Unsafe code
\chapter{Unsafe code}
\label{sec:18-unsafe}
\index{unsafe code}

\section{Purpose}
\label{sec:unsafe.purpose}

\pnum
The \lain{unsafe} block provides a controlled escape hatch for operations
that a conforming implementation cannot statically verify.  Unsafe blocks
are explicitly delimited in the source text and are therefore auditable.
Unsafe code cannot accidentally spread outside the enclosing block.

\section{Syntax}
\label{sec:unsafe.syntax}

\pnum
An unsafe block is written as \lain{unsafe \{ body \}}.  It may appear
wherever a statement is valid (\S\,\ref{sec:stmt.unsafe}).

\section{What unsafe enables}
\label{sec:unsafe.enables}
\index{unsafe!what it enables}

\pnum
Inside an \lain{unsafe} block, the following operations are permitted:

\begin{enumerate}
  \item \textbf{Raw pointer dereference}: \lain{*ptr} reads the value
        pointed to by a raw pointer; \lain{*ptr = v} writes through a
        mutable raw pointer.
  \item \textbf{Address-of}: \lain{\&x} creates a raw pointer to the
        storage of \lain{x}.
  \item \textbf{Pointer and integer casts}: \lain{p as *T},
        \lain{n as *T}, \lain{p as usize}, etc.
  \item \textbf{Direct ADT field access}: accessing a field of an ADT
        variant without going through a \lain{case} arm.
  \item \textbf{Bounds-check suppression}: array indexing inside an
        \lain{unsafe} block does not require a VRA proof (the
        programmer asserts the access is in bounds).
\end{enumerate}

\section{What unsafe does NOT disable}
\label{sec:unsafe.not-disabled}
\index{unsafe!what it does not disable}

\pnum
The following safety features remain active inside \lain{unsafe} blocks:

\begin{longtable}{@{}ll@{}}
\toprule
\textbf{Feature} & \textbf{Disabled by \lain{unsafe}?} \\
\midrule
\endfirsthead
\toprule
\textbf{Feature} & \textbf{Disabled by \lain{unsafe}?} \\
\midrule
\endhead
\bottomrule
\endlastfoot
Pointer dereference                  & \textbf{Yes} \\
Pointer/integer casts                & \textbf{Yes} \\
Direct ADT field access              & \textbf{Yes} \\
Bounds check requirement             & \textbf{Yes} \\
Ownership tracking (\diagcode{E001}--\diagcode{E007}) & No \\
Borrow checking (\diagcode{E004})    & No \\
Move semantics                       & No \\
Type checking                        & No \\
Division-by-zero check (\diagcode{E015}) & No \\
Exhaustiveness checking (\diagcode{E014}) & No \\
Purity enforcement (\diagcode{E011}) & No \\
\end{longtable}

\begin{constraint}
  Ownership and linearity rules (\S\,\ref{sec:11-ownership}) are not
  suspended inside \lain{unsafe} blocks.  Use-after-move, double-move,
  and borrow conflicts are still constraint violations inside
  \lain{unsafe}~(\diagcode{E001}--\diagcode{E007}).
\end{constraint}

\section{Undefined behavior in unsafe}
\label{sec:unsafe.ub}
\index{undefined behavior!in unsafe}

\pnum
The following operations inside an \lain{unsafe} block may produce
\undefbehav{undefined behavior if the programmer's safety assertion is
incorrect}:

\begin{itemize}
  \item Dereferencing a null pointer.
  \item Dereferencing a pointer to freed or uninitialized memory.
  \item Writing to memory through a pointer derived from a read-only
        (shared) source.
  \item Accessing an array element with an out-of-range index.
  \item Accessing a field of an ADT through a tag that does not match
        the active variant.
\end{itemize}

\pnum
A conforming implementation that generates C99 translation output inherits
the undefined behavior rules of ISO/IEC~9899 (C17) for operations within
\lain{unsafe} blocks.

\section{Borrow scope for unsafe blocks}
\label{sec:unsafe.borrow-scope}

\pnum
An \lain{unsafe} block is a regular block scope for the purposes of borrow
checking.  Borrows registered within the block are released on block exit,
and the borrow checker's invariants hold at the block boundaries.

\section{Nesting}
\label{sec:unsafe.nesting}

\pnum
\lain{unsafe} blocks may be nested.  The inner block does not extend the
permissions of the outer block beyond what is already enabled.
\idbehav{Whether the implementation issues a diagnostic for redundant nested
unsafe blocks is implementation-defined}.

% 19-memory.tex — Clause 19: Memory model
\chapter{Memory model}
\label{sec:19-memory}
\index{memory model}

\section{Overview}
\label{sec:mem.overview}

\pnum
Lain's memory model is founded on three invariants:
\begin{enumerate}
  \item \textbf{Deterministic allocation} — every allocation is explicit
        and predictable; the implementation never implicitly allocates
        heap memory.
  \item \textbf{Zero runtime overhead} — there is no garbage collector
        and no reference-counting mechanism.
  \item \textbf{C99 layout compatibility} — the physical layout of all
        objects follows ISO/IEC~9899 (C17) rules exactly.
\end{enumerate}

\pnum
This clause specifies the abstract notions of \textit{object},
\textit{storage duration}, \textit{object representation}, and
\textit{address}.  The abstract machine described here is realised by any
conforming translation output (\S\,\ref{sec:conformance}).

\section{Objects}
\label{sec:mem.objects}
\index{object}

\pnum
An \textit{object} is a contiguous region of storage that holds a value of
a declared type.  Every variable declaration introduces exactly one object.
Fields of a struct introduce sub-objects.  Elements of a fixed-size array
introduce element sub-objects.

\pnum
Every object has:
\begin{itemize}
  \item a \textit{declared type} (\S\,\ref{sec:07-types});
  \item an \textit{ownership mode} — shared, \lain{var}, or \lain{mov}
        (\S\,\ref{sec:own.modes});
  \item a \textit{linear state} — \textsc{Uninit}, \textsc{Unconsumed},
        or \textsc{Consumed} (\S\,\ref{sec:own.states});
  \item a \textit{storage duration} (\S\,\ref{sec:mem.storage});
  \item an \textit{address} — the byte address of its first storage unit,
        aligned as required by the target platform.
\end{itemize}

\section{Storage duration}
\label{sec:mem.storage}
\index{storage duration}

\pnum
Three storage durations are defined:

\begin{description}[leftmargin=2em]
  \item[Automatic duration]
    Objects declared inside a function or block body (including function
    parameters) have automatic storage duration.  Storage is allocated when
    the enclosing scope is entered and released when that scope exits.

  \item[Static duration]
    Objects declared at module scope have static storage duration.  Their
    storage persists for the entire program execution.
    \idbehav{The value of a static-duration object whose declaration
    lacks an initializer is implementation-defined; the reference
    implementation zero-initializes such objects via C99 static
    initialization rules}.

  \item[Heap (explicit) duration]
    Objects allocated through C interop (\lain{malloc} or equivalent)
    have programmer-controlled duration.  The ownership system
    (\S\,\ref{sec:11-ownership}) enforces that \lain{mov} pointers to
    heap memory are consumed exactly once~(\diagcode{E003}).
\end{description}

\begin{constraint}
  A \lain{return} statement shall not return a reference whose referent
  has automatic storage duration~(\diagcode{E010}).
\end{constraint}

\section{Stack allocation}
\label{sec:mem.stack}
\index{stack allocation}

\pnum
All local variables, struct values, and fixed-size arrays are allocated on
the call stack.  No implicit heap allocation is ever performed by the
compiler.

\begin{laincode}
var x int = 42     // stack: sizeof(int32_t) bytes
var arr int[100]    // stack: 100 * sizeof(int32_t) bytes
var p Point         // stack: sizeof(Point) bytes
\end{laincode}

\pnum
Each function call produces a \textit{stack frame} containing:
\begin{itemize}
  \item all local variables declared in the function;
  \item all function parameters (by value or by pointer, per the calling
        convention of \S\,\ref{sec:fn.distinction});
  \item compiler-generated temporaries required for expression evaluation.
\end{itemize}

\pnum
The stack frame is automatically reclaimed when the function returns.
\idbehav{The implementation's behaviour when the call stack is exhausted
is implementation-defined; the reference implementation inherits the
platform's stack overflow signal}.

\section{Heap allocation}
\label{sec:mem.heap}
\index{heap allocation}

\pnum
Heap allocation is performed exclusively through C interop
(\S\,\ref{sec:17-interop}).  The idiomatic pattern declares the C
\texttt{malloc}/\texttt{free} pair as \lain{extern proc}:

\begin{laincode}
c_include "<stdlib.h>"
extern proc malloc(size usize) mov *void
extern proc free(ptr mov *void)
\end{laincode}

\pnum
Because \lain{malloc} returns \lain{mov *void}, the caller receives sole
ownership of the allocated block.  Because \lain{free} consumes
\lain{mov *void}, calling \lain{free} transfers that ownership to the
allocator.  The ownership system therefore enforces statically:
\begin{itemize}
  \item \textbf{No use-after-free}: using the pointer after \lain{free}
        is a use-after-move~(\diagcode{E001}).
  \item \textbf{No double-free}: a second call to \lain{free} on the
        same pointer is a double-consume~(\diagcode{E002}).
  \item \textbf{No leak}: failing to call \lain{free} on a \lain{mov}
        pointer is an unconsumed linear value~(\diagcode{E003}).
\end{itemize}

\section{Value semantics}
\label{sec:mem.value-semantics}
\index{value semantics}

\pnum
The behaviour of assignment and function argument passing is determined
by the type's linearity and the ownership mode:

\begin{longtable}{@{}lll@{}}
\toprule
\textbf{Type} & \textbf{Assignment \texttt{b = a}} & \textbf{Notes} \\
\midrule
\endfirsthead
\toprule
\textbf{Type} & \textbf{Assignment \texttt{b = a}} & \textbf{Notes} \\
\midrule
\endhead
\bottomrule
\endlastfoot
Primitives (\lain{int}, \lain{bool}, \ldots) & Bitwise copy & Always allowed \\
Non-linear struct                             & Field-by-field copy & Allowed \\
Linear struct (has \lain{mov} field)          & Forbidden & Requires \lain{mov} \\
Non-linear array                              & Element-wise copy & Allowed \\
Linear array                                  & Forbidden & Requires \lain{mov} \\
Slice (\lain{[T]})                            & Copy \texttt{\{data, len\}} pair & Always allowed \\
Raw pointer (\lain{*T})                       & Copy address & Always allowed \\
\end{longtable}

\begin{constraint}
  Copying a linear type without an explicit \lain{mov} expression is
  \illformed{}~(\diagcode{E006}).
\end{constraint}

\section{Object representation}
\label{sec:mem.representation}
\index{object representation}

\pnum
The \textit{object representation} of a Lain type is its C99 equivalent
as defined in \S\,\ref{sec:interop.types}.  Every object is a sequence of
\lain{u8} bytes whose count equals the C99 \texttt{sizeof} of the
corresponding C type.

\pnum
\textit{Value equality} is defined per type:
\begin{itemize}
  \item Primitives: bitwise equality of the value representation.
  \item Structs: field-by-field value equality.
  \item Slices: equality of both the data pointer and the length.
  \item Raw pointers: address equality.
  \item ADTs: tag equality and, when tags match, field-by-field equality
        of the active variant.
\end{itemize}

\section{Alignment and layout}
\label{sec:mem.alignment}
\index{alignment}
\index{struct layout}

\pnum
Lain inherits ISO/IEC~9899 (C17) alignment and padding rules:
\begin{itemize}
  \item Struct fields are laid out in declaration order.
  \item Each field is aligned to its natural alignment (the alignment of
        its C99 equivalent type).
  \item The compiler inserts padding between fields as necessary.
  \item The total size of a struct is rounded up to a multiple of its
        largest field alignment.
\end{itemize}

\pnum
Lain does not provide layout control annotations.  There is no
\texttt{\#pragma pack} equivalent, no \texttt{alignas}, and no
\texttt{repr(packed)}.  The layout is always the platform-default C99
layout.

\idbehav{The exact size and alignment of each type is implementation-defined,
following the C99 ABI of the target platform}.

\section{ADT and enum layout}
\label{sec:mem.adt-layout}
\index{ADT layout}
\index{enum layout}

\pnum
An enum type is emitted as a C \texttt{enum} wrapped in a named struct so
that each enum type is nominally distinct in the C translation:

\begin{verbatim}
typedef enum { Color_Red, Color_Green, Color_Blue } Color_Tag;
typedef struct { Color_Tag tag; } Color;
\end{verbatim}

\pnum
An ADT type uses a niche-based encoding whenever a zero-cost
representation is derivable. The compiler performs \textit{aggressive
multi-slot niche optimization} (D-Niche): it explores the full
sentinel space of the payload and packs as many empty variants as
fit, without any tag byte. When the pool is insufficient to cover
every empty variant, the layout falls back to a tagged union and the
compiler emits \diagcode{W120} (best-effort warning; not an error).

\begin{verbatim}
// Tagged-union fallback (only when no niche is derivable):
typedef struct {
    enum { Shape_Circle_tag, Shape_Rect_tag } tag;
    union {
        struct { float r; } Circle;
        struct { float w; float h; } Rect;
    };
} Shape;
\end{verbatim}

\pnum
\textit{Niche encoding (zero-cost).}  When the payload of one variant
contains a bit pattern that cannot occur in valid values---for instance,
\texttt{0x0} for a pointer field, or a value outside a VRA-proven
range for an integer field---that pattern is reserved to encode an
alternative variant.  No explicit tag is emitted.

\begin{verbatim}
// Niche encoding for Option(*File):  size = sizeof(*File)
typedef File *Option_Ptr_File;
//  Some(p)   -> non-null pointer p
//  None      -> 0x0
\end{verbatim}

\pnum
The sentinel pool size depends on the target. For pointer payloads,
the pool is the zero-page-unmappable region of the target's address
space, stepped by the pointer alignment (e.g.\ 16384 slots for
x86\_64-linux with alignment~8). For integer payloads, the pool is
the type's full range minus any VRA-narrowed sub-range:
\lain{type Pct = i32 >= 0 and <= 100} as a payload yields a pool
of $\approx 4.3 \times 10^9$ slots (every integer outside $[0,100]$).
For \lain{bool}, the pool is the 254 byte patterns outside $\{0, 1\}$.

\pnum
Niche allocation cascades through nested enum payloads. If
\lain{Option(*T)} reserves the value~\texttt{0} for \lain{None}, a
surrounding \lain{Result(Option(*T), Err4)} sees the remaining slots
$(8, 16, 24, \ldots)$ and packs its own empty variants into them.
Layout of the entire result still occupies one pointer-sized word.

\pnum
\textit{Multi-payload niche.} An enum with exactly two payload variants
where one payload (the \textit{primary}) has a non-empty sentinel pool
and the other payload (the \textit{secondary}) has a single field whose
type is an \textit{enum of pure-empty variants} is encoded by mapping
each sub-variant into a sentinel of the primary's pool. For
\lain{Result(*T, FileErr)} with \lain{FileErr = enum \{ NotFound,
Permission, OOM \}}, the encoding on a host with pointer alignment~8 is
\lain{Ok(f)} = any valid pointer, \lain{Err(NotFound)} = sentinel~0,
\lain{Err(Permission)} = sentinel~8, \lain{Err(OOM)} = sentinel~16. The
secondary constructor emits a stride-multiplication; the case-match
discriminator emits a corresponding stride-division for binding-style
patterns.

\pnum
The compiler never refuses compilation to satisfy a layout
optimization: when the pool is insufficient, \diagcode{W120} is
emitted with suggestions (constrain the payload, change to a type
with larger sentinel space, or split the enum), and a tag byte is
added. The programmer decides whether to accept the overhead.

\pnum
The CLI flag \texttt{--target=<triple>} selects the target's sentinel
pool size (host auto-detection is the default). Supported triples
include \texttt{x86\_64-linux-gnu}, \texttt{aarch64-linux-gnu},
\texttt{x86\_64-windows-msvc}, and \texttt{cortex-m4-bare} (the
last has zero pool: no MMU means address~0 may be a valid vector
table entry).

\section{No garbage collector}
\label{sec:mem.no-gc}
\index{garbage collector!absence of}

\pnum
Lain has no garbage collector.  Consequently:
\begin{itemize}
  \item execution is never interrupted for memory management;
  \item deallocation occurs at statically deterministic program points;
  \item no reference counting or heap tracing is performed at runtime.
\end{itemize}

\pnum
This makes Lain suitable for real-time, embedded, and safety-critical
applications where garbage-collection pauses are unacceptable.

\section{Memory safety in the safe fragment}
\label{sec:mem.safety}
\index{memory safety}

\pnum
In the \textit{safe fragment} (code outside \lain{unsafe} blocks), the
following memory safety properties hold by construction:

\begin{longtable}{@{}ll@{}}
\toprule
\textbf{Property} & \textbf{Enforcement mechanism} \\
\midrule
\endfirsthead
\toprule
\textbf{Property} & \textbf{Enforcement mechanism} \\
\midrule
\endhead
\bottomrule
\endlastfoot
No use-after-free         & \diagcode{E001} (use-after-move) \\
No double-free            & \diagcode{E002} (double-consume) \\
No leak of linear values  & \diagcode{E003} (unconsumed at scope exit) \\
No dangling references    & \diagcode{E010} (ref to automatic-duration object) \\
No out-of-bounds access   & \diagcode{E014} (VRA proof failure) \\
No null-pointer deref     & No raw pointers in safe fragment \\
No data races             & No concurrency in current version \\
\end{longtable}

\begin{constraint}
  Inside \lain{unsafe} blocks, pointer operations may violate the
  above properties if the programmer's safety assertion is incorrect
  (\S\,\ref{sec:unsafe.ub}).
\end{constraint}

\section{Compiler-internal arena (informative)}
\label{sec:mem.arena}
\index{arena!compiler-internal}

\begin{note}
  The reference implementation uses a bump-pointer arena internally for
  all AST nodes, symbol table entries, and diagnostic objects.  This
  arena is an implementation detail; it is never exposed to user code.
  User programs may implement arena allocators of their own through C
  interop.
\end{note}

% 20-stdlib.tex — placeholder
\chapter{Stdlib}
\label{sec:20-stdlib}

% TODO: write this chapter.


% -------------------------------------------------------------------------
% Back matter — Annexes
% -------------------------------------------------------------------------
\backmatter
\appendix

%% A-grammar.tex — Annex A: Grammar summary (normative)
\chapter{Grammar summary}
\label{annex:A-grammar}
\index{grammar summary}

\pnum
This annex collects all grammar productions defined in this specification.
It is normative.  Productions are listed in the order they are introduced
in the body of the specification.  Each section heading identifies the
originating clause.

\pnum
The grammar uses an extended BNF notation.  \texttt{::=} introduces a
production.  Quoted strings are terminals.  Unquoted names are
non-terminals.  \texttt{\{~X~\}} denotes zero or more repetitions.
\texttt{[~X~]} denotes an optional item.  \texttt{A | B} denotes
alternatives.

% grammar.tex — Complete BNF Grammar for the Lain Programming Language
%
% This file is the single authoritative source for all grammar productions.
% Each rule carries a \label{gram:rule-name} for cross-referencing.
% Productions are grouped by clause; within each group, the order is
% top-down (start symbol first, auxiliary rules after).
%
% Usage in chapter files:
%   \gramref{expression}   → hyperlinked italic "expression"
%   To typeset a rule inline, use the grammar environment from lain-spec.sty.

% =========================================================================
% §5 — Lexical conventions
% =========================================================================

\section*{§5 Lexical grammar}

\subsection*{§5.5 Tokens}
\label{gram:token}
\begin{grammar}
token ::=
    keyword
  | identifier
  | integer-literal
  | float-literal
  | string-literal
  | char-literal
  | bool-literal
  | operator
  | punctuator
\end{grammar}

\subsection*{§5.6 Keywords}
\label{gram:keyword}
\begin{grammar}
keyword ::=
    "func"   | "proc"   | "var"    | "mov"    | "type"
  | "return" | "if"     | "else"   | "while"
  | "for"    | "in"     | "case"   | "import" | "use"
  | "extern" | "c_include" | "unsafe" | "comptime" | "defer"
  | "true"   | "false"  | "and"    | "or"     | "as"
  | "decreasing" | "continue" | "break"
\end{grammar}

\subsection*{§5.7 Identifiers}
\label{gram:identifier}
\begin{grammar}
identifier       ::= identifier-start { identifier-continue }
identifier-start ::= letter | "_"
identifier-continue ::= letter | decimal-digit | "_"
letter           ::= "A" .. "Z" | "a" .. "z"
\end{grammar}

\subsection*{§5.8 Literals}
\label{gram:integer-literal}
\begin{grammar}
integer-literal  ::=
    decimal-literal
  | hex-literal
  | binary-literal
  | octal-literal

decimal-literal  ::= nonzero-digit { decimal-digit | "_" }
                   | "0"
hex-literal      ::= "0x" hex-digit { hex-digit | "_" }
binary-literal   ::= "0b" binary-digit { binary-digit | "_" }
octal-literal    ::= "0o" octal-digit { octal-digit | "_" }

decimal-digit    ::= "0" .. "9"
nonzero-digit    ::= "1" .. "9"
hex-digit        ::= decimal-digit | "a" .. "f" | "A" .. "F"
binary-digit     ::= "0" | "1"
octal-digit      ::= "0" .. "7"
\end{grammar}

\label{gram:float-literal}
\begin{grammar}
float-literal    ::= decimal-literal "." decimal-literal [ float-exponent ]
float-exponent   ::= ( "e" | "E" ) [ "+" | "-" ] decimal-literal
\end{grammar}

\label{gram:string-literal}
\begin{grammar}
string-literal   ::= '"' { string-char } '"'
string-char      ::= any-source-char-except-'"'-or-'\' | escape-sequence
escape-sequence  ::= '\' ( '"' | '\' | "n" | "t" | "r" | "0" )
\end{grammar}

\label{gram:bool-literal}
\begin{grammar}
bool-literal     ::= "true" | "false"
\end{grammar}

\subsection*{§5.9 Operators and punctuators}
\label{gram:operator}
\begin{grammar}
operator         ::=
    "+"  | "-"  | "*"  | "/"  | "%"
  | "==" | "!=" | "<"  | ">"  | "<=" | ">="
  | "and"| "or" | "!" | "in"
  | "&"  | "|"  | "^"  | "<<" | ">>"
  | "="  | "+=" | "-=" | "*=" | "/=" | "%="
  | "&=" | "|=" | "^=" | "<<="| ">>="
  | "->" | "."  | ".."

punctuator       ::=
    "("  | ")"  | "{"  | "}"  | "["  | "]"
  | ":"  | ","  | "@"
\end{grammar}

% =========================================================================
% §6 — Basic concepts
% =========================================================================

\section*{§6 Program structure grammar}
\label{gram:program}
\begin{grammar}
program              ::= { top-level-declaration }
top-level-declaration ::=
    import-declaration
  | type-declaration
  | function-declaration
  | procedure-declaration
  | extern-declaration
  | c-include-declaration

import-declaration   ::= "import" qualified-identifier [ "." "{" identifier { "," identifier } "}" ] newline
type-declaration     ::= "type" identifier [ "(" type-param-list ")" ]
                         ( "=" type | struct-body | enum-body )
type-param-list      ::= type-param { "," type-param }
type-param           ::= identifier "type"
qualified-identifier ::= identifier { "." identifier }
\end{grammar}

% =========================================================================
% §7 — Types
% =========================================================================

\section*{§7 Type grammar}
\label{gram:type}
\label{gram:array-type}
\begin{grammar}
type             ::=
    primitive-type
  | array-type
  | slice-type
  | pointer-type
  | qualified-type
  | constrained-type
  | user-defined-type

primitive-type   ::=
    "int"   | "u8"    | "u16"   | "u32"   | "u64"   | "usize"
  | "i8"    | "i16"   | "i32"   | "i64"   | "isize"
  | "f32"   | "f64"
  | "bool"  | "str"   | "void"

array-type       ::= type "[" integer-literal "]"
slice-type       ::= [ "*" ] type "[]"                    % fat slice; "*" optional
                   | [ "*" ] type "[" size-constraint "]" % sized fat slice
pointer-type     ::= "*" type                             % raw pointer
                   | "*" type "[" integer-literal "]"     % thin ptr + static length
                   | "*" type "[" ":" sentinel-literal "]" % sentinel-terminated
qualified-type   ::= ownership-qualifier type
ownership-qualifier ::= "var" | "mov"
user-defined-type ::= qualified-identifier
constrained-type ::= type relational-op integer-literal
                   { ( "and" | "or" ) type relational-op integer-literal }
relational-op    ::= ">=" | "<=" | ">" | "<" | "=="
\end{grammar}

% =========================================================================
% §8 — Expressions
% =========================================================================

\section*{§8 Expression grammar}
\label{gram:expression}
\begin{grammar}
expression           ::= assignment-expression

assignment-expression ::=
    logical-or-expression
  | lvalue-expression assignment-operator assignment-expression

assignment-operator  ::=
    "="  | "+=" | "-=" | "*=" | "/=" | "%="
  | "&=" | "|=" | "^=" | "<<="| ">>="

logical-or-expression ::= logical-and-expression { "or" logical-and-expression }

logical-and-expression ::= in-expression { "and" in-expression }

in-expression        ::= equality-expression { "in" equality-expression }

equality-expression  ::= relational-expression
                        { ( "==" | "!=" ) relational-expression }

relational-expression ::= shift-expression
                         { ( "<" | ">" | "<=" | ">=" ) shift-expression }

shift-expression     ::= additive-expression
                        { ( "<<" | ">>" ) additive-expression }

additive-expression  ::= multiplicative-expression
                        { ( "+" | "-" ) multiplicative-expression }

multiplicative-expression ::= cast-expression
                             { ( "*" | "/" | "%" ) cast-expression }

cast-expression      ::= unary-expression { "as" type }

unary-expression     ::=
    postfix-expression
  | unary-operator unary-expression

unary-operator       ::= "-" | "!" | "mov" | "var"

postfix-expression   ::=
    primary-expression
  | postfix-expression "." identifier
  | postfix-expression "[" expression "]"
  | postfix-expression "(" [ argument-list ] ")"

argument-list        ::= expression { "," expression }

primary-expression   ::=
    identifier
  | literal
  | "(" expression ")"
  | block-expression
  | if-expression
  | match-expression
  | constructor-expression
  | comptime-expression

literal              ::=
    integer-literal
  | float-literal
  | string-literal
  | bool-literal

block-expression     ::= "{" { statement } expression "}"

if-expression        ::=
    "if" expression block-expression
  [ "else" ( if-expression | block-expression ) ]

match-expression     ::=
    "case" [ "&" ] expression "{"
    { match-arm }
    "}"

match-arm            ::= ( pattern-list | "else" ) ":" arm-body

arm-body             ::= expression newline | block-expression

pattern-list         ::= pattern { "," pattern }

pattern              ::=
    literal
  | literal ".." literal                        % inclusive range pattern
  | string-literal
  | qualified-identifier                         % unit variant
  | qualified-identifier "(" [ binding-list ] ")" % variant with bound fields

binding-list         ::= identifier { "," identifier }

constructor-expression ::=
    qualified-identifier "{" [ field-init-list ] "}"
  | qualified-identifier "(" [ argument-list ] ")"

field-init-list      ::= identifier ":" expression { "," identifier ":" expression }

comptime-expression  ::= "comptime" block-expression

lvalue-expression    ::=
    identifier
  | postfix-expression "." identifier
  | postfix-expression "[" expression "]"
\end{grammar}

% =========================================================================
% §9 — Statements
% =========================================================================

\section*{§9 Statement grammar}
\label{gram:statement}
\begin{grammar}
statement            ::=
    block-statement
  | immutable-declaration
  | var-declaration
  | assignment-statement
  | expression-statement
  | if-statement
  | while-statement
  | bounded-while-statement
  | return-statement
  | defer-statement
  | unsafe-block
  | comptime-if-statement

block-statement      ::= "{" { statement } "}"

immutable-declaration ::= identifier [ type ] "=" expression newline

var-declaration      ::= "var" identifier [ type ] "=" expression newline

assignment-statement ::= lvalue-expression assignment-operator expression newline

expression-statement ::= expression newline

if-statement         ::=
    "if" expression block-statement
  { "else" "if" expression block-statement }
  [ "else" block-statement ]

while-statement      ::= "while" expression block-statement

bounded-while-statement ::=
    "while" expression "decreasing" expression block-statement

return-statement     ::= "return" [ expression ] newline

defer-statement      ::= "defer" block-statement

unsafe-block         ::= "unsafe" block-statement

comptime-if-statement ::=
    "comptime" "if" expression block-statement
  [ "else" block-statement ]
\end{grammar}

% =========================================================================
% §10 — Declarations
% =========================================================================

\section*{§10 Declaration grammar}
\label{gram:function-declaration}
\label{gram:struct-declaration}
\label{gram:enum-declaration}
\begin{grammar}
function-declaration  ::=
    { attribute } "func" identifier
    "(" [ param-list ] ")" return-annotation
    block-statement

procedure-declaration ::=
    { attribute } "proc" identifier
    "(" [ param-list ] ")" return-annotation
    block-statement

attribute            ::= "[" identifier [ "(" argument-list ")" ] "]" newline

return-annotation    ::= type [ relational-op integer-literal ]
                       | ( empty )

param-list           ::= param { "," param }
param                ::=
    [ ownership-qualifier ] identifier type              % value parameter
  | identifier "type"                                    % type parameter (generics)
  | [ ownership-qualifier ] "{" identifier { "," identifier } "}" type
                                                         % field-destructuring parameter

struct-body          ::= "{" { field-declaration } "}"
field-declaration    ::= [ "mov" ] identifier type newline

enum-body            ::= "{" { variant-declaration } "}"
variant-declaration  ::=
    identifier newline
  | identifier "{" field-declaration-list "}" newline

field-declaration-list ::= field-declaration { field-declaration }
\end{grammar}


%% B-diagnostics.tex — Annex B: Diagnostic codes (normative)
\chapter{Diagnostic codes}
\label{annex:B-diagnostics}
\index{diagnostic codes}

\pnum
This annex is normative.  It lists every diagnostic code defined by this
specification, the clause that requires the diagnostic, the category
(constraint violation, warning, or informational), and a representative
message text.  A conforming implementation shall issue a diagnostic for
every constraint violation in the table below.

\pnum
Diagnostic codes are divided into the following ranges:

\begin{longtable}{@{}ll@{}}
\toprule
\textbf{Range} & \textbf{Domain} \\
\midrule
\endfirsthead
\toprule
\textbf{Range} & \textbf{Domain} \\
\midrule
\endhead
\bottomrule
\endlastfoot
\diagcode{E001}–\diagcode{E010}, \diagcode{E016} & Ownership and lifetime \\
\diagcode{E011}–\diagcode{E015} & Purity, safety, and verification \\
\diagcode{E014}                 & Compile-time evaluation \\
\diagcode{E080}–\diagcode{E082} & Bounded while verification \\
\diagcode{E083}–\diagcode{E087} & Array/slice safety and integer boundaries \\
\diagcode{E100}                 & Miscellaneous \\
\diagcode{VRA}                  & Value range analysis \\
\end{longtable}

% -------------------------------------------------------------------------
\section*{Ownership and lifetime diagnostics (\diagcode{E001}–\diagcode{E010}, \diagcode{E016})}
% -------------------------------------------------------------------------

\begin{longtable}{@{}p{1.5cm}p{3.5cm}p{7.5cm}@{}}
\toprule
\textbf{Code} & \textbf{Clause} & \textbf{Message and trigger} \\
\midrule
\endfirsthead
\toprule
\textbf{Code} & \textbf{Clause} & \textbf{Message and trigger} \\
\midrule
\endhead
\bottomrule
\endlastfoot

\diagcode{E001}
  & \S\,\ref{sec:own.diag}
  & \textit{use of moved value '\texttt{x}'} ---
    Issued when an expression reads or passes a variable whose linear
    state is \textsc{Consumed}.
    \\[4pt]

\diagcode{E002}
  & \S\,\ref{sec:own.diag}
  & \textit{value '\texttt{x}' consumed twice} ---
    Issued when a second \lain{mov} expression is applied to a variable
    already in state \textsc{Consumed}.
    \\[4pt]

\diagcode{E003}
  & \S\,\ref{sec:own.diag}
  & \textit{linear value '\texttt{x}' not consumed before end of scope} ---
    Issued when a \lain{mov}-mode variable exits scope in state
    \textsc{Unconsumed}.
    \\[4pt]

\diagcode{E004}
  & \S\,\ref{sec:own.diag}
  & \textit{cannot borrow '\texttt{x}' as mutable; it is already borrowed}
    --- Issued when a mutable borrow conflicts with an existing borrow
    (reads or write).
    \\[4pt]

\diagcode{E005}
  & \S\,\ref{sec:own.diag}
  & \textit{cannot borrow '\texttt{x}' as shared; it is already mutably
    borrowed} --- Issued when a shared borrow is requested while a
    mutable borrow is active.
    \\[4pt]

\diagcode{E006}
  & \S\,\ref{sec:own.diag}
  & \textit{cannot copy linear type '\texttt{T}'} ---
    Issued when assignment or argument passing implicitly copies a type
    that has a \lain{mov}-mode field.
    \\[4pt]

\diagcode{E007}
  & \S\,\ref{sec:own.modes}
  & \textit{cannot assign through a shared reference '\texttt{x}'} ---
    Issued when a write is attempted through a shared (non-\lain{var})
    binding.
    \\[4pt]

\diagcode{E010}
  & \S\,\ref{sec:mem.storage}
  & \textit{cannot return reference to local variable '\texttt{x}'} ---
    Issued when a \lain{return} statement returns a reference whose
    referent has automatic storage duration.
    \\[4pt]

\diagcode{E016}
  & \S\,\ref{sec:own.diag}
  & \textit{linear variable '\texttt{x}' is used inconsistently in
    the branches of \emph{stmt}} --- Issued when a linear value is
    consumed in some control-flow paths of an \lain{if}/\lain{case}
    but not in others, leaving an ambiguous state at the join point.
    Field-level inconsistency on linear struct fields is reported with
    the same code.
    \\[4pt]

\end{longtable}

% -------------------------------------------------------------------------
\section*{Purity, safety, and verification (\diagcode{E011}–\diagcode{E015})}
% -------------------------------------------------------------------------

\begin{longtable}{@{}p{1.5cm}p{3.5cm}p{7.5cm}@{}}
\toprule
\textbf{Code} & \textbf{Clause} & \textbf{Message and trigger} \\
\midrule
\endfirsthead
\toprule
\textbf{Code} & \textbf{Clause} & \textbf{Message and trigger} \\
\midrule
\endhead
\bottomrule
\endlastfoot

\diagcode{E011}
  & \S\,\ref{sec:fn.purity}
  & \textit{impure operation in pure function '\texttt{f}'} ---
    Issued when a \lain{func} body contains a call to a \lain{proc}, an
    unbounded \lain{while} without a \lain{decreasing} measure, or a
    direct/mutual recursion without a measure.
    \\[4pt]

\diagcode{E012}
  & \S\,\ref{sec:interop.pointers}
  & \textit{pointer cast outside unsafe block} ---
    Issued when a pointer-to-pointer or integer-to-pointer cast occurs
    outside an \lain{unsafe} block.
    \\[4pt]

\diagcode{E013}
  & \S\,\ref{sec:basics.declarations}
  & \textit{use of undeclared identifier '\texttt{x}'} ---
    Issued when an identifier is referenced before its declaration in
    scope.
    \\[4pt]

\diagcode{E014}
  & \S\,\ref{sec:expr.index},
    \S\,\ref{sec:13-vra}
  & \textit{index '\texttt{i}' not proven in bounds of '\texttt{arr}'} ---
    Issued when VRA cannot prove that an array index expression is
    within $[0,\, n-1]$.  Also issued for non-exhaustive \lain{case}
    expressions when combined with the exhaustiveness diagnostic (shared
    code path).
    \\[4pt]

\diagcode{E015}
  & \S\,\ref{sec:expr.arith}
  & \textit{division by zero} ---
    Issued when VRA can prove statically that a division or modulo
    expression has a zero divisor.
    \\[4pt]

\end{longtable}

% -------------------------------------------------------------------------
\section*{Compile-time evaluation (\diagcode{E014})}
% -------------------------------------------------------------------------

\begin{longtable}{@{}p{1.5cm}p{3.5cm}p{7.5cm}@{}}
\toprule
\textbf{Code} & \textbf{Clause} & \textbf{Message and trigger} \\
\midrule
\endfirsthead
\toprule
\textbf{Code} & \textbf{Clause} & \textbf{Message and trigger} \\
\midrule
\endhead
\bottomrule
\endlastfoot

\diagcode{E014}
  & \S\,\ref{sec:gen.recursion-depth}
  & \textit{comptime evaluation depth limit exceeded} ---
    Issued when compile-time function evaluation recurses beyond the
    implementation's depth limit (256 in the reference implementation).
    \\[4pt]

\end{longtable}

% -------------------------------------------------------------------------
\section*{Bounded while verification (\diagcode{E080}–\diagcode{E082})}
% -------------------------------------------------------------------------

\begin{longtable}{@{}p{1.5cm}p{3.5cm}p{7.5cm}@{}}
\toprule
\textbf{Code} & \textbf{Clause} & \textbf{Message and trigger} \\
\midrule
\endfirsthead
\toprule
\textbf{Code} & \textbf{Clause} & \textbf{Message and trigger} \\
\midrule
\endhead
\bottomrule
\endlastfoot

\diagcode{E080}
  & \S\,\ref{sec:stmt.while.bounded}
  & \textit{termination measure not proven non-negative} ---
    Issued when the compiler cannot prove that the \lain{decreasing}
    expression is $\geq 0$ whenever the loop condition holds.
    \\[4pt]

\diagcode{E081}
  & \S\,\ref{sec:stmt.while.bounded}
  & \textit{termination measure does not depend on a loop variable} ---
    Issued when the \lain{decreasing} expression consists entirely of
    compile-time constants and therefore cannot strictly decrease.
    \\[4pt]

\diagcode{E082}
  & \S\,\ref{sec:stmt.while.bounded}
  & \textit{termination measure not proven to decrease in loop body} ---
    Issued when the compiler cannot verify that every execution of the
    loop body strictly decreases the measure by at least 1.
    \\[4pt]

\end{longtable}

% -------------------------------------------------------------------------
\section*{Miscellaneous (\diagcode{E100})}
% -------------------------------------------------------------------------

\begin{longtable}{@{}p{1.5cm}p{3.5cm}p{7.5cm}@{}}
\toprule
\textbf{Code} & \textbf{Clause} & \textbf{Message and trigger} \\
\midrule
\endfirsthead
\toprule
\textbf{Code} & \textbf{Clause} & \textbf{Message and trigger} \\
\midrule
\endhead
\bottomrule
\endlastfoot

\diagcode{E100}
  & \S\,\ref{sec:05-lexical}, \S\,\ref{sec:07-types}
  & \textit{ill-formed token / semantic constraint violation} ---
    A catch-all code used for lexical and syntactic errors without a
    dedicated code, including: missing semicolons (\S\,\ref{sec:lex.disambiguation}),
    using \lain{==} on enum types (\S\,\ref{sec:types.enum}), and other
    single-token syntactic errors.
    \\[4pt]

\end{longtable}

% -------------------------------------------------------------------------
\section*{Array/slice safety and integer boundaries (\diagcode{E083}–\diagcode{E087})}
% -------------------------------------------------------------------------

\begin{longtable}{@{}p{1.5cm}p{3.5cm}p{7.5cm}@{}}
\toprule
\textbf{Code} & \textbf{Clause} & \textbf{Message and trigger} \\
\midrule
\endfirsthead
\toprule
\textbf{Code} & \textbf{Clause} & \textbf{Message and trigger} \\
\midrule
\endhead
\bottomrule
\endlastfoot

\diagcode{E083}
  & \S\,\ref{sec:types.struct}
  & \textit{linear field '\texttt{f}' missing \texttt{mov} annotation} ---
    Issued when a struct or enum variant field has a linear type but lacks
    the \lain{mov} keyword.  (Q-008 transitivity rule.)
    \\[4pt]

\diagcode{E085}
  & \S\,\ref{sec:expr.index}, \S\,\ref{sec:expr.subslice}
  & \textit{index or slice endpoint not proven in bounds} ---
    Issued when VRA cannot prove a non-negativity or upper-bound requirement
    for an array/slice index or sub-slice endpoint.
    \\[4pt]

\diagcode{E086}
  & \S\,\ref{sec:vra.boundary}
  & \textit{integer overflow at boundary: range exceeds target type} ---
    Issued when the VRA range of a source value provably exceeds the
    natural range of the target type at a type boundary (declaration,
    assignment, call argument, return, struct constructor).
    \\[4pt]

\diagcode{E087}
  & \S\,\ref{sec:expr.subslice}, \S\,\ref{sec:vra.sized-slice-callsite}
  & \textit{sized slice constraint violated} ---
    Issued (1) when both endpoints of a sub-slice \lain{arr[a..b]} are
    integer literals and $a > b$; (2) at a call site where VRA can prove
    that the argument does not satisfy the formal parameter's sized slice
    constraint (\lain{T[k]}, \lain{T[>= k]}, \lain{T[> k]}, or
    \lain{T[ref.len]}).
    \\[4pt]

\end{longtable}

% -------------------------------------------------------------------------
\section*{Value range analysis (\diagcode{VRA})}
% -------------------------------------------------------------------------

\begin{longtable}{@{}p{1.5cm}p{3.5cm}p{7.5cm}@{}}
\toprule
\textbf{Code} & \textbf{Clause} & \textbf{Message and trigger} \\
\midrule
\endfirsthead
\toprule
\textbf{Code} & \textbf{Clause} & \textbf{Message and trigger} \\
\midrule
\endhead
\bottomrule
\endlastfoot

\diagcode{VRA}
  & \S\,\ref{sec:13-vra}
  & \textit{value range proof required} ---
    Issued when VRA cannot prove that an operation satisfies its
    required constraint and more specific code is not available.
    Appears as a diagnostic prefix in messages such as:
    ``\texttt{[VRA] cannot prove index in bounds}'' and
    ``\texttt{[VRA] return value may not satisfy constraint}''.
    \\[4pt]

\end{longtable}

\section*{Additional diagnostics (\diagcode{E008}, \diagcode{E009}, \diagcode{E017}--\diagcode{E125})}

The following diagnostics are emitted by the reference implementation and are
normative. Full clause cross-references are pending an annex-completion pass.

\begin{longtable}{@{}ll@{}}
\toprule
\textbf{Code} & \textbf{Message and trigger} \\
\midrule
\endhead
\bottomrule
\endlastfoot
\diagcode{E008} & cannot move a value that is currently borrowed \\[2pt]
\diagcode{E009} & cannot mutate a field through a raw (shared) pointer --- use a \lain{var} borrow \\[2pt]
\diagcode{E017} & passing to a \lain{var} parameter requires explicit \lain{var} at the call site \\[2pt]
\diagcode{E018} & a function can reach its end without returning a value \\[2pt]
\diagcode{E019} & use of a partially-initialized value \\[2pt]
\diagcode{E060} & pointer dereference outside an \lain{unsafe} block \\[2pt]
\diagcode{E063} & a possibly-marker union value used without testing it first \\[2pt]
\diagcode{E064} & a \lain{|}-union cannot be niche-optimized (no spare bit pattern) \\[2pt]
\diagcode{E084} & access to a \lain{[private]} declaration from another module \\[2pt]
\diagcode{E101} & calling a \lain{proc} from a \lain{comptime} context (only \lain{func} is allowed) \\[2pt]
\diagcode{E102} & expected an attribute name after the opening bracket \\[2pt]
\diagcode{E103} & unknown attribute \\[2pt]
\diagcode{E104} & a type contains itself by value (infinite size) \\[2pt]
\diagcode{E105} & identifier is defined in another module --- add an \lain{import} \\[2pt]
\diagcode{E121} & struct-field packing supports only iN/uN fields totaling $\le 64$ bits \\[2pt]
\diagcode{E122} & a function pointer must be initialized with a matching function \\[2pt]
\diagcode{E123} & wrong number of arguments in a call through a function pointer \\[2pt]
\diagcode{E124} & type application on a non-generic type \\[2pt]
\diagcode{E125} & direct ADT field access outside \lain{unsafe} --- destructure with \lain{case} \\[2pt]
\end{longtable}

%% C-rationale.tex — Annex C (placeholder)
\chapter{Rationale}
\label{annex:C-rationale}

% TODO: Iter 19 — write this annex.

%% D-comparison.tex — Annex D (placeholder)
\chapter{Comparison}
\label{annex:D-comparison}

% TODO: Iter 19 — write this annex.


% -------------------------------------------------------------------------
% Index
% -------------------------------------------------------------------------
\printindex

\end{document}
